% Vorlage fuer Abschlussarbeiten am Fachbereich Mathematik der TU Darmstadt.
% Geeignet fuer Bachelorarbeiten und Masterarbeiten.

\documentclass[12pt,a4paper,twoside]{report}
% Benutzt man die Option draft, so kann man die Umbrueche ueberpruefen.

% Hier werden alle benoetigten Pakete und Einstellungen geladen. Auch
% hier sind Sie frei diese direkt in der Praeambel zu laden.

% Folgende Pakete sind Minimalvoraussetzung.

% Anpassung von LaTeX an die deutsche Sprache. Zum Beispiel wird
% \chapter{} direkt als "`Kapitel"' uebersetzt usw. Ausserdem
% wird die Silbentrennung nach deutscher Rechtschreibung vorgenommen.
% Sonderzeichen, z.B. Umlaute, koennen direkt eingegeben werden.

%\usepackage[ngerman]{babel}
\usepackage[ngerman,english]{babel}

% Kodierung fuer unixoide Systeme und Windowssysteme.

\usepackage[utf8]{inputenc}

% Schriftart Times New Roman.

%\usepackage{times}
%\usepackage{lmodern}
\usepackage{mathptmx}

% Erweitert den Zeichenvorrat, so dass z.B. auch Umlaute im PDF-Dokument
% gefunden werden.

\usepackage[T1]{fontenc}

% Zur Einbindung von Bildern.

\usepackage{graphicx}

% Erweiterte enumerate-Umgebung.

\usepackage{enumerate}

% Verschiedene Pakete, die nuetzlich sind um Mathematik in LaTeX zu setzen.

\usepackage{mathtools}
\usepackage{amsmath, amssymb, amsthm, dsfont}
\usepackage{relsize}
\usepackage{hyperref}

%%%%%%%%%%%%%%%%
% Seitenlayout %
%%%%%%%%%%%%%%%%

% DIV# gibt den Divisor f�r die Layoutberechnung an.
% Vergr��ern des Divisors vergr��ert den Textbereich.
% BCOR#cm gibt die Breite des Bundstegs an.
\usepackage[DIV14,BCOR2cm]{typearea}

% Abstand obere Blattkante zur Kopfzeile ist 2.54cm - 15mm
%\setlength{\topmargin}{-15mm}

% Keine Einrueckung nach einem Absatz.

%\parindent 0pt

% Abstand zwischen zwei Abs\"atzen.

%\parskip 12pt

% Zeilenabstand.
\usepackage{bibgerm}

\linespread{1.3}

% Inhaltsverzeichnis erstellen.

\usepackage{makeidx}
\makeindex

\usepackage[inline]{showlabels}
\allowdisplaybreaks


% Hier werden Makros und Umgebungen eingebunden. Diese werden separat
% in der Datei befehle.tex definiert. Sie sind frei diese Befehle auch
% direkt in der Praeambel zu definieren.
% Makros

\newcommand{\N}{\mathbb{N}} % natuerliche Zahlen
\newcommand{\Z}{\mathbb{Z}} % ganze Zahlen
\newcommand{\Q}{\mathbb{Q}} % rationale Zahlen
\newcommand{\R}{\mathbb{R}} % reelle Zahlen
\newcommand{\C}{\mathbb{C}}
\DeclareMathOperator{\Ret}{Re} % Guter Realteil

%%Mengen
\newcommand{\BB}{\mathrm{B}} 
\newcommand{\EE}{\mathrm{E}} 
\newcommand{\verCone}{\Gamma_\mathrm{V}} 
\newcommand{\ShenCone}{\Gamma} 

%%Funktionenr�ume
\newcommand{\CC}{\mathrm{C}} 
\newcommand{\Ell}{\mathrm{L}} 
\newcommand{\Li}{\mathfrak{L}}
\newcommand{\WW}{\mathrm{W}}
\newcommand{\HH}{\mathrm{H}}

%%Mathematische Operatoren

\newcommand{\diam}{\operatorname{diam}}
\newcommand{\K}{\mathcal{K}}
\newcommand{\slp}{\mathcal{S}}
\newcommand{\MM}{\mathcal{M}}
\newcommand{\supp}{\operatorname{supp}}
\newcommand{\dlp}{\mathcal{D}}
\DeclareMathOperator{\pv}{p. v.}
\renewcommand{\Im}{\operatorname{Im}}
\renewcommand{\Re}{\operatorname{Re}}

%Differentialoperatoren
\let\divsymb=\div % rename builtin command \div to \divsymb
\renewcommand{\div}[1]{\mathrm{div\,} #1} % for divergence
\renewcommand{\d}[1]{\ensuremath\, {\operatorname{d}\!{#1}}}
\newcommand{\DD}{\ensuremath\,{\operatorname{D}}}

\newcommand{\ii}{\mathrm{i}}
\newcommand{\e}{\mathrm{e}}

% Umgebungen f�r Definitionen, S�tze, usw.

\theoremstyle{definition}
\newtheorem{defn}{Definition}[chapter]
\newtheorem{Beispiel}{Beispiel}[chapter]
\newtheorem*{RBsp}{Regelbeispiel}

\theoremstyle{plain}
\newtheorem{thm}{Theorem}[chapter]
\newtheorem{lem}[thm]{Lemma}
\newtheorem{cor}[thm]{Corollary}
\newtheorem{prop}[thm]{Proposition}

\theoremstyle{remark}
\newtheorem{rem}[thm]{Remark}

\def\Xint#1{\mathchoice
  {\XXint\displaystyle\textstyle{#1}}%
  {\XXint\textstyle\scriptstyle{#1}}%
  {\XXint\scriptstyle\scriptscriptstyle{#1}}%
  {\XXint\scriptscriptstyle\scriptscriptstyle{#1}}%
\!\int}
\def\XXint#1#2#3{{\setbox0=\hbox{$#1{#2#3}{\int}$ }
\vcenter{\hbox{$#2#3$ }}\kern-.6\wd0}}
\def\ddashint{\Xint=}
\def\dashint{\Xint-}


\begin{document}
% Auf der Titelseite und im Inhaltsverzeichnis sollen keine
% Seitenzahlen dargestellt werden.
\pagestyle{empty}

% Hier wird die Titelseite eingebunden.
\begin{titlepage}
  \begin{center}
    \includegraphics[width=0.5\linewidth]{TU_Darmstadt_Logo.pdf}
    \vfill
    
    \large{Fachbereich Mathematik}
    \vfill
    
    \large{Masterthesis}
    \vfill

    \huge{On Resolvent Estimates in $\Ell^p$ for the Stokes Operator in Lipschitz Domains}
    \vfill
    
		\large
    Fabian Gabel

    \large 15.04.2018
    \vfill
\begin{tabular}{rl}
    Gutachter:& Prof. Dr. Robert Haller-Dintelmann
    \\
    Betreut durch:& Dr. Patrick Tolksdorf
\end{tabular}
  \end{center}
\end{titlepage}


% Inhaltsverzeichnis erstellen.
\tableofcontents

% Ab sofort werden Seitenzahlen in der Kopfzeile angezeigt.
\pagestyle{headings}

% Einbinden der Einleitung.
\phantomsection
\chapter*{Introduction}
\addcontentsline{toc}{chapter}{Introduction}

In the solution theory for nonlinear partial differential equations, an integral part of the solution process is often to develop a semigroup theory for the linearization of the equation.
In the case of the famous \emph{Navier-Stokes equations} which for a given domain $\Omega \subseteq \R^d$, $d \geq 2$, describe the behavior of a Newtonian fluid over time, the linearization is given by the \emph{Stokes equation}s
\begin{align*}
  \partial_t u - \Delta u + \nabla \pi &= 0 \quad\text{in } \Omega\,, \;t > 0\,,\\
  \div(u) &= 0 \quad\text{in } \Omega\,,\; t > 0\,, \\
  u(0) &= a \quad\text{in } \Omega\,, \\
  u &= 0 \quad\text{on } \partial\Omega\,,\; t > 0\,,
\end{align*}
where $u \colon \R^+ \times \Omega \to \R^d$ stands for the velocity field and $\pi \colon \R^+ \times \Omega \to \R$ represents the pressure of the fluid.
The so-called \emph{Stokes semigroup} $(\e^{-tA})_{t \geq 0}$ describes the evolution of the velocity $u$ and the \emph{Stokes operator} $A$ corresponds to the term ``$-\Delta u + \nabla \pi$'' in the Stokes equations. 

Having a semigroup makes it possible to look for \emph{mild solutions} to the Navier-Stokes equations using a variation of constants formula to construct an iteration method.
This approach was introduced by Fujita and Kato \cite{fujitaKato,katoFujita} and builds mainly on resolvent estimates for the Stokes operator $A$ and the analyticity property of the Stokes semigroup.
Fujita and Kato applied their methods to find solutions in $\Ell^2$ spaces, given that the initial values $a$ were to be found in domains of fractional powers of the Stokes operator.
Fractional power domains are in general difficult to describe which is why Fujita and Kato suggested that an $\Ell^p$ theory for the Navier-Stokes equations could overcome this problem, see \cite[p.\@~313]{fujitaKato}.
%The interest in an $\Ell^p$ theory of the Stokes equations became eminent as they were looking to find spaces with pertinence conditions that are easier to verify.
Resolvent estimates in $\Ell^p$ for the Stokes operator and the analyticity of the Stokes semigroup are crucial ingredients of this classical functional analytic approach to the solution theory of the Navier-Stokes equations.

On smooth domains, the first milestone in the direction of an $\Ell^p$ theory was laid by Giga \cite{giga} and Solonnikov \cite{solonnikov} who proved that on bounded smooth domains the Stokes operator generates a bounded analytic semigroup in $\Ell^p(\Omega)$ for $1 < p < \infty$.
This result was then used by Giga and Miyakawa \cite{gigaMiyakawa} to prove the existence of mild solutions given $\Ell^3$ initial data.
Kato \cite{katoExtend} then extended this theory to the whole space and later Giga \cite{gigaAdjust} adjusted the approach for bounded smooth domains.

A natural assumption in real world applications is that the boundary of the domain $\Omega$ is not smooth but merely Lipschitz continuous. 
On Lipschitz domains, the analysis of partial differential equations is more complicated as classical localization techniques fail due to the lack of smoothness of the boundary.
%In order to extend the advances of the $\Ell^p$ theory on smooth domains to Lipschitz domains, the first existence and unique results were given by Mitrea and Monniaux.
In order to overcome these problems, new techniques had to be developed which gave rise to the work of Fabes, Kenig and Verchota~\cite{fabesKenigVerchota} who constructed solutions to the $\Ell^2$ Dirichlet problem of the Stokes equations by the method of layer potentials.

As on smooth domains, the further expansion of the $\Ell^p$ theory by means of the classical iteration method was reliant on the analyticity of the Stokes semigroup.
Due to the boundedness properties of the Helmholtz projection, Taylor conjectured in his paper \cite{taylor} from 2000 that the analyticity of the Stokes semigroup would only hold in a neighborhood of $p = 2$, namely for $p \in  (3/2 - \varepsilon, 3 + \varepsilon)$ in the three dimensional case.
His conjecture was supported by the fact that, not much later, Deuring showed in \cite{deuring} the existence of a three dimensional Lipschitz domain such that the needed resolvent estimates would fail for sufficiently large $p$.

Twelve years passed until a positive result could be given to Taylor's conjecture:  Shen showed in his seminal paper \cite{Shen2012} that in bounded Lipschitz domains $\Omega \subseteq \R^d$, $d \geq 3$, for all $p \in (2d/( d + 1) - \varepsilon , 2d/(d - 1) + \varepsilon)$ the Stokes operator generates a bounded analytic semigroup on $\Ell^p_\sigma(\Omega)$.
Shen's result on the analyticity of the semigroup thus made the Fujita-Kato approach available for the study of the Navier-Stokes equations on Lipschitz domains which, together with other methods, was used by Tolksdorf to prove the existence of mild solutions in $\Ell^3$ to the Navier-Stokes equations \cite{tolksdorf}.

In the present thesis, we will study Shen's approach to the analyticity problem of the Stokes semigroup in $\Ell^p$ for bounded Lipschitz domains in $\R^d$.
While Shen's result was only formulated for $d \geq 3$, we will show that his approach can be extended to the two dimensional case by proving suitable estimates on the fundamental solutions of the Stokes resolvent problem.

Except for the first chapter, in which we will gather the fundamentals and formulate the problem under consideration, the rest of this thesis will closely follow the structure of Shen's work \cite{Shen2012}. 

In Chapter \ref{chap:2}, we will derive the central estimates of the fundamental solutions to the Stokes resolvent problem on $\R^d$ by taking advantage of their explicit representation formula.

Chapter \ref{chap:3} will introduce the method of boundary layer potentials for the solution of the Stokes resolvent problem on bounded Lipschitz domains.
We will discover the central relation between boundary values that are attained nontangentially and the singular integral operators that provide a representation formula for solutions to the Stokes resolvent problem. 

In Chapter \ref{chap:4}, we take a more general look on solutions to the Stokes resolvent problem and derive Rellich-type estimates for these solutions.

Based on the results of Chapter 3 and 4, in Chapter \ref{chap:5}, we will study the $\Ell^2$ Dirichlet problem of the Stokes resolvent system. In this chapter, we will see that for given boundary values in $\Ell^2(\partial\Omega; \C^d)$, there exists a unique solution to the $\Ell^2$ Dirichlet problem which is given via a double layer potential.

Finally, in Chapter \ref{chap:6}, we will tackle the problem of analyticity of the Stokes operator. In this chapter, we will derive the necessary resolvent estimates on $\Ell^p$ by making use of Shen's extrapolation theorem which can be considered a refined version of the Calder\'on-Zygmund Lemma.

This last paragraph is reserved to mention some of the persons that were involved in the course of this thesis.
I thank Prof. Dr. Robert Haller-Dintelmann for taking the responsibility of supervising this thesis.
Furthermore, I want to express my gratitude to Dr. Patrick Tolksdorf. 
This thesis would not have been possible without his generous support in innumerable occasions.
I also want to thank Sebastian Bechtel for proofreading parts of this thesis. 
Last but not least, I want to thank my family for supporting me during my studies.



% Einbinden der einzelnen Kapitel (zur �bersicht eine Datei je Kapitel)

\chapter{Fundamentals}

The purpose of this chapter is to collect basic definitions that will be used throughout the subsequent chapters. Furthermore we want to formulate the main problem regarding the resolvent estimates of the Stokes operator.
Throughout this chapter we let $d$ always denote a natural number greater or equal to 2.

\section{Lipschitz-Domains}

In this first section we will establish the fundamental notions regarding bounded Lipschitz domains.

\begin{defn}
  \label{defn:lipschitzDomain}
  Let $\Omega \subset \R^d$ be a bounded open connected set.
  We call $\Omega$ a \emph{bounded Lipschitz domain} if there exist $r_0, M > 0$ such that for all $x \in \partial\Omega$ there exists a function $\eta_x \colon \R^{d - 1} \to \R$ which is Lipschitz continuous and fulfills $\eta_x(0) = 0$ and $\|\nabla \eta_x \|_{\Ell^\infty(\R^{d - 1})} \leq M$, and a rotation $R_x \colon \R^d \to \R^d$ such that for all $0 < r \leq r_0$ 
  \begin{align*}
    R_x[ \Omega - \{x\} ] \cap D(r) &= D_{\eta_x}(r) \\
    R_x[\partial\Omega - \{x\}] \cap D(r) &= I_{\eta_x}(r),
  \end{align*}
  where
  \begin{align*}
    D(r) &\coloneqq \{ (x', x_d) \colon |x'| < r, |x_d| < 10 d (M + 1) r\} \\
    D_{\eta_x}(r) &\coloneqq \{ (x', x_d) \colon |x'| < r, \eta_x(x') < x_d < 10 d(M+ 1)r\} \\
    I_{\eta_x}(r) &\coloneqq \{ (x', x_d) \colon |x'| < r, \eta_x(x') = x_d \} .
  \end{align*}
\end{defn}
  It is common to refer to sets of the form $D_{\eta_x}$ as \emph{Lipschitz cylinders}. If the point $x$ in the definition of Lipschitz cylinders is not of particular importance we will denote the Lipschitz cylinder by $D_{\eta}(r)$.

  If $\Omega$ is a bounded Lipschitz domain, $x \in \partial\Omega$ and $0 < r \leq r_0$, then we may define $U_{x,r} \coloneqq \{ x\} + R_x^{-1} D(r)$, where $R_x$ is the rotation corresponding to $x$ from Definition \ref{defn:lipschitzDomain}.
  This is all we need to define the Lipschitz character of a bounded Lipschitz domain $\Omega$ as suggested by Pipher and Verchota in \cite{pipherVerchota}.

  \begin{defn}
    Let $\Omega \subset \R^d$, $d \geq 2$, be a bounded Lipschitz domain and $x_1, \dots,x_n \in \partial\Omega$ be such that $\{U_{x_i,r_0} \}_{i = 1}^n$ covers $\partial\Omega$.  
    Furthermore let $M$ be the constant from Definition \ref{defn:lipschitzDomain}.
    Then a constant $C > 0$ is said to depend on the \emph{Lipschitz character of} $\Omega$ if it depends on $M$ and $n$.
  \end{defn}

  That the Lipschitz character is indeed a fruitful concept will be emphasized by the following theorem.
  This result is a crucial ingredient in the proof of the Rellich estimates in Chapter \ref{chap:4} as it provides a useful approximating property of Lipschitz domains. In short it enables us to approximate a bounded Lipschitz domain $\Omega$ by a sequence $(\Omega_j)$ of $\CC^\infty$ domains in such a way that estimates on $\Omega_j$ with bounding constants that only depend on the Lipschitz characters may be extended to $\Omega$ when taking the limit.
  The original proof of this Theorem goes back to Ne\v{c}as \cite{necas} and Verchota \cite{verchota} . 
  The presented version of this theorem appeared in Brown \cite{brown}.

\begin{thm}[Ne\v{c}as, Verchota]
  \label{thm:smoothApproximation}
  Let $\Omega$ be a Lipschitz domain.
  Then there exists a sequence of $\CC^\infty$-domains $(\Omega_k)$ with uniform Lipschitz characters, corresponding homeomorphisms $\Lambda_k \colon \partial\Omega \to \partial\Omega_k$, functions $\vartheta_k \colon \partial\Omega \to \R^+$ and a smooth compactly supported vector field $h \colon \R^d \to \R^d$ which satisfy the following properties:
  \begin{enumerate}[i)]
    \item There exists a covering of $\partial\Omega$ by coordinate cylinders which also serve as coordinate cylinders for $\partial\Omega_k$.
    \item The homeomorhpisms $\Lambda_k \colon \partial \Omega \to \partial\Omega_k$ satisfy
      \begin{align*}
        \sup_{Q \in \partial\Omega} |Q - \Lambda_k(Q)| \to 0\quad\text{as } k \to \infty
      \end{align*}
      and $\Lambda_k(P)$ approaches $P$ nontangentially meaning that for all $k \in \N$
      \begin{align*}
        | P - \Lambda_k(P) | < (1 + \beta) \operatorname{dist}(\Lambda_k(P), \partial\Omega)
      \end{align*}
      for some constant $\beta$ depending only on $d$ and the Lipschitz character of $\Omega$.
    \item The normals $\nu_k$ of $\partial\Omega_k$ satisfy $\lim_{k \to \infty} \nu_k(\Lambda_k(P)) = \nu(P)$ a. e. for all $P \in \partial\Omega$
    \item The functions $\vartheta_k$ satisfy $\delta \leq \vartheta_k \leq \delta^{-1}$ for some $\delta > 0$, $\vartheta^k \to 1$ pointwise a. e. and
      \begin{align*}
        \int_E \vartheta_k(Q) \d \sigma(Q) = \int_{\Lambda_k(E)} \d \sigma_k(Q),
      \end{align*}
      where $E \subset \partial \Omega$ is measurable and $\sigma_k$ denotes the surface measure on $\Omega_k$.
    \item The vector field $\alpha$ satisfies $\langle \alpha, \nu_k \rangle \geq c > 0$ a.e. on each $\partial\Omega_k$ where $\nu_k$ denotes the unit inner normal to $\partial\Omega_k$.
  \end{enumerate}
\end{thm}

The next concept we introduce will allow us to talk about boundary values of functions which are defined on $\Omega$ by considering their nontangential behavior.
The first step will be to define nontangential approach regions.
Unfortunately, in the literature there exist at least two different concepts which will be introduced in the next definitions.
  In the following, by a cone we mean an open, circular, truncated cone with only one convex component.

\begin{defn}[Regular family of cones, Verchota]
  \label{defn:regularFamilyOfCones}
  Let $\Omega \subset \R^d$ be a bounded Lipschitz domain.
  If $q \in \partial\Omega$, then $\Gamma(q)$ will denote a cone with vertex $q$ and one component in $\Omega$.
  Assigning to each $q \in \partial\Omega$ one cone $\Gamma(q)$ the family $\{ \Gamma(q) \colon q \in \partial\Omega \}$ will be called \emph{regular} if there exist $x_1, \dots,x_{n_0} \in \partial\Omega$, $\tilde r > 0$ and rotations $\tilde R_{x_1}, \dots, \tilde R_{x_{n_0}}$ such that
  \begin{align*}
    \partial\Omega \subset \bigcup_{i = 1}^{n_0} \{ x_i \} + \tilde R_{x_i}^{-1} D(4 \tilde r / 5),
  \end{align*}
  and such that there exist Lipschitz continuous functions $\tilde \eta_{x_i} \colon \R^{d - 1} \to \R$ such that for all $\tilde r \leq r \leq \nu \tilde r$ with 
  \begin{align*}
    \nu \coloneqq 1 + [ 1 + [10 d(M + 1)]^2]^{1/2}
  \end{align*}
  we have
  \begin{align*}
    \tilde R_{x_i}[\Omega - \{x_i\}] \cap D(r) &= D_{\tilde \eta_{x_i}}(r) \\
    \tilde R_{x_i}[\partial\Omega - \{x_i\}] \cap D(r) &= I_{\tilde \eta_{x_i}}(r).
  \end{align*}
  In addition for all $i$ there exist cones $\alpha_i$, $\beta_i$ and $\gamma_i$ with vertex at the origin and axis along the $x_d$-axis such that
  \begin{align*}
    \alpha_i \subset \overline{\beta_i} \setminus\{0\} \subset \gamma_i
  \end{align*}
  and such that for all $q \in [ \{ x_i\} + \tilde R_{x_i}^{-1} D(4 \tilde r / 5)] \cap \partial\Omega$, we have
  \begin{align*}
    \tilde R_{x_i}^{-1} \alpha_i + \{ q\} \subset \Gamma(q) &\subset \overline{\Gamma(q)} \setminus \{q\} \subset \tilde R_{x_i}^{-1} \beta_i + \{ q\}, \\
    \tilde R_{x_i}^{-1} \gamma_i + \{q\} &\subset [\{ x_i\} + \tilde R_{x_i}^{-1} D(\tilde r)] \cap \Omega.
  \end{align*}
  We will sometimes denote a regular cone as above by $\verCone(q)$.
\end{defn}

For the existence of such families of cones see the Appendix of Verchota \cite{verchota}.

In Verchota cones $\verCone(q)$ we have the properties that for all $\Omega$ there exists a constant $C > 0$ depending only on the Lipschitz character such that for all $q$, $p \in \partial\Omega$ and any $x \in \Gamma_V(p)$ we have that
\begin{align}
  |x - q| &\geq C |x - p| \label{eq:verCone1}\\ 
  |x - q| &\geq C |p - q|.\label{eq:verCone2}
\end{align}
For a proof see Verchota \cite[p. 9f.]{verchota}

\begin{defn}[Nontangential approach region, Shen]
  \label{defn:nontangentialApproachRegion}
   For $\alpha > 1$ and $q \in \partial\Omega$ we define 
   \begin{align*}
     \ShenCone_\alpha(q) \coloneqq \big\{ x \in \Omega \setminus \partial\Omega \colon |x - q| < \alpha \operatorname{dist}(x, \partial\Omega) \big\}
   \end{align*}
   If $\alpha$ is chosen sufficiently large (see Shen \cite{shen2017}) we call $\{ \ShenCone_\alpha(q) \colon q \in \partial\Omega\}$ a \emph{family of nontangential approach regions}.
\end{defn}

  Note that in Shen cones $\ShenCone_\alpha(q)$, we have that for $q, y \in \partial\Omega$ and $x \in \ShenCone_\alpha(q)$
\begin{align}
  \label{eq:shenConeEstimate}
  |q - y| 
  &\leq |q - x| + |x - y| 
  \leq \alpha \operatorname{dist}(x, \partial\Omega) + |x - y|  \nonumber\\
  &\leq (\alpha + 1) |x - y|
\end{align}
where $\alpha$ is the constant from Definition \ref{defn:nontangentialApproachRegion}. 
%It is reasonable to choose $\alpha$ in a way that also condition (ii) from Theorem \ref{thm:smoothApproximation} holds.

Depending on the type of cones used one may introduce similar concepts of nontangential convergence and nontangential maximal functions.

\begin{defn}
  For a function $u$ in $\Omega$ and a fixed family of nontangential approach regions $\{\Gamma_\alpha\}$, we define the nontangential maximal function $(u)_\alpha^*$ by
\begin{align}
  \label{eq:defnNontangMaxFunction}
  (u)_\alpha^*(q) = \sup\big\{ |u(x)| \colon x \in \ShenCone_\alpha(q)\big\}
\end{align}
for $q \in \partial\Omega$.
  For a fixed regular family of cones $\{\verCone(q)\}$ we define the nontangential maximal function $N(u)(q)$ via
  \begin{align*}
    N(u)(q) = \sup\big\{ |u(x)| \colon x \in \verCone(q)\big\}.
  \end{align*}
\end{defn}

Note that Tolksdorf \cite{tolksdorf} and Shen \cite{shen2017} show that the choice of $\alpha$ for the nontangential maximal function as in \ref{eq:defnNontangMaxFunction} does not affect their $p$-norms in an unpredictable way. In fact their $p$-norms for different $\alpha_1$ and $\alpha_2$ stay comparable with a constant only depending on $d$, $\alpha_1$, $\alpha_2$ and the Lipschitz character.
We will therefore for a given bounded Lipschitz domain always assume that $\alpha > 1$ has been choosen big enough such that on the one hand condition (ii) from Theorem \ref{thm:smoothApproximation} is fulfilled and that on the other hand $\alpha$ is large enough such that $\{ \Gamma_\alpha(q) \colon q \in \partial\Omega \}$ is a family of nontangential aproach regions.
In the following we will thus ignore the parameter $\alpha$ in cones and nontangential maximal functions and tacitly assume that it was chosen appropriately.
We further note that the functions $(u)^*$ and $N(u)$ will not be comparable in general, see the discussion in Tolksdorf.

The above mentioned constructions of cones are not limited to cones that lay in the interior of the domain $\Omega$.
In fact the same construction can be carried out for the exterior domain $\R^d \setminus \overline{\Omega}$ yielding cones that lay outside of $\Omega$. While Verchota's cones from Definition \ref{defn:regularFamilyOfCones} can be mirrored along the $x_d = 0$ plane in a suitable local coordinate system, Shen's cones from Definition \ref{defn:nontangentialApproachRegion} have to be modified in a natural way to give cones lying inside of $\R^d \setminus \overline \Omega$, namely
\begin{align*}
  \Gamma_\alpha^{\,\mathrm{ext}}(q) \coloneqq \big\{ x \in \R^d \setminus \overline \Omega \colon |x - q| < \alpha \operatorname{dist}(x, \partial\Omega)\big\}.
\end{align*}

As the name \emph{nontangential approach region} suggests, for functions $u$ living on $\Omega$ or $\R^d \setminus \overline\Omega$ there will be a notion of convergence of function values $u(x)$ as $x$ goes to a point on $p \in \partial\Omega$.
The idea is to restrict the set of directions from which one can approach $p$ by only allowing sequences of points lying in cones $\Gamma(q)$.

\begin{defn}[Nontangential convergence]
  Let $\Omega$ be a bounded Lipschitz domain and $\{\ShenCone(q) \colon q \in \partial\Omega\}$ be a family of nontangential approach regions with its exterior counterpart $\{\ShenCone^{\,\mathrm{ext}}(q) \colon q \in \partial\Omega\}$.
  Let furthermore $u$ be a function on $\R^d \setminus \partial\Omega$ and $f$ a function on $\partial\Omega$.
  We say that $u = f$ \emph{in the sense of nontangential convergence from the inside} if 
  \begin{align*}
    \lim_{\substack{ x \to q \\ x \in \ShenCone(q)}} u(x) = f(q), \quad\text{for a. e. } q \in \partial\Omega
  \end{align*}
  and we say that $u = f$ \emph{in the sense of nontangential convergence from the outside} if 
  \begin{align*}
    \lim_{\substack{ x \to q \\ x \in \ShenCone^{\,\mathrm{ext}}(q)}} u(x) = f(q), \quad\text{for a. e. } q \in \partial\Omega.
  \end{align*}
  If both of the above limits exist and coincide we say that $u = f$ \emph{in the sense of nontangential convergence}.
  %We will sometimes denote a cone coming from a nontangential approach region as above by $\ShenCone(q)$.
\end{defn}

Usually the nontangential limits taken from inside and outside the domain will differ.
For functions $u$ on $\R^d \setminus \partial\Omega$ we will therefore often use the notation $u_+$ to denote the \emph{inner} nontangential limit and $u_-$ for the respective \emph{outer} nontangential limit.

\section{The Stokes Operator}

In this section, we will introduce the Stokes operator on $\Ell^2(\Omega; \C^d)$ and $\Ell^p(\Omega; \C^d)$ for general $p$ and establish a relation to the \emph{Dirichlet problem for the Stokes resolvent system}
\begin{align}
  -\Delta u + \nabla \phi + \lambda u &= f  \quad\text{in } \Omega, \nonumber\\
  \div u &= 0 \quad\text{in } \Omega, \label{eq:stokesResolventSystem} \\
  u &= 0 \quad\text{on } \partial\Omega, \nonumber
\end{align}
where $\lambda \in \Sigma_\theta \coloneqq \{ z \in \C \colon \lambda \neq 0 \text{ and } |\arg(z) | < \pi - \theta \}$ and $\theta \in (0, \pi/2)$.

We beginn by defining the relevant function spaces.
Let $\Omega \subseteq \R^d$ be a bounded Lipschitz domain and $1 < p < \infty$. 
We define
\begin{align*}
  \CC_{c, \sigma}^\infty(\Omega) \coloneqq \{ \varphi \in \CC_c^\infty(\Omega; \C^d) \colon \div(\varphi) = 0 \},
\end{align*}
which can serve as a suitable space of test functions.
We can now close this space in $\Ell^p(\Omega)$ and the Sobolev space $\WW_{1,p}(\Omega)$ which gives
\begin{align*}
  \Ell_\sigma^p \coloneqq \overline{\CC_{c,\sigma}^\infty(\Omega)}_{\Ell^p}
\end{align*}
and
\begin{align*}
  \WW_{0,\sigma}^{1,p}(\Omega) \coloneqq \overline{\CC_{c,\sigma}^\infty(\Omega)}_{\WW^{1,p}},
\end{align*}
respectively.
If $p = 2$, we will use the symbol $\HH_{0,\sigma}^1(\Omega)$ to denote $\WW_{0,\sigma}^{1,2}(\Omega)$ in order to emphasize that this space is a Hilbert space.

In order to define the Stokes operator, we introduce the following sesquilinear form
\begin{align*}
  \aform \colon \HH_{0,\sigma}^1(\Omega) \times \HH_{0,\sigma}^1(\Omega) \to \C, \quad (u,v) \mapsto \int_\Omega \nabla u \cdot \overline{\nabla v} \d x.
\end{align*}
Note that for $u \in \HH_{0,\sigma}^1(\Omega)$ the gradient $\nabla u$ is a matrix and an element of the space $\Ell^2(\Omega; \C^{d \times d})$.

\begin{defn}
  \label{defn:stokes}
  The \emph{Stokes operator} $A_2$ on $\Ell_\sigma^2(\Omega)$ is given via
  \begin{align*}
    \Dom(A_2) &\coloneqq \Big\{ u \in \HH_{0,\sigma}^1(\Omega) \colon \exists ! f \in \Ell^2_\sigma(\Omega) \text{ s.t. } \forall v \in \HH_{0,\sigma}^1(\Omega) \colon \aform(u,v) = \int_\Omega f \cdot \overline v \d x \Big\} \\
    A_2 u &\coloneqq f,
  \end{align*}
  where $u \in \Dom(A_2)$ and $f$ comes from the definition of the domain.
\end{defn}

The following theorem from Mitrea and Monniaux \cite[Thm 4.7]{mitreaMonniaux} shows that our definition of the Stokes operator and the one used in Shen's paper \cite{Shen2012} coincide. 
Another advantage of this characterization is the immediate link of the Stokes operator to the Stokes system.

\begin{thm}
  If $\Omega \subseteq \R^d$, $d \geq 2$, is a bounded Lipschitz domain and $A_2$ is the Stokes operator on $\Ell^2_\sigma(\Omega)$ then
  \begin{align*}
    \Dom(A_2) = \big\{ u \in H^1_{0,\sigma}(\Omega) \colon \exists \pi \in \Ell^2(\Omega) \text{ s.t. } -\Delta u + \nabla \pi \in \Ell^2_\sigma(\Omega) \big\},
  \end{align*}
  where the expression $\Delta u + \nabla \pi \in\Ell^2_\sigma(\Omega)$ needs to be understood in the distributional sense.
  For $u \in \Dom(A_2)$ and the corresponding pressure $\pi$ we have
  \begin{align*}
    A_2 u = -\Delta u + \nabla \pi.
  \end{align*}
\end{thm}

The following proposition summarizes some facts about the Stokes operator on $\Ell^2_\sigma(\Omega)$. A proof can be found in Tolksdorf \cite{tolksdorf}.

\begin{prop}
  Let $\Omega \subset \R^d$ be a bounded Lipschitz domain and $A_2$ the Stokes operator as in Definition \ref{defn:stokes}. Then we have
  \begin{enumerate}[a)]
    \item $A_2$ is closed with dense domain. Furthermore $0 \in \rho(A_2)$.
    \item $\sigma(A) \subset [0,\infty)$ and for all $\theta \in (0,\pi]$ there exists $C > 0$ such that
      \begin{align}
        \label{eq:resolventEstimateL2}
        \| \lambda(\lambda + A)^{-1} \|_{\Li(\Ell^2_\sigma)} \leq C, \quad\text{for all } \lambda \in \C \setminus \overline \Sigma_\theta.
      \end{align}
      In particular $-A_2$ generates a bounded analytic semigroup on $\Ell^2_\sigma(\Omega)$.
  \end{enumerate}
\end{prop}

With these results at hand we can now give a quick recap of the solution theory to \eqref{eq:stokesResolventSystem}.
Let $f \in \Ell^2_\sigma$ and $\lambda \in \Sigma_\theta$, $\theta \in (0, \pi/2)$.
By the previous proposition we know that there exists a unique $u \in \Dom(A_2) \subseteq \HH^1_{0,\sigma}(\Omega)$ and some $\pi \in \Ell^2(\Omega)$ such that
\begin{align*}
  -\Delta u + \nabla \pi + \lambda u = A_2 u = f.
\end{align*}
For general $f = \Ell^2(\Omega)$ we use the \emph{Helmholtz projection} $\PP_2$ to get
\begin{align*}
  \Delta u + \nabla \pi + \lambda u + (I - \PP_2) f = \PP_2 f + (I - \PP_2) f = f,
\end{align*}
where $u$ and $\pi$ now correspond to $\PP_2 f \in \Ell^2_\sigma(\Omega)$. On bounded Lipschitz domains the orthogonal complement to $\PP_2[\Ell^2(\Omega; \C^d)] = \Ell^2_\sigma(\Omega)$ is characterized via
\begin{align*}
  \Ell^2_\sigma(\Omega)^\perp = \big\{ f \in \Ell^2(\Omega; \C^d) \colon f = \nabla \phi, \text{ for some } \phi \in \Ell^2(\Omega) \big\}.
\end{align*}
A proof of this fact can be found in the book of Sohr \cite[Lemma 2.5.3]{sohr}.
Using this result we find $g \in \Ell^2(\Omega)$ such that $\nabla g = (I - \PP_2)f$ in the sense of distributions and we see that
\begin{align*}
  -\Delta u + \nabla( \pi + g ) + \lambda u = f.
\end{align*}
Furthermore we may deduce from the resolvent estimate \eqref{eq:resolventEstimateL2} that the solution $u$ which apparently is not affected by the additional part $(I-\PP_2)f$ fulfills the inequality
\begin{align*}
  |\lambda|^{-1} \| u\|_{\Ell^2(\Omega; \C^d)} = |\lambda|^{-1} \|(A_2 + \lambda)^{-1} \PP_2 f \|_{\Ell^2(\Omega; \C^d)} \leq C \| f \|_{\Ell^2(\Omega; \C^d)},
\end{align*}
where $C$ depends only on $\theta$.

In order to develop an $\Ell^p$-theory for system \eqref{eq:stokesResolventSystem}, one way is to study the Stokes operator on subspaces of $\Ell^p(\Omega; \C^d)$. 

\begin{defn}
  \label{defn:stokeslp}
  Let $\Omega \subseteq \R^d$, $d \geq 2$ be a bounded Lipschitz domain and $1 < p < \infty$.
  If $p > 2$ we define the Stokes operator $A_p$ via its \emph{part of } $A_2$ \emph{in } $\Ell^p_\sigma(\Omega)$.
  \begin{align*}
    \Dom(A_p) &\coloneqq \Big\{ u \in \Dom(A_2) \cap \Ell^p_\sigma(\Omega) \colon A_2 u \in \Ell^p_\sigma(\Omega) \Big\} \\
    A_p u &\coloneqq A_2 u, \quad u \in \Dom(A_p).
  \end{align*}
  If $p < 2$ and $A_2$ is closable in $\Ell^p_\sigma(\Omega)$, then $A_p \coloneqq \overline{A_2}$.
\end{defn}

\chapter{Estimating Fundamental Solutions}

The purpose of this section is to study fundamental solutions of the Stokes resolvent problem and to deduce related estimates which will be crucial for  the next chapters.

Let $\lambda = r \e^{\ii \tau}$ with $0 < r < \infty$ and $-\pi + \theta < \tau < \pi - \theta$ and set $k = \sqrt{r} \e^{\ii(\pi + \tau)/2}$.
Then 
$$k^2 = -\lambda\quad\text{ and }\quad \frac{\theta}{2} < \arg(k) < \pi - \frac{\theta}{2}$$
as
\begin{align*}
  \arg(k) = \frac{\pi + \tau}{2} &> \frac{\pi}{2} + \frac{-\pi + \theta}{2} = \frac{\theta}{2} 
  \intertext{on the one hand and}
  &< \frac{\pi}{2} + \frac{\pi - \theta}{2} = \pi  - \frac{\theta}{2}
\end{align*}
on the other hand.
This gives rise to the following estimate
\begin{align}
  \label{eq:imaginaryPartEstimate}
  \Im(k) > \sin(\theta/2) \sqrt{|\lambda|}.
\end{align}
Indeed, we have $$\Im( k) = \sqrt{r} \sin( \frac{\pi + \tau}{2} ) = \sqrt{|\lambda|} \sin( \frac{\pi + \tau}{2} )\quad\text{and}\quad \frac{\theta}{2} < \frac{\pi + \tau}{2} < \pi - \frac{\theta}{2}$$
which gives for $\tau$ with $\frac{\pi + \tau}{2} \leq \frac{\pi}{2}$ that $\sin(\frac{\pi + \tau}{2}) \geq \sin(\frac{\theta}{2} )$ and for $\tau$ with $\frac{\pi + \tau}{2} > \frac{\pi}{2}$ that $\sin(\frac{\pi + \tau}{2}) > \sin(\pi - \frac{\theta}{2} ) = \sin(\frac{\theta}{2} )$.

Before diving into fundamental solutions of the Stokes resolvent problem, we will first consider a fundamental solution for the (scalar) Helmoltz equation in $\R^d$
$$
-\Delta u + \lambda u = 0.
$$
One fundamental solution with pole at the origin is given by
\begin{align}
  \label{eq:definitionFundamentalHelmholtz}
  G(x; \lambda) = \frac{\ii}{4 ( 2\pi )^{\frac{d}{2} - 1}} \cdot \frac{1}{|x|^{d - 2}} \cdot (k |x|)^{\frac{d}{2} - 1} H_{\frac{d}{2} - 1}^{(1)} (k|x|),
\end{align}
where $H_{\nu}^{(1)}(z)$ is the Hankel function of the first kind which can be written as
\begin{align}
  \label{eq:integralRepresentationHankel}
  H_\nu^{(1)}(z) = \frac{2^{\nu + 1} \e^{\ii(z - \nu \pi)} z^\nu}{\ii \sqrt{\pi} \Gamma(\nu + \frac{1}{2})} \int_0^\infty \e^{2 z \ii s} s^{\nu - \frac{1}{2}} (1 + s)^{\nu - \frac{1}{2}} \d s.
\end{align}
This formula holds for $\nu > -\frac{1}{2}$ and $0 < \arg(z) < \pi$.
We will usually set 
\begin{align*}
\nu = \frac{d}{2} - 1 \quad\text{and}\quad z = k|x|. 
\end{align*}
Therefore, the formula \eqref{eq:integralRepresentationHankel} will hold for all dimensions $d \geq 2$ and all $x \in \R^d$.
In the case $d = 2$ formula \eqref{eq:definitionFundamentalHelmholtz} simplifies to 
\begin{align}
  \label{eq:2dDefinitionFundamentalHelmholtz}
  G(x;\lambda) = \frac{\ii}{4} H_{0}^{(1)}(k|x|),
\end{align}
in the case $d = 3$ one has an even easier formula, namely
\begin{align}
  \label{eq:3dDefinitionFundamentalHelmholtz}
  G(x; \lambda) = \frac{\e^{\ii k |x|}}{4 \pi |x|}.
\end{align}

Our first estimate is concerned with derivatives of the fundamental solution for the (scalar) Helmholtz equation.

\begin{lem}
  \label{lem:estimateHelmholtzDerivatives}
  Let $\lambda \in \Sigma_\theta$.
  Then
  \begin{align}
    \label{eq:estimateHelmholtzDerivatives}
    |\nabla_x^l G(x; \lambda)| \leq \frac{ C_l \e^{-c \sqrt{|\lambda|} |x|}}{|x|^{d - 2 + l}}
  \end{align}
  for any integer $l \geq 0$ if $d \geq 3$ and for $l \geq 1$ if $d = 2$.
  Here, $c > 0$ depends only on $\theta$ and $C_l$ depends only on $d$, $l$ and $\theta$.
\end{lem}

\begin{proof}
  We start with the case $l = 0$ and $d \geq 3$.
  Let $\Im(z) > 0$ and $\nu - \frac{1}{2} \geq 0$.
  Then \eqref{eq:integralRepresentationHankel} gives
  \begin{align*}
    |H_\nu^{(1)}(z)|
    &\leq C \e^{-\Im(z)} |z|^\nu \int_0^\infty \e^{-2s \Im(z)} s^{\nu - \frac{1}{2}} (1 + s)^{\nu - \frac{1}{2}} \d s.
  \end{align*}
  Since by the substitution rule
  \begin{align*}
    \e^{\frac{-\Im(z)}{2}} \int_0^\infty \e^{-s \Im(z)} s^{\nu - \frac{1}{2}} ( 1 + s )^{\nu - \frac{1}{2}} \d s
    &= \int_0^\infty \e^{-\Im(z) ( s + \frac{1}{2})} s^{\nu - \frac{1}{2}} ( 1 + s )^{\nu - \frac{1}{2}} \d s \\
    &= \int_{\frac{1}{2}}^\infty \e^{-\Im(z) t } (t^2 - \frac{1}{4} ) ^{\nu - \frac{1}{2}} \d t \\
    &\leq \int_0^\infty e^{-\Im(z) t} t^{2\nu - 1} \d t \\
    &= \int_0^\infty e^{-u} (\frac{ u}{\Im(z)} )^{2 \nu - 1} (\Im(z))^{-1} \d u \\
    &= (\Im(z))^{-2 \nu} \int_0^\infty \e^{-u} u^{2 \nu - 1} \d u,
  \end{align*}
  we can estimate
  \begin{align*}
    |z|^\nu |H_\nu^{(1)}(z)| \leq C |z|^{2 \nu} |\Im(z)|^{-2\nu} \e^{-\frac{\Im(z)}{2}},
  \end{align*}
  which for $z = k|x|$ gives
  \begin{align}
    \label{eq:zHEstimate}
    |kx|^\nu |H_\nu^{(1)}(k |x|)| \leq \sin(\theta/2)^{-2\nu} \e^{-\frac{1}{2} \sin(\theta/2) \sqrt{|\lambda|} |x|},
  \end{align}
  where we used \eqref{eq:imaginaryPartEstimate} to estimate
  \begin{align*}
    (|kx|)^{2\nu} \cdot |\Im(k|x|)|^{-2\nu} 
    = |\lambda|^\nu \cdot |\Im(k)|^{-2\nu} 
    \leq \sin(\theta/2)^{-2\nu}.
  \end{align*}
  Using \eqref{eq:definitionFundamentalHelmholtz}, we estimate for $d \geq 3$ setting $\nu = \frac{d}{2} - 1$
  \begin{align*}
    |G(x; \lambda)| 
    \leq C  |x|^{2 - d} e^{-c \sqrt{|\lambda|} |x|}.
  \end{align*}
  This gives the estimate for $l = 0$ and $d \geq 3$.

  Using the relation
  \begin{align*}
    \frac{\d{}}{\d z} \{ z^{-\nu} H_\nu^{(1)}(z) \} = -z^{-\nu} H_{\nu + 1}^{(1)} (z),
  \end{align*}
  we can inductively establish the estimate \eqref{eq:estimateHelholtzDerivatives} for $l \geq 1$ and $d \geq 2$:
  For $1 \leq j \leq d$, we calculate
  \begin{align*}
    |\nabla_x^{} G(x; \lambda)|
    &\leq C  \cdot \big\{|x|^{1 - d} \cdot (k|x|)^{\frac{d}{2} - 1} H_{\frac{d}{2} - 1}^{(1)}(k|x|) - |x|^{d - 2} (k|x|)^{\frac{d}{2} - 1} H_{\frac{d}{2}}^{(1)}(k|x|) k\big\} \\
    &\leq C \cdot |x|^{1 - d}\big\{(k|x|)^{\frac{d}{2} - 1} H_{\frac{d}{2} - 1}^{(1)}(k|x|) -  (k|x|)^{\frac{d}{2}} H_{\frac{d}{2}}^{(1)}(k|x|) \big\},
  \end{align*}
  where the first summand does not arise in the case $d = 2$ due to \eqref{eq:2dDefinitionFundamentalHelmholtz}.
  The terms in the bracket can now be estimated individually by \eqref{eq:zHEstimate}.
  The extension of this proof to orders of differentiation $l \geq 2$ is straightforward using the Leibniz product rule for higher derivatives.
\end{proof}

In the derivation of the next estimates we will use the following useful interior estimate which we write down for further use.

\begin{lem}
  \label{lem:interiorEstimatePoisson}
  Let $w$ be a solution to $\Delta w = f$ in $\BB(x,r)$. Then
  \begin{align}
    \label{eq:interiorEstimatePoisson}
    |\nabla^l w(x)| \leq C r^{-l} \sup_{\BB(x,r)} |w| + C \max_{0 \leq j \leq l - 1} \sup_{\BB(x,r)} r^{j - l + 2} |\nabla^j f|.
  \end{align}
\end{lem}

\begin{proof}
  Todo
\end{proof}

We will need the following asymptotic expansions for the function $z^\nu H_\nu^{(1)}(z)$ in $\C \setminus (-\infty, 0]$.
\begin{align}
  z^{\nu}H_\nu^{(1)}(z) &= \frac{2^\nu \Gamma(\nu)}{\ii \pi} + \frac{\ii}{\pi} z^2 \log(z) + \omega z^2 + O(|z|^4 |\log(z)) \text{ if } d = 4, \label{eq:asymptoticd4}\\
  z^{\nu} H_{\nu}^{(1)}(z) &= \frac{2^\nu \Gamma(\nu)}{\ii\pi} + \frac{2^\nu \Gamma(\nu - 1)}{4 \pi \ii} z^2 + \omega z^3 + O(|z|^4) \text{ if } d = 5, \label{eq:asymptoticd5}\\
  z^\nu H_\nu^{(1)}(z) &= \frac{2^\nu \Gamma(\nu)}{\ii \pi} + \frac{2^\nu \Gamma(\nu - 1)}{4\pi \ii}z^2 + O(|z|^4 |\log z|) \text{ if } d = 6,\label{eq:asymptoticd6} \\
  z^\nu H_\nu^{(1)}(z) &= \frac{2^\nu \Gamma(\nu)}{\ii \pi} + \frac{2^\nu \Gamma(\nu - 1)}{4 \pi \ii} z^2 + O(|z|^4) \text{ if } d \geq 7.\label{eq:asymptoticd7}
\end{align}

The derivation of these asymptotic expansions is based on asymptotic expansions of the Bessel functions of the first and the second kind and can be found in Tolksdorf \cite{tolksdorfDiss}.

We will denote the fundamental solution for $-\Delta$ in $\R^d$ with pole at the origin by
\begin{align*}
  G(x;0) \coloneqq \begin{cases} -\frac{1}{2\pi} \log(|x|), &\quad\text{for } d = 2, \\c_d \frac{1}{|x|^{d - 2}}, &\quad\text{for } d > 2,  \end{cases}
\end{align*}
where
\begin{align*}
  c_d = \frac{1}{(d - 2) \omega_d}, \quad\text{with}\quad \omega_d = \frac{2\pi^{\frac{d}{2}}}{\Gamma(\frac{d}{2})} = |\mathbb{S}^{d - 1}|.
\end{align*}
Since
\begin{align*}
  (d - 2) \omega_d 
  = 2 ( \frac{d}{2} - 1) \frac{2 \pi^{\frac{d}{2}}}{\Gamma(\frac{d}{2})}
  =  2 (\frac{d}{2} - 1) \frac{2 \pi^{\frac{d}{2}}}{\Gamma(\frac{d}{2} - 1)(\frac{d}{2} - 1)}
  = \frac{4 \pi^{\frac{d}{2}}}{\Gamma(\frac{d}{2} - 1)},
\end{align*}
we will also sometimes use
\begin{align*}
  c_d \coloneqq \frac{\Gamma(\frac{d}{2} - 1)}{4 \pi^{\frac{d}{2}}}.
\end{align*}
The leading coefficient of the asymptotic expansions for $d \geq 3$ will be denoted as
\begin{align}
  \label{eq:Defnad}
  a_d \coloneqq \frac{2^{\frac{d}{2} - 1} \Gamma(\frac{d}{2} - 1)}{\ii \pi}.
\end{align}
The coefficients $a_d$ and $c_d$ are related in the following way
\begin{align*}
  c_d 
  %= \frac{\Gamma(\frac{d}{2} - 1) 2^\nu}{\ii \pi} 
  %=\frac{4 (2 \pi)^{\frac{d}{2} - 1}}{\ii} \frac{\Gamma(\frac{d}{2} - 1)}{4 \pi^{\frac{d}{2}}} 
  =\frac{\ii} {4 (2 \pi)^{\frac{d}{2} - 1}}\; a_d 
\end{align*}
This allows us to write for $d \geq 3$
\begin{align}
  \label{eq:HelmholtzLaplaceDifference}
  G(x;\lambda) - G(x; 0) = \frac{i}{4(2\pi)^{\frac{d}{2} - 1}} \cdot \frac{1}{|x|^{d - 2}} \Big\{ (k|x|)^{\frac{d}{2} - 1} H_{\frac{d}{2} - 1}^{(1)}(k|x|) - a_d \Big\}.
\end{align}

The following lemma is helps us to estimate derivatives of \eqref{eq:HelmholtzLaplaceDifference} and the $2$-dimensional counterpart.

\begin{lem}
  \label{lem:HelmholtzLaplaceDifference}
  Let $\lambda \in \Sigma_\theta$.
  Then
  \begin{align}
    \label{eq:HelmholtzLaplaceDifferenceEstimate}
    |\nabla_x^l \{ G(x; \lambda) - G(x; 0) \}| \leq C |\lambda| |x|^{4 - d - l}
  \end{align}
  if $d \geq 5$ and $l \geq 0$, where $C$ depends only on $d$, $l$ and $\theta$.
  If $d = 3$ or $4$, estimate \eqref{eq:HelmholtzLaplaceDifferenceEstimate} holds for $l \geq 1$ and if $d = 2$ the estimate holds for $l \geq 3$.
\end{lem}

\begin{proof}
  \begin{enumerate}[(a)]
    \item In this part we will show that the desired estimates \eqref{eq:HelmholtzLaplaceDifferenceEstimate} and \eqref{eq:HelmholtzLaplaceDifferenceEstimate} hold if we assume that $|\lambda| |x|^2 > \frac{1}{2}$.
    In this case, Lemma \ref{lem:estimateHelmholtzDerivatives} gives
    \begin{align*}
      |\nabla_x^l \{ G(x; \lambda) - G(x; 0) \} |
      \leq C \Big\{ \frac{\e^{-c \sqrt{|\lambda|} |x|}}{|x|^{d - 2 + l}} + \frac{1}{|x|^{d - 2 + l}} \Big\} \leq C \frac{|\lambda|}{|x|^{d - 4 + l}},
    \end{align*}
    where $C$ depends only on $d$, $l$ and $\theta$.
    Therefore, for the remaining proof we will suppose $|\lambda||x|^2 \leq \frac{1}{2}$.
  \item In this step, we show that we can restrict ourselves to proving \eqref{eq:HelmholtzLaplaceDifferenceEstimate} in three cases: (1) $d \geq 5$ and $l = 0$; (2) $d = 3$ or $4$ and $l = 1$; (3) $d = 2$ and $l = 3$.
    
    Suppose \eqref{eq:HelmholtzLaplaceDifferenceEstimate} holds in case (1) and let $l > 1$.
    If we set $ w(x) = G(x;\lambda) - G(x;0) $, we have $\Delta_x w = \lambda G(x; \lambda)$ in $\R^d \setminus \{0\}$.
    For $f = \lambda G(x; \lambda)$ estimate \eqref{eq:interiorEstimatePoisson} now gives
    \begin{align*}
      |\nabla^l w(x) |
      &\leq C r^{-l} \sup_{\BB(x,r)} |w| + C \max_{0 \leq j \leq l - 1} \sup_{\BB(x,r)} r^{j - l + 2} |\nabla^j f| \\
      &\leq C r^{-l} \sup_{y \in \BB(x,r)} |\lambda| |y|^{4 - d} + C \sum_{j = 0}^{l - 1} \sup_{y \in \BB(x,r)} r^{j - l + 2} |\lambda| |y|^{2 - d - j} \\
      &= C r^{-l} |\lambda| |x - r \frac{x}{|x|} |^{4 - d} + C \sum_{j = 0}^{l - 1} r^{j - l + 2} |\lambda| |x - r \frac{x}{|x|} |^{2 - d - j},
    \end{align*}
    for all $0 < r < |x|$, where we used \eqref{eq:HelmholtzLaplaceDifferenceEstimate} with $l = 1$ for the first summand and \eqref{eq:estimateHelmholtzDerivatives} to estimate the second summand.
    Setting $r = \frac{|x|}{2}$ now gives
    \begin{align*}
      |\nabla^l w(x)| 
      &\leq C |\lambda |x|^{-l} |x|^{4 - d} + C \sum_{j = 0}^{l - 1} |x|^{j - l + 2} |\lambda| |x|^{2 - d - j} \\
      &\leq C |\lambda| |x|^{4 - d - l}.
    \end{align*}

    The proof for case (2) is completely analogous if one sets $w(x) = \nabla_x (G(x; \lambda) - G(x; 0))$ and $f(x) = \lambda \nabla_x G(x; \lambda)$.
    Also case (3) is proven in a similar fashion. We will give the proof for the sake of completeness.

    For $w$ and $f$ as in case (2) by \eqref{eq:HelmholtzLaplaceDifferenceEstimate} we get
    \begin{align*}
      |\nabla^l w(x) |
      &\leq C r^{-l} \sup_{\BB(x,r)} |w| + C \max_{0 \leq j \leq l - 1} \sup_{\BB(x,r)} r^{j - l + 2} |\nabla^j f| \\
      &\leq C r^{-l} \sup_{y \in \BB(x,r)} |\lambda| |y|(|\log |\lambda| |y|^2 | + 1 ) + C \sum_{j = 0}^{l - 1} \sup_{y \in \BB(x,r)} r^{j - l + 2} |\lambda| |y|^{- j - 1} \\
      &\leq S_1 + S_2,
      \intertext{wheras}
      S_1 
      &\leq C r^{-l} |\lambda| |x + r \frac{x}{|x|} | ( |\log |\lambda| |x - r \frac{x}{|x|}|^2 | + |\log( |\lambda||x + r \frac{x}{|x|}|^2 )|+ 1) \\
      &\leq C  |\lambda| |x|^{1-l}(|\log( |\lambda| |x|^2)| + 1)
      \intertext{
        if we choose $r = \frac{|x|}{2}$. For $S_2$ we calculate as before, using estimate \eqref{eq:estimateHelmholtzDerivatives}
  }
      S_2
      &\leq C \sum_{j = 0}^{l - 1} C |x|^{j - l + 2} |\lambda| |x|^{-j - 1} \\
      &\leq C |\lambda| |x|^{1 - l}.
    \end{align*}

  \item In this step we prove \eqref{eq:HelmholtzLaplaceDifferenceEstimate} for $d \geq 5$ and $l = 0$. 
    First, note that for the functions
    \begin{align*}
      &g(x) \coloneqq (k|x|)^{\frac{d}{2} - 1} H_{\frac{d}{2} - 1}^{(1)}(k|x|), \quad g(0) = a_d, \\
      &h(z) \coloneqq z^{\frac{d}{2} - 1} H_{\frac{d}{2} - 1}^{(1)}(z), \quad h(0) = a_d,
    \end{align*}
    the mean value theorem yields the estimate
    \begin{align*}
      |g(x) - g(0)| \leq |x| \sup_{\BB(0,|x|)} |\nabla f(y)| \leq |x| |k| \sup_{\BB(0,|x|)} \frac{\d{}}{\d z} |h(k|x|)|.
    \end{align*}
    Due to \eqref{eq:HelmholtzLaplaceDifference} we estimate
    \begin{align}
      |G(x; \lambda) - G(x; 0) |
      &\leq C |x|^{2 - d} \cdot |k| |x| \max_{\substack{|z| \leq |k| |x| \\ \Im(z) > 0}} \big| \frac{\d{}}{\d z} \big\{ z^{\frac{d}{2} - 1} H_{\frac{d}{2} - 1}(z) \big\} \big| \nonumber\\
      &= C |x|^{2 - d} \cdot |k| |x| \max_{\substack{|z| \leq |k||x| \\ \Im(z) > 0}} |z^{\frac{d}{2} - 1} H_{\frac{d}{2} - 2}^{(1)}(z) |, \label{eq:HelmholtzLaplaceDifferenceGt5}
    \end{align}
    where for the last equality we used the relation
    \begin{align}
      \frac{\d{}}{\d z} \big\{ z^\nu H_\nu^{(1)}(z) \big\} = z^\nu H_{\nu - 1}^{(1)}(z).
    \end{align}
    Since the asymptotic expansions yield that $|z^{\nu} H_\nu^{(1)}(z)| \leq C_\nu$ for $\nu > 0$ and $|z| \leq 1$ with $\Im(z) > 0$ it follows from \eqref{eq:HelmholtzLaplaceDifferenceGt5} that
    \begin{align*}
      |G(x;\lambda) - G(x; 0)| \leq C |x|^{2 - d} \cdot |k| |x| \cdot |k| |x| \max_{\substack{|z| \leq |k||x| \\ \Im(z) > 0}} |z^{\frac{d}{2} - 2} H_{\frac{d}{2} - 2}^{(1)}(z) |
      \leq C |\lambda| |x|^{4 - d}
    \end{align*}
    
  \item Now we consider the case $d = 4$ and $l = 1$.
    The asymptotic expansion \eqref{eq:asymptoticd4} gives that
    \begin{align}
      \label{eq:mwt4d}
      \Big| \frac{\d{}}{\d z} \bigg\{ \frac{z H_1^{(1)}(z) - a_4}{z^2} \bigg\} \Big| \leq C |z|^{-1}
    \end{align}
    for all $|z| \leq \frac{1}{2}$ with $\Im(z) > 0$.
    Since
    \begin{align*}
      \frac{G(x; \lambda) - G(x; 0)}{\lambda} = - \frac{C (z H_1^{(1)}(z) - a_4)}{z^2},
    \end{align*}
    where $z = k |x|$, it follows from \eqref{eq:mwt4d} and the mean value theorem that
    \begin{align*}
      \Big|\frac{\nabla_x \{ G(x; \lambda) - G(x; 0) \}}{\lambda} \Big|
      \leq C |k| \frac{\d{}}{\d z} \Big( \frac{z H_1^{(1)}(z) - a_4}{z^2} \Big) \Big|_{z = k|x|}
      \leq C |k| |k|^{-1} |x|^{-1}.
    \end{align*}
    Which after rearrangement of the involved terms gives the claim.

  \item For the case $d = 3$ and $l = 1$, equation \eqref{eq:HelmholtzLaplaceDifference} reads
    \begin{align*}
      G(x;\lambda) - G(x; 0) = \frac{\e^{\ii k |x|}}{4\pi |x|} - \frac{c_3}{|x|} = \frac{\e^{\ii k|x|} - 1}{4 \pi |x|}.
    \end{align*}

    Now we calculate
    \begin{align*}
      \frac{\partial}{\partial x_j} \Big\{ \frac{\e^{\ii k|x|} - 1}{|x|} \Big\} 
      &= \frac{\partial}{\partial x_j} \Big\{ \frac{\e^{\ii k |x|} - 1 - \ii k |x|}{|x|} \Big\} 
      = \frac{\partial}{\partial x_j} \Big\{ \sum_{n = 2}^\infty \frac{(\ii k |x|)^n}{n!} \cdot \frac{1}{|x|} \Big\} \\
      &= \sum_{n = 2}^\infty \frac{(\ii k)^n}{n!} (n - 1) \cdot \frac{x_j}{|x|} |x|^{n - 2} 
    \end{align*}
    which in turn implies 
    \begin{align*}
      \Big| \frac{\partial}{\partial x_j} \Big\{ \frac{\e^{\ii k |x|} - 1}{|x|} \Big\} \Big|
      &\leq |\lambda| \sum_{n = 2}^\infty \frac{n - 1}{n!} |k|^{n - 2} |x|^{n - 2} 
      \leq C \, |\lambda|
    \end{align*}
    since $|\lambda||x| < \frac{1}{2}$.

  \item For the last case $d = 2$ and $l = 3$, we will directly calculate the estimate using the asymptotic expansion of $H_0^{(1)}(z)$ with $z = k|x|$:
    \begin{align*}
      H_0^{(1)}(z) 
      &= J_0(z) + \ii Y_0(z) \\
      &= \sum_{l = 0}^\infty \frac{(-1)^l}{(l!)^2 4^l} z^{2l} \big( 1 - \frac{2\ii \log(2)}{\pi} \big) 
      - \frac{2 \ii}{\pi} \sum_{l = 0}^\infty \frac{(-1)^l}{(l!)^2 4^l} \psi(l + 1) \cdot z^{2l} \\
      &\quad + \frac{2\ii}{\pi} \sum_{l = 0}^\infty \frac{(-1)^l}{(l!)^2 4^l} z^{2l} \log(z)
    \end{align*}
    The first complex derivative of $H_0^{(1)}(z)$ reads
    \begin{align*}
      \frac{\d{}}{\d z} H_0^{(1)}(z) = 
    \end{align*}

  \end{enumerate}
\end{proof}

\begin{rem}
  In the situation of Lemma \ref{lem:HelmholtzLaplaceDifference} one can show by considering the asymptotic expansions and if $|\lambda| |x|^2 \leq (1/2)$ that 
  \begin{align*}
    |G(x; \lambda) - G(x; 0) | \leq \begin{cases}
      C \sqrt{|\lambda|} \quad&\text{if } d = 3, \\
      C |\lambda| \{ |\log(|\lambda| |x|^2) | + 1 \} \quad&\text{if } d = 4.
    \end{cases}
  \end{align*}
  Also using the asymptotic expansions it can be shown that if $d = 2$
  \begin{align*}
    |\nabla_x^l \{ G(x; \lambda) - G(x; 0) \} | \leq C |\lambda| |x|^{2 - l} \{ |\log(|\lambda| |x|^2 ) | + 1 \}
  \end{align*}
  for $l \in \{1, 2\}$.  
\end{rem}

  The fundamental solution to the (scalar) Helmholtz equation and the Laplace equation form the main ingredient for the following matrix of fundamental solutions $(\Gamma(x;\lambda) = (\Gamma_{\alpha\beta}(x;\lambda))_{d \times d}$ with pole at the origin to the Stokes resolvent problem with $\lambda \in \Sigma_\theta$:
  \begin{align}
    \label{eq:fundamentalMatrixStokes}
    \Gamma_{\alpha\beta}(x;\lambda) = G(x; \lambda) \delta_{\alpha\beta} - \frac{1}{\lambda} \frac{\partial^2}{\partial x_\alpha \partial x_\beta} \{ G(x; \lambda) - G(x; 0) \}, \quad \alpha,\beta = 1,\dots,d.
  \end{align}
  Having the formula at sight, the following observations are obvious:
  \begin{align*}
    \Gamma_{\alpha\beta}(x; \lambda) = \Gamma_{\beta\alpha}(x; \lambda), \quad 
    \overline{\Gamma_{\alpha\beta}(x; \lambda)} = \Gamma_{\alpha\beta}(x; \bar\lambda)
    \quad\text{and}\quad
    \Gamma_{\alpha\beta}(x; \lambda) = \Gamma_{\alpha\beta}(-x; \lambda).
  \end{align*}
  For the pressure, we define the vector of fundamental solutions
  \begin{align}
    \label{eq:fundamentalVectorPressure}
      \Phi_\beta(x) = -\frac{\partial}{\partial x_\beta} \{ G(x; 0) \} = \frac{x_\beta}{\omega_d |x|^d}.
  \end{align}
  We note that $\Phi_\beta(x) = \Phi_\beta(-x)$.

  Using the fact that $\Delta_x G(x; \lambda) = \lambda G(x; \lambda)$ in $\R^d \setminus \{0\}$, one can see that on $\R^d \setminus \{0\}$ and for all $1 \leq \beta \leq d$
  \begin{align}
    \label{eq:solutionStokesSystem}
    \begin{cases}
      (-\Delta_x + \lambda) \Gamma_{\alpha\beta}(x;\lambda) + \frac{\partial}{\partial x_\alpha} \{ \Phi_\beta(x) \} &= 0 \quad\text{for } 1 \leq \alpha \leq d, \\
      \frac{\partial}{\partial x_\alpha} \{ \Gamma_{\alpha\beta} (x; \lambda) \} &= 0.
    \end{cases}
  \end{align}
  Note that in the last equation the summation convention was used.

  We now keep up to the spirit of this exhausting section by proving further estimates, this time for the fundamental solutions to the Stokes resolvent problem.

  \begin{thm}
    \label{thm:fundamentalMatrixEstimate}
    Let $\lambda \in \Sigma_\theta$.
    Then for any $d \geq 3$ and $l \geq 0$
    \begin{align}
      \label{eq:fundamentalMatrixEstimate}
      | \nabla_x^l \Gamma(x; \lambda) | &\leq \frac{C}{(1 + |\lambda||x|^2) |x|^{d - 2 + l}} 
    \end{align}
      where $C$ depends only on $d$, $l$ and $\theta$. For $d = 2$ and $l \geq 1$ the same estimate holds.
  \end{thm}

  \begin{proof}
    Let $|\lambda| |x|^2 > \frac{1}{2}$. 
    Then there exist constants $C_a$, $C_b$, $C_c$ such that
    \begin{align*}
      \e^{-c \sqrt{|\lambda|} |x|} ( 1 + |\lambda| |x|^2) &\leq C_a \\
      1 &\leq \frac{C_b |\lambda| |x|^2}{1 + |\lambda| |x|^2} \\
      \e^{-c \sqrt{|\lambda|} |x|} &\leq \frac{ C_c |\lambda| |x|^2}{1 + |\lambda| |x|^2},
    \end{align*}
    where $c$ is the constant from Lemma \ref{lem:estimateHelmholtzDerivatives}.
    Using these estimates and Lemma \ref{lem:estimateHelmholtzDerivatives} gives
    \begin{align*}
      |\nabla_x^l \Gamma(x; \lambda) |
      &\leq |\nabla_x^l G(x; \lambda)| + \frac{1}{|\lambda|} |\nabla_x^{l + 2} G(x; \lambda) | + \frac{1}{|\lambda|} |\nabla_x^{l + 2} G(x; 0)| \\
    &\leq\frac{C_l \e^{-c \sqrt{|\lambda|} |x|}}{|x|^{d - 2 + l}} + \frac{1}{|\lambda|} \frac{C_{l + 2} \e^{-\sqrt{|\lambda|} |x|}}{|x|^2 |x|^{d - 2 + l}} + \frac{1}{|\lambda|} \frac{C}{|x|^2 |x|^{d - 2 + l}} \\
    %&\leq \frac{C}{(1 + |\lambda||x|^2)} \frac{1}{|x|^{d - 2 + l}} + \frac{C_{l + 2} |\lambda||x|^2}{(1 + |\lambda||x|^2)} \frac{1}{|\lambda| |x|^2} \frac{1}{|x|^{d - 2 + l}} + \frac{C}{1 + |\lambda||x|^2} \frac{1}{|x|^{d - 2 + l}}
    &\leq \frac{C}{1 + |\lambda| |x|^2}\; \frac{1}{|x|^{d - 2 + l}}
    \end{align*}
    Now let $|\lambda| |x|^2 \leq \frac{1}{2}$.
    Then by \ref{lem:estimateHelmholtzDerivatives} and \ref{lem:HelmholtzLaplaceDifference} we get
    \begin{align*}
      |\nabla_x^l \Gamma(x; \lambda) |
      &\leq |\nabla_x^l G(x; \lambda) | + \frac{1}{|\lambda|} |\nabla_x^{l + 2} ( G(x; \lambda) - G(x; 0))| \\
      &\leq \frac{C}{|x|^{d - 2 + l}} + \frac{1}{|\lambda|} \cdot C |\lambda| |x|^{4 - d - (l + 2)} \\
      &\leq \frac{C}{|x|^{d - 2 + l}} \frac{(1 + |\lambda| |x|^2)}{(1 + |\lambda| |x|^2)} \\
      &\leq \frac{C}{(1 + |\lambda| |x|^2) |x|^{d - 2 + l}}
    \end{align*}
    which gives the claim.
    %If $d = 2$ the steps are analogous considering the different structure of the estimate \eqref{eq:HelmholtzLaplaceDifferenceEstimate2d}.
  \end{proof}

 %Regarding the structure of the above proof, the derived estimate for the case $d = 2$ seems natural at first sight. 
 %One important difference compared to $d \geq 3$ is that it is not possible to derive an estimate of the form $|\nabla_x^l \Gamma(x; \lambda)| \leq C |x|^{-l}$.

  If $\lambda = 0$, the matrix of fundamental solutions to the Stokes problem in $\R^d$ with pole at the origin is given by $\Gamma(x; 0) = (\Gamma_{\alpha\beta}(x; 0))_{d \times d}$, where
  \begin{align}
    \label{eq:fundamentalSolutionStokes}
    \Gamma_{\alpha\beta}(x; 0) &= \frac{1}{2 \omega_d} \big\{ \frac{\delta_{\alpha\beta}}{(d - 2) |x|^{d - 2}} + \frac{x_{\alpha} x_\beta}{|x|^d} \big\}
    \intertext{if $d \geq 3$ and}
    \label{eq:fundamentalSolutionStokes2d}
    \Gamma_{\alpha\beta}(x; 0) &= \frac{1}{2 \omega_2} \big\{ - \delta_{\alpha\beta} \log(|x|) + \frac{x_\alpha x_\beta}{|x|^2} \big\}  \end{align}
  for $d = 2$.
  Note that the given fundamental solution for the case $d = 2$ differs from the one given by Mitrea and Wright \cite{mitreaWright} by having summands with alternating signs.
  The alternatig sign is necessary for $\Gamma_{\alpha\beta}$ to be divergence free.

  One important technique in the following chapter will be to reduce problems formulated for $\Gamma(x; \lambda)$ to problems formulated in $\Gamma(x; 0)$ and the difference $\Gamma(x; \lambda) - \Gamma(x; 0)$.
  Under this aspect it seems reasonable to study estimates of the difference of fundamental solutions.
  To this end it is helpful to rewrite parts of the fundamental solution.
  Using the fact that for $d \geq 5$ or $d = 3$
  we have
  \begin{align*}
    \frac{\partial^2}{\partial x_\alpha \partial x_\beta} \big( \frac{1}{|x|^{d - 4}} \big)
    = -(d - 4) \frac{\partial}{\partial x_\alpha} \frac{x_\beta}{|x|^{d - 2}}
    = -(d - 4) \frac{\delta_{\alpha\beta}}{|x|^{d - 2}} + \frac{(d - 4)(d - 2) x_\alpha x_\beta}{|x|^d}
  \end{align*}
  This allows us to express
  \begin{align*}
    \frac{x_\alpha x_\beta}{|x|^d} = \frac{\delta_{\alpha \beta}}{(d - 2)|x|^{d - 2}} + \frac{1}{(d - 4)(d - 2)} \frac{\partial^2}{\partial x_\alpha \partial x_\beta} \big( \frac{1}{|x|^{d - 4}} \big)
  \end{align*}
  which considering \eqref{eq:fundamentalSolutionStokes} gives 
  \begin{align}
    \label{eq:fundamentalSolutionStokes35}
    \Gamma_{\alpha\beta}(x; 0) &= G(x; 0)\delta_{\alpha\beta} + \frac{1}{2\omega_d(d - 4)(d - 2)} \frac{\partial^2}{\partial x_\alpha \partial x_\beta} \big( \frac{1}{|x|^{d - 4}} \big).
  \intertext{A similar trick works for $d = 4$.
  Since $\omega_4 = 2 \pi^2$, we have}
    \Gamma_{\alpha \beta}(x; 0)
    &= \frac{1}{2 \omega_4} \frac{1}{|x|^2} \delta_{\alpha\beta} - \frac{1}{8\pi^2} \big( \frac{\delta_{\alpha\beta}}{|x|^2} - \frac{2 x_\alpha x_\beta}{|x|^4} \big) \nonumber\\
    \label{eq:fundamentalSolutionStokes4}
    &= G(x; 0) \delta_{\alpha\beta} - \frac{1}{8\pi^2} \frac{\partial^2}{\partial x_\alpha \partial x_\beta}( \log(|x|))
  \intertext{In the case $d = 2$ this game shows that since
    $$
  \frac{1}{8\pi} \frac{\partial^2}{\partial_\alpha \partial_\beta} (|x|^2 \log(|x|))
  = \frac{\delta_{\alpha\beta}}{4\pi} \log(|x|) + \frac{1}{4\pi} \frac{x_\alpha x_\beta}{|x|^2} + \frac{\delta_{\alpha\beta}}{8\pi},
$$
we can write}
    \label{eq:fundamentalSolutionStokes2}
    \Gamma_{\alpha \beta}(x; 0) &= G(x; 0)\delta_{\alpha\beta} - \frac{\delta_{\alpha\beta}}{8\pi} - \frac{\partial^2}{\partial x_\alpha \partial x_\beta} (|x|^2 \log(|x|))
  \end{align}
  This ends the preparatory step and brings us to the next theorem.
  
  \begin{thm}
    \label{thm:differenceFundamentalSolutionStokes}
    Let $\lambda \in \Sigma_\theta$.
    Suppose that $|\lambda| |x|^2 \leq \frac{1}{2}$.
    Then
    \begin{align}
      |\nabla_x \{ \Gamma(x; \lambda) - \Gamma(x; 0) \} |
      \leq 
      \begin{cases}
        C |\lambda| |x|^{3 - d} &\quad\text{if } d \geq 7 \text{ or } d = 5, \\
        C |\lambda| |x|^{3 - d} | \log(|\lambda| |x|^2) | &\quad\text{if } d = 4 \text{ or } 6, \\
        C \sqrt{|\lambda|} |x|^{-1} &\quad\text{if } d = 3,\\
        C |\lambda| |x| ( |\log(|\lambda| |x|^2)| + 1 ) &\quad\text{if } d = 2,
      \end{cases}
    \end{align}
    where $C$ depends only on $d$ and $\theta$.
  \end{thm}

\begin{proof}
  We will split the proof in several parts.
  We start by considering the cases $d = 3$ and $d \geq 5$.
      Taking into account \eqref{eq:fundamentalSolutionStokes35} we have
      \begin{align*}
        &\Gamma_{\alpha\beta}(x; \lambda) - \Gamma_{\alpha\beta}(x; 0) \\
        &\qquad= \{ G(x; \lambda) - G(x; 0) \}\delta_{\alpha\beta}  \\
        &\qquad\qquad- \frac{1}{\lambda} \frac{\partial^2}{\partial x_\alpha \partial x_\beta} \Big\{ G(x; \lambda) - G(x; 0) + \frac{\lambda}{2 \omega_d (d - 4) (d - 2) |x|^{d - 4}} \Big\}
      \end{align*}
  As the first term can already be estimated via Lemma \ref{lem:HelmholtzLaplaceDifference}, we will only be concerned about the second one.
  If $d = 3$, a direct calculation will yield the desired result:
  We start by noting that $\omega_3 = 4 \pi$ gives
  \begin{align*}
    \frac{\e^{\ii k |x|}}{4\pi |x|} - \frac{1}{4 \pi |x|} - \frac{(\ii k)^2}{2 \omega_3 |x|^{-1}}
    &= \frac{1}{4 \pi |x|} \Big( \e^{\ii k |x|} - 1 - \frac{ (\ii k)^2 |x|^2}{2} \Big)\\
    &= \frac{1}{4 \pi |x|} \Big( \ii k |x| + \sum_{n = 3}^\infty \frac{(\ii k |x| )^n}{n!} \Big) \\
    &= \frac{1}{4\pi} \Big( \ii k + \sum_{n = 3}^\infty \frac{ (\ii k )^n |x|^{n - 1}}{n!} \Big).
  \end{align*}
  Taking the first derivative of this expression we get
  \begin{align*}
    \frac{\partial}{\partial x_\beta} \dots = 
    \frac{x_\beta}{4 \pi} \sum_{n - 3}^\infty \frac{(ik)^n (n - 1) }{n!} |x|^{n - 3}
  \end{align*}
  and differentiating with respect to $x_\alpha$ yields
  \begin{align*}
    \frac{\partial}{\partial x_\alpha} \dots 
    = \frac{\delta_{\alpha\beta}}{4 \pi} \sum_{n - 3}^\infty \frac{(ik)^n (n - 1) }{n!} |x|^{n - 3}
    + \frac{x_\beta x_\alpha}{4\pi} \sum_{n = 4}^\infty \frac{(\ii k)^n (n - 1) (n - 3)}{n!} |x|^{n - 5}.
  \end{align*}
  As we are interested in estimating the gradient of the difference of $\Gamma(x; \lambda)$ and $\Gamma(x; 0)$ we have to consider one additional derivative. This leaves us with
\begin{align*}
    \frac{\partial}{\partial x_\gamma} \dots
    &= \frac{\delta_{\alpha\beta} x_\gamma + \delta_{\beta\gamma} x_\alpha + \delta_{\alpha\gamma} x_\beta}{4\pi} \sum_{k = 4}^\infty \frac{(\ii k)^n (n - 1) (n - 3)}{n!}  |x|^{n - 5} \\
    &\qquad + \frac{x_\beta x_\alpha x_\gamma}{4\pi} \sum_{\substack{n = 4}}^\infty \frac{(\ii k)^n (n - 1) (n - 3) (n - 5)}{n!} |x|^{n - 7}.
\end{align*}
We can now prove the stated estimate
\begin{align*}
  | \frac{1}{\lambda} \frac{\partial^3}{\partial x_\gamma \partial x_\alpha \partial x_\beta} \dots  |
  &\leq \frac{1}{|k|^2\pi}  \sum_{k = 4}^\infty \frac{|k|^n (n - 1) (n - 3)(1 + (n - 5))}{n!} |x|^{n - 4} \\
  &\leq \frac{1}{|k|^2 |x|\pi} |k|^3 \sum_{k = 4}^\infty \frac{(n - 1)(n - 3)(1 + (n - 5))}{n!} |k|^{n - 3} |x|^{n - 3}\\
  &\leq C \frac{1}{|k| |x|}.
\end{align*}
This gives the claim for $d = 3$.
If $d \geq 5$, equation \eqref{eq:HelmholtzLaplaceDifference} gives
\begin{align*}
  &G(x; \lambda)- G(x; 0) + \frac{\lambda}{2 \omega_2 (d - 4) (d - 2) |x|^{d - 4}} \\
  &\qquad= \frac{\ii}{4 (2\pi)^{\frac{d}{2} - 1}} \frac{1}{|x|^{d - 2}} \big\{ z^{\frac{d}{2} - 1} H_{\frac{d}{2} - 1}^{(1)}(z) - a_d - b_d z^2\big\},
\end{align*}
where $z = k|x|$, $a_d$ was calculated in \eqref{eq:Defnad} and $b_d$ is given by
\begin{align*}
 b_d 
 &= \frac{2\ii (2\pi)^{\frac{d}{2} - 1}}{\omega_d (d - 4) (d - 2)}  
 = -\frac{2\ii (2\pi)^{\frac{d}{2} - 1} \Gamma(\frac{d}{2})}{2\pi^{\frac{d}{2}} (d - 2)(d - 4)}
 = \frac{2^{\frac{d}{2} - 1} }{\pi \ii (d - 4)} \frac{\Gamma(\frac{d}{2})}{(d - 2)} \\
 &= \frac{2^{\frac{d}{2} - 1} }{2\pi \ii } \frac{\Gamma(\frac{d}{2} - 1)}{(d - 4)} 
 = \frac{2^{\frac{d}{2} - 1}}{4 \pi \ii } \frac{\Gamma(\frac{d}{2} - 1)}{(\frac{d}{2} - 1 - 1)}
 = \frac{2^{\frac{d}{2} - 1} \Gamma(\frac{d}{2} - 2)}{4 \pi \ii }.
\end{align*}
If $d \geq 7$ this shows that $b_d$ is the second coefficient of the asymptotic expansion \eqref{eq:asymptoticd7} and thus we can estimate
\begin{align}
  \label{eq:estimateDerivativesd7}
  \big| \frac{\mathrm{d}^l}{\d z^l} \Big\{ z^{\frac{d}{2} - 1} H_{\frac{d}{2} - 1}^{(1)}(z) - a_d - b_d z^2 \Big\} \big| \leq C |z|^{4 - l}
\end{align}
for $0 \leq l \leq 3$ and $z \in \C \setminus (-\infty, 0]$.
For better readability we set
\begin{align*}
  g(z) = z^{\frac{d}{2} - 1} H_{\frac{d}{2} - 1}^{(1)}(z) - a_d - b_d z^2
\end{align*}
and consider the function $f(x) = g(k |x|)$ on $\R^d \setminus \{0\}$.
The derivatives of $f$ read
\begin{align*}
  \frac{\partial}{\partial x_\beta} f(x)
  &= (\frac{\d{}}{\d z} g)(k|x|) \frac{k x_\beta}{|x|} \\
  \frac{\partial^2}{\partial x_\alpha \partial x_\beta} f(x)
  &= (\frac{\mathrm{d}^2}{\d z^2} g)(k|x|) \frac{k^2 x_\alpha x_\beta}{|x|^2} + (\frac{\d{}}{\d z}) g(k |x|) k \Big\{ \frac{\delta_{\alpha\beta}}{|x|} - \frac{x_\beta x_\alpha}{|x|^3} \Big\} \\
  \frac{\partial^3}{\partial x_\gamma \partial x_\alpha \partial x_\beta} f(x)
  &= (\frac{\mathrm{d}^3}{\d z^3}g)(k |x|) \frac{k^3 x_\alpha x_\beta x_\gamma}{|x|^3} \\
  &\quad + (\frac{\mathrm{d}^2}{\d z^2}g)(k|x|) 
  k^2 \Big\{ \frac{x_\alpha \delta_{\beta\gamma} + x_\beta \delta_{\alpha\gamma} + x_\gamma \delta_{\alpha\beta}}{|x|^2} 
  - \frac{3 x_\alpha x_\beta x_\gamma}{|x|^4} \Big\} \\
  &\quad + (\frac{\d{}}{\d z}g)(k|x|) k \Big\{-\frac{\delta_{\alpha\beta} x_\gamma}{|x|^3} - \frac{x_\alpha \delta_{\beta\gamma} + x_\beta \delta_{\alpha\gamma}}{|x|^3} + \frac{3 x_\alpha x_\beta x_\gamma}{|x|^5} \Big\}.
\end{align*}
If we now look for estimates on the absolute value of the derivatives, we see that by \eqref{eq:estimateDerivativesd7}
\begin{align*}
% |\nabla_x f(x) | 
% &\leq C |k|^{4 - 1 + 1} |x|^{4 - 1} = C |k|^{4} |x|^{3}\\
% |\nabla_x^2 f(x) | 
% &\leq C \{ |k|^{4 - 2 + 2} |x|^{4 - 2} + |k|^{4 - 1 + 1} |x|^{4 -1 -1} \}  = C |k|^4 |x|^2 \\
% |\nabla^3_x f(x) | 
% &\leq C \{ |k|^{4 - 3 + 3} |x|^{4 - 3} + |k|^{4 - 2 + 2} |x|^{4 - 2 -1} + |k|^{4 - 1 + 1} |x|^{4 - 1 - 2}
  |\nabla_x^l f(x)| \leq C |k|^4 |x|^{4 - l}, \quad 1 \leq l \leq 3,
\end{align*}
where $C$ only depends on $l$.
We can now finally uncover the desired estimate via
\begin{align*}
  &\Big|\frac{1}{\lambda} \nabla_x^3 \Big\{ \frac{\ii}{4 (2\pi)^{\frac{d}{2} - 1}} \frac{1}{|x|^{d - 2}} \big\{ z^{\frac{d}{2} - 1} H_{\frac{d}{2} - 1}^{(1)}(z) - a_d - b_d z^2\big\} \Big\} \Big| \\
  &\quad\leq C \frac{1}{|k|^2} \sum_{l = 0}^3 | \nabla_x^{3 - l}(\frac{1}{|x|^{d - 2}})| | \nabla_x^l f(x) | 
  \leq C \sum_{l = 0}^3 |x|^{-d + 2 - 3 + l} |k|^2 |x|^{4 - l} = C |\lambda| |x|^{3 - d},
\end{align*}
where $C$ is a constant only depending on $d$.
If $d = 6$, this shows that the asymptotic expansion \eqref{eq:asymptoticd6} gives us similar to \eqref{eq:estimateDerivativesd7}
\begin{align}
  \label{eq:estimateDerivativesd6}
  \big| \frac{\mathrm{d}^l}{\d z^l} \Big\{ z^{\frac{d}{2} - 1} H_{\frac{d}{2} - 1}^{(1)}(z) - a_d - b_d z^2 \Big\} \big| \leq C |z|^{4 - l} | \log(z) |,
\end{align}
for $0 \leq l \leq 3$ and $z \in \C \setminus (-\infty, 0]$.
Using as before the expressions for derivatives of $f$, we can estimate
\begin{align*}
  |\nabla_x^l f(x) | \leq C |k|^4 |x|^{4 - l} |\log(|\lambda| |x|^2) |, 
\end{align*}
which by a calculation analogous to the case $d \geq 7$ yields
\begin{align*}
  | \nabla_x \{ \Gamma(x; \lambda) - \Gamma(x; 0) \}| \leq C |\lambda| |x|^{3 - d} | \log(|\lambda| |x|^2) |.
\end{align*}
For $d = 5$ write
\begin{align*}
  &\frac{\partial^2}{\partial x_\alpha \partial x_\beta} \Big\{ G(x; \lambda) - G(x; 0) + \frac{\lambda}{2 \omega_d (d - 4) (d - 2) |x|{d - 4}} \Big\} \\
  &\qquad= \frac{\partial^2}{\partial x_\alpha x_\beta} \Big\{ \frac{\ii}{4(2\pi)^{\frac{3}{2}}} \cdot \frac{1}{|x|^3} \big[ z^{\frac{3}{2}} H_{\frac{3}{2}}^{(1)}(z) - a_5 - b_5 z^2 - w z^3 \big] \Big\},
\end{align*}
where $w\in\C$ can be an any constant if we set $z = k|x|$.
Now, for the appropriate choice of $w \in \C$ the asymptotic expansion \eqref{eq:asymptoticd5} gives the same estimate as \eqref{eq:estimateDerivativesd7} which proves the claim for $d = 5$.

In the case $d = 4$ we use \eqref{eq:fundamentalSolutionStokes4} to reformulate
\begin{align*}
  &\Gamma_{\alpha\beta}(x; \lambda) - \Gamma_{\alpha\beta}(x; 0) \\
  &\quad= \{ G(x; \lambda) - G(x; 0) \}\delta_{\alpha\beta}  
  - \frac{1}{\lambda} \frac{\partial^2}{\partial x_\alpha \partial x_\beta} \Big\{ G(x; \lambda) - G(x; 0) - \frac{\lambda \log(|x|)}{8 \pi^2} \Big\} \\
  &\quad= \{ G(x; \lambda) - G(x; 0) \} \delta_{\alpha\beta} \\
  &\qquad \frac{\ii}{\lambda} \frac{\partial^2}{\partial x_\alpha \partial x_\beta} \Big\{ \frac{1}{8 \pi |x|^2} [ z H_1^{(1)}(z) - a_4 - w z^2 - b_4 z^2 \log(z)] \Big\},
\end{align*}
where $z = k |x|$, $b_4 = \frac{\ii}{\pi}$ and $w \in \C$ is an arbitrary constant. Using the asymptotic expansion \eqref{eq:asymptoticd4} and the appropriate constant $w \in \C$ we get the estimate
\begin{align*}
  \big| \frac{\mathrm{d}^l}{\d z^l} \Big\{ z H_1^{(1)}(z) - a_4 - w z^2 - b_4 z^2 \log(z) \Big\} \big| \leq C |z|^{4 - l} |\log(z)|.
\end{align*}

For $d = 2$ the estimate follows from a direct calculation.
\end{proof}

We can now use the assumption $|\lambda||x|^2 \leq \frac{1}{2}$  to unify the structure of the estimates for $d \geq 2$.
\begin{cor}
  \label{cor:differenceFundamentalSolutionStokes}
  Let $\lambda \in \Sigma_\theta$. Suppose that $|\lambda||x|^2 \leq \frac{1}{2}$.
  Then for all $d \geq 2$
  \begin{align*}
    |\nabla_x\{ \Gamma(x; \lambda) - \Gamma(x; 0)\} | \leq C \sqrt{|\lambda|} |x|^{2 - d},
  \end{align*}
  where $C$ depends only on $d$ and $\theta$.
\end{cor}

\begin{proof}
  We just extend the estimates given in Theorem \ref{thm:differenceFundamentalSolutionStokes}.
  Let $d \geq 7$ or $d = 5$. Since $|\lambda|^{\frac{1}{2}} \leq C{|x|}$ we have
  \begin{align*}
    C |\lambda| |x|^{3 - d} \leq C |\lambda|^{\frac{1}{2}} |x|^{2 - d}.
  \end{align*}
  For $d = 4,6$ we have
  \begin{align*}
    C |\lambda| |x|^{3 - d} | \log(|\lambda||x|^2)|
    = C |\lambda|^{\frac{1}{2}} |x|^{2 - d} \cdot |\lambda|^{\frac{1}{2}} |x| | \log(|\lambda| |x|^2) |
    \leq C |\lambda|^{\frac{1}{2}} |x|^{2 - d},
  \end{align*}
  since $|\lambda|^{\frac{1}{2}} |x| |\log(|\lambda| |x|^2)|$ is bounded for $|\lambda| |x|^2 \leq \frac{1}{2}$.
  For $d = 2$ the same argument applies to the expression $|\lambda|^{\frac{1}{2}} |x| (|\log(|\lambda| |x|^2)| + 1)$.
\end{proof}


\chapter{Single and Double Layer Potentials}
\label{chap:3}

In this chapter, we will deal with \emph{single} and \emph{double layer potentials}.
Both will serve as ``representation formulas'' for solutions to the Stokes resolvent problem.
We will study their properties as they will serve as the crucial ingredient to solving the Neumann and Dirichlet boundary problems associated to the Stokes resolvent problem.
In this chapter we will always assume that $\Omega$ is a bounded Lipschitz domain in $\R^d$ with $d \geq 2$ and $1 < p < \infty$.
We will also tacitly use the summation convention.

Let $\lambda \in \Sigma_\theta$. 
For $f \in \Ell^2(\partial\Omega; \C^d)$, the single layer potential $u = \slp_\lambda(f)$ is defined by
\begin{align}
  \label{eq:defSingleLayer}
  u_j(x) 
  = \int_{\partial\Omega} \Gamma_{jk}(x - y; \lambda) f_k(y) \d \sigma(y),
\end{align}
where $\Gamma_{jk}$ is the fundamental solution to the Stokes reolvent problem given by \eqref{eq:fundamentalMatrixStokes}.
For the pressure, respectively, we define the single layer potential $\phi = \slp_\Phi(f)$ by
\begin{align}
  \label{eq:defSingleLayerPressure}
  \phi(x) = \int_{\partial\Omega} \Phi_k(x - y) f_k(y) \d \sigma(y),
\end{align}
where $\Phi_k$ is given by \eqref{eq:fundamentalVectorPressure}.
As we have already shown, $(u,\phi)$ defines a solution to the Stokes resolvent problem \eqref{eq:stokesResolventProblem}.

We define two further integral operators
\begin{align}
  T_\lambda^*(f)(P) &= \sup_{t > 0} \big| \int_{\substack{y \in \partial\Omega \\ |y - P| > t}} \nabla_x \Gamma(P - y; \lambda) f(y) \d \sigma(y) \big| \label{eq:supTOperator}\\
  T_\lambda(f)(P) &= \pv \int_{\partial\Omega} \nabla_x \Gamma(P - y; \lambda) f(y) \d \sigma(y) \label{eq:pvTOperator}
\end{align}
for $P \in \partial\Omega$ which will be used to prove boundedness of maximal operators related to $u$.

\begin{lem}
  \label{lem:lpBoundednessT}
  Let $1 < p < \infty$ and $T_\lambda(f), T_\lambda^*(f)$ be defined by \eqref{eq:supTOperator} and \eqref{eq:pvTOperator}.
  Then $T_\lambda(f)(P)$ exists for almost everywhere $P \in \partial\Omega$ and
  \begin{align}
    \label{eq:lpBoundednessT}
    \|T_\lambda(f) \|_{\Ell^p(\partial\Omega)} 
    \leq \|T_\lambda^*(f) \|_{\Ell^p(\partial\Omega)}
    \leq C_p \, \|f\|_{\Ell^p(\partial\Omega)},
  \end{align}
  where $C_p$ depends only on $d$, $\theta$, $p$, and the Lipschitz character of $\Omega$.
\end{lem}

\begin{proof}
  If $\lambda = 0$, the Lemma is known \cite{fabesKenigVerchota} as a consequence of the seminal result of Coifman et al. \cite{coifmanEtAl}.
  One idea of the proof in the case $\lambda \in \Sigma_\theta$ will thus be to nourish from this result and to consider the difference $\Gamma(x - y; \lambda) - \Gamma(x - y; 0)$ as a well-disposed integral kernel.

  We start with the second inequality of \ref{eq:lpBoundednessT}.
  To this end, let $t > 0$ and additionally assume that $t^2 |\lambda| \geq \frac{1}{2}$. 
  Theorem \ref{thm:fundamentalMatrixEstimate} gives
  \begin{align*}
    \Big| \int_{|y - P| > t} \nabla \Gamma(P - y; \lambda) f(y) \d\sigma(y) \Big|
    &\leq C \int_{|P - y| > t} \frac{|f(y)|}{|\lambda| |P - y|^{d + 1}} \d \sigma(y) 
  \end{align*}
  Choose now $N \in \N$ such that $2^N t \leq \diam(\Omega) < 2^{N + 1} t$.
  We now exhaust the domain of integration by suitable annuli and calculate
  \begin{align*}
    &\sum_{k = 0}^N \int_{\BB(P, 2^{k + 1}t) \cap \partial\Omega} \frac{1}{|\lambda| 2^{k(d + 1)} t^{d + 1}} |f(y)| \d \sigma(y) \\
    &\quad\leq \frac{1}{|\lambda|t^2}\frac{1}{2^{1 -d}}\sum_{k = 0}^N \frac{1}{2^{2k}} \frac{1}{(2^{k+1} t)^{d - 1}} \int_{\BB(P, 2^{k + 1}t \cap \partial \Omega)}  |f(y)| \d\sigma(y) \\
    &\quad\leq C \sum_{k = 0}^N \frac{1}{2^{2k}} M_{\partial\Omega}(f)(P) \\
    &\quad\leq C  M_{\partial\Omega}(f)(P)
    %&\quad\leq C |\lambda|^{\frac{1}{2}} \frac{2^{N + 1}t - t}{1}  M_{\partial\Omega}(f)(P) \\
    %&\quad\leq C |\lambda|^{\frac{1}{2}} t
  \end{align*}
  where for the second inequality we used Lemma \ref{lem:compareBoundaryWithBall} to estimate
  \begin{align*}
    \frac{1}{(2^{k+1} t)^{d-1}} \leq C (\sigma(\BB(P, 2^{k + 1}d) \cap \partial\Omega))^{-1}.
  \end{align*}
  which gives the claimed estimate with a constant $C$ that depends on $d$, $\theta$ and the Lipschitz character of $\Omega$.
  Now let $t^2 |\lambda| < \frac{1}{2}$.
  We then split the integral as follows
  \begin{align*}
    &\Big| \int_{|y - P| > t} \nabla \Gamma(P - y; \lambda) f(y) \d\sigma(y) \Big| \\
    &\quad\leq \Big| \int_{|y - P| \geq (2|\lambda|)^{-1/2}} \nabla\Gamma(P - y; \lambda)f(y) \d\sigma\Big|
     + \Big| \int_{t < |y - P| < (2|\lambda|)^{-1/2}} \nabla\Gamma(P - y; \lambda)f(y) \d\sigma\Big|.
  \end{align*}
  The first summand can be estimated like in the step before, if we substitute $t$ by $(2|\lambda|)^{-1/2}$
  For the second term we use the principle of the nutrient zero and estimate
  \begin{align*}
     &\Big| \int_{t < |y - P| < (2|\lambda|)^{-1/2}} \nabla\Gamma(P - y; \lambda)f(y) \d\sigma\Big| \\
     &\qquad\leq \int_{t < |y - P| < (2|\lambda|)^{-1/2}} |\nabla\Gamma(P - y; \lambda) - \nabla\Gamma(P - y; 0) |f(y)| \d\sigma \\
     &\qquad\quad + \Big| \int_{t < |y - P| < (2|\lambda|)^{-1/2}} \nabla\Gamma(P - y; 0) f(y)  \d\sigma \Big|.
  \end{align*}
  We don't need to worry about the second summand here since the corresponding estimate is already coverd by the case of $\lambda = 0$ and therefore
  \begin{align*}
     &\Big| \int_{t < |y - P| < (2|\lambda|)^{-1/2}} \nabla\Gamma(P - y; 0) f(y)  \d\sigma \Big| \\
     &\quad\leq \Big| \int_{|y - P| > t} \nabla\Gamma(P - y; 0) f(y)  \d\sigma \Big| +  \Big| \int_{|y - P| > (2|\lambda|)^{-1/2}} \nabla\Gamma(P - y; 0) f(y)  \d\sigma \Big| \\
     &\quad\leq 2 T_0^*(f)(P).
   \end{align*}
   For the first summand we make use of Theorem \ref{thm:differenceFundamentalSolutionStokes} and more precisely of Corollary \ref{cor:differenceFundamentalSolutionStokes} which unifies all estimates.
   We then calculate
   \begin{align*}
     &\int_{t < |y - P| < (2|\lambda|)^{-1/2}} |\nabla\Gamma(P - y; \lambda) - \nabla\Gamma(P - y; 0) |f(y)| \d\sigma \\
     &\qquad\leq \int_{t < |y - P| < (2|\lambda|)^{-1/2}} |\lambda|^{\frac{1}{2}} |y - P|^{2 - d} |f(y)| \d\sigma,
     \intertext{and as before we choose adequate $N$ such that $2^{N + 1} t > (2|\lambda|)^{-1/2} \geq 2^N t$ which leads to}
     &\qquad\leq |\lambda|^{\frac{1}{2}} \sum_{k = 0}^N \int_{2^kt < |y - P| < 2^{k + 1} t} |y - P|^{2 - d} |f(y)| \d\sigma \\
     &\qquad\leq  |\lambda|^{\frac{1}{2}} t^{2 - d} \sum_{k = 0}^N 2^{k(2 - d)} \int_{\BB(P, 2^{k + 1}t)} |f(y)| \d\sigma\\
     &\qquad\leq 2^d  |\lambda|^{\frac{1}{2}} t \sum_{k = 0}^N 2^{k - 1}\; 2^{(k + 1)(1 - d)} t^{1 - d} \int_{\BB(P, 2^{k + 1}t)} |f(y)| \d\sigma \\
     &\qquad \leq C  |\lambda|^{\frac{1}{2}} t\frac{2^{N} - 1}{1} M_{\partial\Omega}(f)(P) \\
     &\qquad \leq C |\lambda|^{\frac{1}{2}} (2|\lambda|)^{-\frac{1}{2}} M_{\partial\Omega}(f)(P).
   \end{align*}
   Taking now the supremum over all $t > 0$ we see that
   \begin{align*}
     T_\lambda^*(f)(P)
     &\leq C (M_{\partial\Omega}(f)(P) + T_0^*(f)(P)),
   \end{align*}
   for all $P \in \partial\Omega$. Once again using the result for $\lambda = 0$ and the $L^p$-boundedness of the Hardy-Littlewood maximal operator we see that
   \begin{align*}
     \|T_\lambda^*(f) \|_{\Ell^p(\partial\Omega)} \leq C\, \|f\|_{\Ell^p(\partial\Omega)}
   \end{align*}

   To conclude the first inequality in \eqref{eq:lpBoundednessT}, we want to use a standard result from harmonic analysis \cite[2.1.14]{grafakos}.
   First we will show that the integral operator
   \begin{align*}
     T_\lambda(f)(P) = \lim_{t \to 0} \int_{\substack{y \in \partial\Omega \\ |y - P| > t}} \nabla_x \Gamma(P - y; \lambda) f(y) \d\sigma(y)
   \end{align*}
   exists for almost every $P \in \partial\Omega$ and all $f \in \CC(\partial\Omega; \C^d)$.
   In a first step, we can split this operator formally in
   \begin{align*}
     T_\lambda(f)(P) = T_0(f)(P) + \lim_{t \to 0} \int_{\substack{y \in \partial\Omega \\ |y - P| > t}} \nabla_x \{ \Gamma(P - y; \lambda) - \Gamma(P - y; 0)\}  f(y) \d\sigma(y)
   \end{align*}
   The right expression is well defined for $f \in \CC_0^\infty$, once we prove integrability of 
   \begin{align*}
     |\nabla \{\Gamma(P - y; \lambda) - \Gamma(P - y; 0)\} |
   \end{align*}
   on $\partial\Omega$.
   To this end we first note that it suffices to consider the integral
   \begin{align*}
     \int_{|P - y| \leq \varepsilon} | \nabla\{ \Gamma(P - y; \lambda) - \Gamma(P - y; 0) \} | \d\sigma(y), 
   \end{align*}
   for $\varepsilon \leq \min(2|\lambda|^{-1/2}, r_0/4)$ as the integrand is smooth away from $0$ and the domain of integration is bounded.
   Now Corollary \ref{cor:differenceFundamentalSolutionStokes} and Tolksdorf 4.3.2 give that this can be estimated by
   \begin{align*}
     \int_{|P - y| \leq \varepsilon} |\lambda|^{1/2} |P - y|^{2 - d} \d\sigma(y) \leq C |\lambda|^{1/2} \varepsilon \leq C.
   \end{align*}
   Based on the preceding calculation we conclude that for all $f \in \CC(\partial\Omega,\C^d)$ the operator $T_\lambda(f)(P)$ exists whenever $T_0(f)(P)$ exists.
   $T_0(f)(P)$ exists for almost everywhere $P \in \partial\Omega$ because of Fabes, Kenig and Verchota \cite{fabesKenigVerchota}.
   As furthermore $T_\lambda^*(f)(P)$ is bounded on $\Ell^p(\partial\Omega)$  we may now apply Theorem 2.1.14 from Grafakos \cite{grafakos} to conclude that $T_\lambda(f)(P)$ exists now for all $f \in \Ell^p(\partial\Omega; \C^d)$ and almost everywhere $P \in \partial\Omega$.
   The desired $\Ell^p$ estimate on $T_\lambda(f)$ now follows from the observation that 
   \begin{align*}
     | T_\lambda(f)(P) | \leq T_\lambda^*(f)(P)
   \end{align*}
   for almost everywhere $P \in \partial\Omega$.
\end{proof}

%For a function $u$ in $\Omega$, we define the nontangential maximal function $(u)^*$ by
%\begin{align}
%  \label{eq:defnNontangMaxFunction}
%  (u)^*(P) = \sup\{ |u(x)| \colon x \in \Omega \text{ and } |x - P| < C \operatorname{dist}(x, \partial\Omega)\}
%\end{align}
%for $P \in \partial\Omega$, where $C > 2$ is a fixed and sufficiently large constant depending on $d$ and the Lipschitz character of $\Omega$.
%Note that in Shen cones we have that for $P, y \in \partial\Omega$ and $x \in \ShenCone(P)$
%\begin{align}
%  \label{eq:shenConeEstimate}
%  |P - y| 
%  &\leq |P - x| + |x - y| 
%  \leq C \operatorname{dist}(x \partial\Omega) + |x - y|  \nonumber\\
%  &\leq (C + 1) |x - y|
%\end{align}
%where $C$ is the constant from \eqref{eq:defnNontangMaxFunction}

We can now prove the boundedness of certain nontangential maximal operators.
\begin{lem}
  \label{lem:nontangentialMaximalFunctions}
  Let $1 < p < \infty$ and $(u,\phi)$ be given by \eqref{eq:defSingleLayer} and \eqref{eq:defSingleLayerPressure}.
  Then 
  \begin{align}
    \| (\nabla u)^* \|_{\Ell^p(\partial\Omega)}  +  \| (\phi)^* \|_{\Ell^p(\partial\Omega)} \leq C_p \|f\|_{\Ell^p(\partial\Omega)},
  \end{align}
  where $C_p$ depends only on $d$, $\theta$, $p$ and the Lipschitz character of $\Omega$.
  Let furthermore $d \geq 3$. Then
  \begin{align*}
     \|(u)^*\|_{\Ell^p(\partial\Omega)} \leq C_p \|f\|_{\Ell^p(\partial\Omega)},
  \end{align*}
  where $C_p$ depends only on $d$, $\theta$, $p$ and the Lipschitz character of $\Omega$.
\end{lem}

\begin{proof}
  A proof of the estimate $\|(\phi)^*\|_{\Ell^p(\partial\Omega)} \leq C_p \|f\|_{\Ell^p(\partial\Omega)}$ can be found in Verchota \cite{verchota}.
  The proof for $\|(\nabla u)^*\|_{\Ell^p(\partial\Omega)}$ works in the same way. 
  We will provide a proof for the sake of completeness.
  To immitate the proof of Verchota, we will work with the corresponding type of cones.
  Therefore the results for $\nabla u$ and $\phi$ will at first only be established for the type of maximal operators defined by Verchota.
  The transferability to Shen's maximal operators is given by Tolksdorf \cite{tolksdorfDiss} as the solution $(u,\phi)$ has a representation as a single layer potential.

  Let $x \in \verCone(P)$ and set $t = |x - P|$.
  Then,
  \begin{align*}
    |(\nabla u)(x) |
    &= \big| \int_{\partial\Omega} \nabla \Gamma_{jk} (x - y; \lambda) f_k \d \sigma(y)\big| \\
    &\leq \big| \int_{|y - P| > t} \nabla\Gamma_{jk}(x - y; \lambda) f_k \d \sigma(y) \big| + \big| \int_{|y - P| \leq t} \nabla \Gamma_{jk}(x - y; \lambda) f_k \d \sigma(y)\big| \\
    &= I_1 + I_2.
  \end{align*}
  We will now estimate $I_1$ and $I_2$ seperately.
  Note that in Verchota cones we have that for all $Q \in \partial\Omega$ we have $|x - Q| > C |x - P|$, where $C$ is a constant only depending on $d$ and the Lipschitz character of $\Omega$.
  By Theorem \ref{thm:fundamentalMatrixEstimate} we know that
  \begin{align*}
    I_2 
    &\leq C \int_{|y - P| \leq t} \frac{1}{|x - y|^{d - 1}} |f(y)| \d \sigma(y) \\
    &\leq \frac{C}{t^{n - 1}} \int_{|y - P| \leq t} |f(y)| \d \sigma(y)
    \leq C M_{\partial\Omega} (f)(P).
  \end{align*}
%  For $I_1$, we calculate
%  \begin{align*}
%    &\Big| \int_{| y - P | > t} \nabla \Gamma_{j k} (x - y; \lambda) f_k(y) - \nabla \Gamma_{jk}(P - y; \lambda) f_k(y) + \nabla \Gamma_{jk}(P - y; \lambda) f_k(y) \d \sigma(y) \Big| \\
%    &\quad\leq \Big| \int_{|y - P| > t} \nabla\Gamma_{jk} (x - y; \lambda) f_k(y) - \nabla\Gamma_{jk}(P - y; \lambda) f_k(y) \d \sigma(y) \Big| \\
%    &\qquad+ \Big| \int_{|y - P| > t} \nabla\Gamma_{jk}(P - y; \lambda) f_k(y) \d\sigma(y) \Big|.
%  \end{align*}
   For $I_1$, we calculate
   \begin{align*}
     &\Big| \int_{| y - P | > t} \nabla \Gamma_{j k} (x - y; \lambda) f_k(y) - \nabla \Gamma_{jk}(x - y; 0) f_k(y) + \nabla \Gamma_{jk}(x - y; 0) f_k(y) \d \sigma(y) \Big| \\
     &\quad\leq \Big| \int_{|y - P| > t} \nabla\Gamma_{jk}(x - y; \lambda) f_k(y) - \nabla\Gamma_{jk}(P - y; \lambda) f_k(y) \d \sigma(y) \Big| \\
     &\qquad+ \Big| \int_{|y - P| > t} \nabla\Gamma_{jk}(P - y; \lambda) f_k(y) \d\sigma(y) \Big|.
   \end{align*}
  The second summand can directly be estimated by $T_\lambda^*(f)(P)$.
  For the second one we apply the mean value theorem and derive using once again Theorem \ref{thm:fundamentalMatrixEstimate}
  \begin{align*}
    &\int_{|y - P| > t} \big| \nabla\Gamma_{jk}(x - y; \lambda) - \nabla\Gamma_{jk}(P - y; \lambda) \big| |f(y)| \d\sigma(y) \\ 
    &\qquad\leq \int_{|y - P| > t} |\nabla^2 \Gamma_{jk}(s - y; \lambda)| |x - P| |f(y) | \d\sigma(y) \\
    &\qquad\leq C t \int_{|y - P| > t} \frac{1}{|s - y|^{d}} |f(y)| \d\sigma(y) \\
    &\qquad\leq C t \int_{|y - P| > t} \frac{1}{|y - P|^{d}} |f(y)| \d\sigma(y) \\
    &\qquad\leq C \int_{\partial\Omega} \frac{t}{(t + |y - P|)^d} |f(y)| \d\sigma(y).
  \end{align*}
  where $s$ is an element on the line connecting $x$ and $P$ and we used the property of Verchota-cones that $|s - y| \geq C |y - P|$.
  Note that Verchota cones are convex.
  As in Verchota \cite{verchotaDiss} the integral may now be bounded by the Hardy-Littlewood maximal operator due to an application of a suitable Lemma from Dahlberg \cite{dahlberg} as the kernel $\frac{t}{(t + |y - P|)^d}$ is uniformly integrable on $\partial\Omega$.
  Summing up we have shown that
  \begin{align*}
    |(\nabla u)(x)| \leq C \big\{ M_{\partial\Omega} f(P) + T_\lambda^*(f)(P) \}
  \end{align*}
  where $C$ only depends on $d$, $\theta$ and the Lipschitz character of $\Omega$.
  We thus may take the supremum over all $x \in \verCone$ and conclude the desired estimate by the well known mapping properties of the Hardy-Littlewood maximal operator and the respective results from Lemma \ref{lem:lqBoundednessT}.
%  By exhausting the domain of integration using annuli, we can estimate this integral by $M_{\partial\Omega}f(P)$:
%  Choose $N$ such that $2^N t \leq \diam(\Omega) < 2^{N + 1}t$.
%  Then
%  \begin{align*}
%    &\int_{|y - P| > t} \frac{1}{|y - P|^{d - 1}} |f(y)| \d\sigma(y) \\
%    &\quad= \sum_{k = 0}^N \int_{2^{k + 1}t > |y - P| \geq 2^k t} \frac{1}{|y - P|^{d - 1}} |f(y)| \d\sigma(y) \\
%    &\quad\leq \sum_{k = 0}^N \frac{1}{2^{k(d - 1)} t^{d - 1}} \int_{|y - P| < 2^{k + 1} t} |f(y)| \d\sigma(y) \\
%    &\quad \leq  C \sum_{k = 0}^N 2^{k - d + 1)} M_{\partial\Omega}(f)(P)
%    &\quad\leq
%  \end{align*}
%  Taking the supremum over all $x \in \Omega$ the claim follows.

  We will now work on the proof of the estimate for $(u)^*$ for $d \geq 3$.
  In order to derive $\Ell^p$ estimates on this maximal operator we will work directly with the Definition of the single layer potential \eqref{eq:defSingleLayer}.
  For $P \in \partial\Omega$, estimate \eqref{eq:fundamentalMatrixEstimate} together with the estimate for Shen cones \eqref{eq:shenConeEstimate} gives that for all $x \in \ShenCone(P)$
  \begin{align*}
    |u^*(x)| 
    &\leq  C \int_{\partial\Omega} \frac{1}{|x - y|^{d - 2}} |f(y)| \d\sigma(y) 
    \leq  C \int_{\partial\Omega} \frac{1}{|P - y|^{d - 2}} |f(y)| \d\sigma(y),
  \end{align*}
  where $C$ only depends on $d$, $\theta$ and the Lipschitz character of $\Omega$.
  Passing to the maximal operator yields the inequality
  \begin{align*}
    u^*(P) \leq  C \int_{\partial\Omega} \frac{1}{|P - y|^{d - 2}} |f(y)| \d\sigma(y),
  \end{align*}
  We are now left with the task to estimate the integral 
  \begin{align*}
    \int_{\partial\Omega} \frac{1}{|P - y|^{d - 2}} \d\sigma(y)
  \end{align*}
  uniformly for all $P \in \partial\Omega$, as the rest can be handled using the Young inequality.
  Let $r_0$ be the Radius from the definition of Lipschitz cylinders.
  Then
  \begin{align*}
    \int_{\partial\Omega} \frac{1}{|P - y|^{d - 2}} \d\sigma(y)
    &\leq \int_{\partial\Omega \cap \BB(P; r_0/4)} \frac{1}{|P - y|^{d - 2}} \d\sigma(y) + \int_{\partial\Omega \setminus \BB(P; r_0/4)} \frac{1}{|P - y|^{d - 2}} \d\sigma(y). \\
    & \leq C r_0/4 +  \sigma(\partial\Omega) r_0^{2 - d} 4^{d - 2}.
  \end{align*}
  where $C$ only depends on $d$ and the Lipschitz character of $\Omega$.
\end{proof}

The next Lemma deals with \emph{trace formulas} for $\nabla u$ and $\phi$. We can now finally talk about boundary values as the existence of nontangential limits guarantees that there exists something on $\partial\Omega$ that is related to the function inside $\Omega$.

\begin{lem}
  \label{lem:traceFormulas}
  Let $(u,\phi)$ be given by \eqref{eq:defSingleLayer} and \eqref{eq:defSingleLayerPressure} with $f \in \Ell^p(\partial\Omega; \C^d)$ and $1 < p < \infty$.
  Then
  \begin{align}
    \big( \frac{\partial u_i}{\partial x_j} \big)_{\pm}(x) 
    &= \pm \frac{1}{2} \{ n_j(x) f_i(x) - n_i(x) n_j(x) n_k(x) f_k(x) \} \nonumber\\
    &\quad+ \pv \int_{\partial\Omega} \frac{\partial}{\partial x_j} \{ \Gamma_{ik} (x - y; \lambda) \} f_k(y) \d\sigma(y), \label{eq:traceFormula} \\
    \phi_\pm(x) &= \mp \frac{1}{2} n_k(x) f_k(x) + \pv\int_{\partial\Omega} \Phi_k(x - y) f_k(y) \d \sigma(y) \nonumber
  \end{align}
  for almost everywhere $x \in \partial\Omega$.
  The subscripts $+$ and $-$ indicate nontangential limits taken inside $\Omega$ and outside $\overline\Omega$, respectively.
\end{lem}

\begin{proof}
  The correctness of the trace formulas \eqref{eq:traceFormula} is known for the case $\lambda = 0$ since Fabes, Kenig and Verchota \cite{fabesKenigVerchota}.
  This fact will now be reused for $\lambda \in \Sigma_\theta$.
  We insert a $0$ to the nontangential limit as
  \begin{align*}
    (\nabla u_j)_\pm(x) = 
    (\nabla v_j)_\pm(x) + (\nabla u_j - \nabla v_j)_\pm(x),
  \end{align*}
  where $v_j(x) = \int_{\partial\Omega} \Gamma_{jk}(x - y; 0) f_k(y) \d\sigma(y)$.
  Because of \cite{fabesKenigVerchota} we know that the first nontangential limit exists and is given by \eqref{eq:traceFormula} with $\lambda = 0$.
  It therefore remains to show that
  \begin{align*}
    (\nabla u_j - \nabla v_j)_\pm(x) = \int_{\partial\Omega} \nabla \{ \Gamma_{jk}(x - y; \lambda) - \Gamma_{jk}(x - y; 0) \} f_k(y) \d\sigma(y)
  \end{align*}
  for all $x \in \partial \Omega$.
  To this end let $(x_l)_{l \in \N}$ a sequence in $\ShenCone(x)$ with $\lim_{l \to \infty} x_l = x$.
  Furthermore let us note that for almost everywhere $x \in \partial\Omega$ we have that 
  \begin{align*}
    \int_{\partial\Omega} \frac{1}{|x - y|^{d - 2}} |f(y)| \d\sigma(y) < \infty.
  \end{align*}
  This is a consequence of the fact that
  \begin{align*}
    \sup_{x \in \partial\Omega} \Big| \int_{\partial\Omega} \frac{1}{|x - y|^{d - 2}} \d\sigma(y) \Big| < \infty
  \end{align*}
  and an application of Young's inequality which can be found in Tolksdorf \cite{tolksdorfDiss}:
  Let $x \in \partial\Omega$.
  Then
  \begin{align*}
    &\int_{\partial\Omega} \frac{1}{|x - y|^{d - 2}} \d\sigma(y) \\
    &\quad\leq \int_{\partial\Omega \cap \BB(x,r_0/4)} \frac{1}{|x - y|^{d - 2}} \d\sigma(y) + \int_{\partial\Omega \setminus \BB(x, r_0/4)} \frac{1}{|x - y|^{d - 2}} \d\sigma(y) \\
    &\quad\leq C r_0 + r^{2 - d} 4^{d - 2} \sigma(\partial\Omega)
  \end{align*}
  by Lemma \ref{lem:compareBoundaryWithBall}.
  Now Young's inequality gives us the desired result.
  In the next step we will show that
  \begin{align*}
    \frac{1}{|x -y|^{d - 2}} |f(y)|
  \end{align*}
  gives a suitable function for dominated convergence.
  Set $\varepsilon = (4 |\lambda|^2)^{-1}$ and without loss of generality assume that $\supp f \subseteq \BB(x,\varepsilon)$.
  Furthermore assume that $|x_l - x| < \varepsilon$ for all $l \in \N$.
  Then $|x_l - y| \leq (2|\lambda|^2)^{-1}$ and Corollary \ref{cor:differenceFundamentalSolutionStokes} give
  \begin{align*}
    & (\nabla u_j - \nabla v_j)(x_l) \int_{\partial\Omega} \nabla \{ \Gamma_{jk}(x_l - y; \lambda) - \Gamma_{jk}(x_l - y; 0)\} f_k(y) \d\sigma(y) \\
    &\quad\leq \int_{\partial\Omega} \frac{1}{\sqrt{|\lambda|}|x_l - y|^{d - 2}} |f(y)| \d\sigma(y) \\
    &\quad\leq \frac{C}{\sqrt{|\lambda|}} \int_{\partial\Omega} \frac{1}{|x - y|^{d - 2}} |f(y)| \d\sigma(y) < \infty.
  \end{align*}
  Now dominated convergence gives the claim for $x_l \to x$.
  Note that it does not affect the proof if the sequence $x_l$ lays inside $\Omega$ or outside $\overline\Omega$.
\end{proof}

The previous Lemma enables us to talk about boundary values of partial derivatives. 
The next theorem will now give a similar result but for conormal derivatives which are defined by
\begin{align*}
  \frac{\partial u}{\partial \nu} = \frac{\partial u}{\partial n} - \phi n.
\end{align*}
We will also be working with tangential derivatives which are defined via
\begin{align*}
  DEFINE
\end{align*}

\begin{thm}
  \label{thm:jumpConditions}
  Let $\lambda \in \Sigma_\theta$ and $\Omega$ be a bounded Lipschitz domain in $\R^d$, $d \geq 3$. 
  Let $(u,\phi)$ be given by \eqref{eq:defSingleLayer} and \eqref{eq:defSingleLayerPressure} with $f \in \Ell^p(\partial\Omega; \C^d)$ and $1 < p < \infty$.
  Then $\nabla_{\mathrm{tan}} u_+ = \nabla_{\mathrm{tan}} u_-$ and
  \begin{align}
    \label{eq:nontangentialConormalDerivative}
    \Big( \frac{\partial u}{\partial \nu} \Big)_\pm = \Big( \pm \frac{1}{2} I + \K_\lambda \Big) f
  \end{align}
  on $\partial\Omega$, with $\K_\lambda$ a bounded operator on $\Ell^p(\partial\Omega; \C^d)$ with
  \begin{align*}
    \| \K_\lambda f \|_{\Ell^p(\partial\Omega)} \leq C_p \|f\|_{\Ell^p(\partial\Omega)},
  \end{align*}
  where $C_p$ depends only on $d$, $\theta$, $p$ and the Lipschitz character of $\Omega$.
\end{thm}

\begin{proof}
  For the the $j$th component of the tangential derivative of $u_i$, $1\leq i,j \leq d$, we calculate using the results from Lemma \ref{lem:traceFormulas}
  \begin{align*}
    ((\nabla_{\mathrm{tan}} u_i)_+)_j
    &= (\frac{\partial u_i}{\partial x_j})_+ - \langle (\nabla u_i)_+, n \rangle n_j \\
    &= (\frac{\partial u_i}{\partial x_j} )_+ - (\frac{\partial u_i}{\partial x_k} )_+ n_k n_j \\
    &= \frac{1}{2} \{ n_j f_i - n_i n_j n_k f_k\} - \frac{1}{2} \{ n_k f_i - n_i n_k n_l f_l \} n_k n_j  \\
    &\quad+ \pv \int_{\partial\Omega} \frac{\partial}{\partial x_j} \{ \Gamma_{ik} (x - y; \lambda) \} f_k(y) \d\sigma(y) \\
    &\quad+ \pv \int_{\partial\Omega} \frac{\partial}{\partial x_k} \{ \Gamma_{il} (x - y; \lambda) \} f_l(y) \d\sigma(y) n_k n_j.
  \end{align*}
  As the first two summands add up to zero, the entire expression does not depend on the direction of the nontangential limit. 
  This gives
  \begin{align*}
    (\nabla_{\mathrm{tan}} u)_+ = (\nabla_{\mathrm{tan}} u)_-
  \end{align*}
  We calculate for the $j$th component of the nontangential limit of the conormal derivative of $u$ at $x \in \partial\Omega$ using the results from Lemma \ref{lem:traceFormulas}
  \begin{align*}
    &\big(\frac{\partial u_j}{\partial x_i}\big)_+(x) n_i - \phi_+(x) n_j\\
    &\quad= \frac{1}{2} \{ n_i f_j(x) - n_j n_i n_k f_k(x) \} n_i + \pv \int_{\partial\Omega} \frac{\partial}{\partial{x_i}} \big\{ \Gamma_{jk}(x - y; \lambda) \big\} f_k(y) \d\sigma(y) n_i \\
    &\qquad+ \frac{1}{2} n_k f_k(x) n_j - \pv \int_{\partial\Omega} \Phi_k(x - y) f_k(y) \d\sigma(y) n_j \\
    &\quad= \frac{1}{2} f_j(x) + (\K_\lambda f)_j(x),
  \end{align*}
  where $n$ denotes the normal vector at $x$ and
  \begin{align}
    \label{eq:defnKlambda}
    (\K_\lambda f) (x)
    &= \pv \int_{\partial\Omega} \nabla_x \Gamma(x - y; \lambda)  f(y) \d\sigma(y) n - \pv \int_{\partial\Omega} \Phi_k(x - y) f_k(y) \d\sigma(y) n.
  \end{align}
  We note that $\K_\lambda$ essentially consists of two boundary layer potentials. 
  The $\Ell^p$-boundedness of the first one was proven in Lemma \ref{lem:lpBoundednessT}.
  The $\Ell^p$-boundedness of the second boundary layer potential follows in an analogous way using the fact that the operators
  \begin{align*}
    A^*(f)(P) = \sup_{t > 0} \Big| \int_{\substack{y \in \partial\Omega \\ |y - P| > t}} \frac{P - y}{|P - y|^n} f(y) \d\sigma(y)\Big| , \quad P \in \partial\Omega
  \end{align*}
  are bounded by Lemma 1.2 of Verchota \cite{verchotaDiss}.
\end{proof}

Similar to $\K_\lambda$ for $\lambda = 0$ we have
\begin{align}
  \label{eq:defnK0}
    (\K_0 f)(x)= \pv \int_{\partial\Omega} \nabla_x \Gamma(x - y; 0)  f(y) \d\sigma(y) n - \pv \int_{\partial\Omega} \Phi_k(x - y) f_k(y) \d\sigma(y) n,
\end{align}
as was shown by Fabes, Kenig and Verchota \cite[(0.12)]{fabesKenigVerchota}.

We now note a fact that will be crucial for solving the $\Ell^2$ Dirichlet problem in Chapter 5 and will fortify the hopes of translating results for $\lambda = 0$ to $\lambda \in \Sigma_\theta$.
\begin{lem}
  \label{lem:compactness}
  Let $\lambda \in \Sigma_\theta$ and $d \geq 3$.
  Then the operator $\K_\lambda - \K_0$ on $\Ell^2(\partial\Omega; \C^d)$ is compact.
\end{lem}

\begin{proof}
  The idea of this proof is similar to the one in Tolksdorf \cite[Lemma 4.3.5]{tolksdorfDiss}.
  Let $f \in \Ell^2(\partial\Omega; \C^d)$.
  Let's denote $\K \coloneqq \K_\lambda - \K_0$. 
  We will now try to approximate $\K$ by compact operators in the operator norm.
  To this end we define for all $\varepsilon > 0$
  \begin{align*}
    (\K^{(\varepsilon)}f)(x) \coloneqq \int_{\partial\Omega \setminus \BB(x,\varepsilon)} \nabla \{ \Gamma(x - y; \lambda) - \Gamma(x - y; 0) \} f(y) \d \sigma(y), \quad x \in \partial\Omega.
  \end{align*}
  We can now estimate by Young's inequality \ref{lim:young} 
  \begin{align*}
    \| (\K^{(\varepsilon)} f)(x) \|_{\Ell^2(\partial\Omega)}
    \leq \sup_{p \in \partial\Omega} \| \nabla \{ \Gamma(p - \cdot; \lambda) - \Gamma(p - \cdot; 0) \} 1_{\BB(p,\varepsilon)} \|_{\Ell^1(\partial\Omega)} \|f\|_{\Ell^2(\partial\Omega)}.
  \end{align*}
  Our goal is to show that
  \begin{align*}
    \sup_{p \in \partial\Omega} \| \nabla \{ \Gamma(p - \cdot; \lambda) - \Gamma(p - \cdot; 0) \} 1_{\BB(p,\varepsilon)} \|_{\Ell^1(\partial\Omega)} \to 0 \quad \text{as} \quad \varepsilon \to 0.
  \end{align*}
  To this end, let $\varepsilon$ be small enough such that we can apply the estimates from Corollary \ref{cor:differenceFundamentalSolutionStokes} to calculate for some $p \in \partial\Omega$
  \begin{align*}
     &\| \nabla \{ \Gamma(p - \cdot; \lambda) - \Gamma(p - \cdot; 0) \} 1_{\BB(p,\varepsilon)} \|_{\Ell^1(\partial\Omega)}  \\
     &\qquad\qquad\leq C \int_{\partial\Omega \cap \BB(p, \varepsilon)} \sqrt{|\lambda|} |p - y|^{2 - d} \d\sigma(y) 
     \leq C \sqrt{|\lambda|} \varepsilon
  \end{align*}
  where for the last step we applied Lemma \ref{lem:comparability}.
  For $\varepsilon \to 0$ this gives us $\K^{(\varepsilon)} \to \K$ in the operator norm.
  The last step is to verify the compactness of $\K^(\varepsilon)$.
  We note that the integral kernel of $\K^{(\varepsilon)}$ is bounded which gives us that in particular the kernel is an element of $\Ell^2(\partial\Omega \times \partial\Omega; \C^{d \times d}$.
  The compactness of $\K^{(\varepsilon)}$ now follows from Weidmann \cite[Thm. 6.11]{weidmann}.
\end{proof}

Our next step is to introduce the \emph{double layer potential} $u(x) = \dlp_\lambda(f)(x)$ for the Stokes resolvent problem, where
\begin{align}
  \label{eq:defDoubleLayer}
  u_j(x) = \int_{\partial\Omega} \Big\{ \frac{\partial}{\partial y_i} \{ \Gamma_{jk}(y - x; \lambda) \} n_i(y) - \Phi_j(y - x) n_k(y) \Big\} f_k(y) \d\sigma(y).
\end{align}
The corresponding pressure $\phi(x) = \dlp_{\phi}(f)(x)$ is defined via
\begin{align}
  \label{eq:defDoubleLayerPressure}
  \phi(x)
  &= \frac{\partial^2}{\partial x_i \partial x_k} \int_{\partial\Omega} G(y - x; 0) n_i(y) f_k(y) \d\sigma(y) + \lambda \int_{\partial\Omega} G(y - x; 0) n_k(y) f_k(y) \d\sigma(y).
\end{align}
Using \ref{eq:fundamentalVectorPressure} and \ref{eq:solutionStokesSystem} one can show that $(u,\phi)$ defines again a solution to the Stokes resolvent problem in $\R^d \setminus \partial\Omega$.

The next theorem will give us a suitable operator which maps a given function $f \in \Ell^p(\partial\Omega; \C^d)$ to boundary values of $u = \dlp_\lambda(f)$.

\begin{thm}
  \label{thm:nontangentialLimitDoubleLayer}
Let $\lambda \in \Sigma_\theta$ and $\Omega$ be a bounded Lipschitz domain in $\R^d$, $d \geq 3$.
Let $u$ be given by \eqref{eq:defDoubleLayer} for $f \in \Ell^p(\partial\Omega; \C^d)$, $1 < p < \infty$.
Then
\begin{align}
  \label{eq:lpBoundednessUNontangentialMax}
  \|(u)^*\|_{\Ell^p(\partial\Omega)} \leq C_p \|f\|_{\Ell^p(\partial\Omega)}
\end{align}
where $C_p$ depends only on $d$, $p$, $\theta$ and the Lipschitz character of $\Omega$.
Furthermore 
\begin{align}
  \label{eq:nontangentialLimitDoubleLayer}
  u_\pm = \Big(\mp \frac{1}{2} I + \K_{\bar\lambda}^* \Big) f,
\end{align}
where $K_{\bar\lambda}^*$ is the adjoint of the operator $K_{\bar\lambda}$ in \eqref{eq:nontangentialConormalDerivative}
\end{thm}

\begin{proof}
  The estimate for $(u)^*$ is a direct consequence of Lemma \ref{lem:nontangentialMaximalFunctions}, in particular of the estimates on $(\nabla u)^*$ and $(\phi)^*$.

  For the proof of \eqref{eq:nontangentialLimitDoubleLayer}, we begin by determining the adjoint of the operator $\K_{\bar\lambda}$.
  To this end we fill first work with truncated operators $\K_{\lambda}^{(\varepsilon)}$ which are defined as
  \begin{align*}
    (\K_\lambda^{(\varepsilon)} f) (x)
    &= \int_{\partial\Omega} 1_{\EE(x,\varepsilon)} \nabla_x \Gamma(x - y; \lambda)  f(y) \d\sigma(y) n - \int_{\partial\Omega} 1_{\EE(x,\varepsilon)} \Phi_k(x - y) f_k(y) \d\sigma(y) n,
  \end{align*}
  for $x \in \partial\Omega$ and $\EE(x,\varepsilon) \coloneqq \R^d \setminus \BB(x,\varepsilon)$.
  Now for $f \in \Ell^p(\partial\Omega; \C^d)$ and $g \in \Ell^q(\partial\Omega; \C^d)$ with $1/p + 1/q = 1$ we calculate
  \begin{align*}
    \langle \K_{\bar\lambda}^{(\varepsilon)}f, g \rangle
    &= \int_{\partial\Omega} (\K_{\bar\lambda}^{(\varepsilon)} f_j)(x) \overline{g_j(x)} \d\sigma(x) \\
    &= \int_{\partial\Omega} \int_{\partial\Omega} \frac{\partial}{\partial x_i} \{ \Gamma_{jk} (x - y; \overline\lambda) \} f_k(y) 1_{\EE(x,\varepsilon}(y) \d\sigma(y) n_i(x) \overline{g_j(x)} \d\sigma(x) \\
    &\quad+ \int_{\partial\Omega} \int_{\partial\Omega} \Phi_k(x -y)f_k(y) 1_{\EE(x,\varepsilon)}(y) \d\sigma(y) n_j(x) \overline{g_j(x)} \d\sigma(x).
  \end{align*}
  Note that $1_{\EE(x,\varepsilon)}(y) = 1_{\EE(y,\varepsilon)}(x)$.
  Now an application of Fubini and factoring out $f_k(y)$ gives that the lengthy expression is equal to
  \begin{align*}
    \int_{\partial\Omega} f_k(y) \int_{\partial\Omega} &\Big\{ \frac{\partial}{\partial x_i} \{ \Gamma_{jk}(x - y; \bar\lambda)\} n_i(x)
    - \Phi_k(x - y)  n_j(x) \Big\}1_{\EE(y, \varepsilon)}(x) \overline{g_j(x)} \d\sigma(x) \d\sigma(y).
  \end{align*}
  Therefore we see that the adjoint of the truncated operator $\K_{\bar\lambda}^{(\varepsilon)}$ is given by
  \begin{align*}
    \big( ( K_{\bar\lambda}^{(\varepsilon)})^* g\big)_k(y)
    = \int_{\partial\Omega} \Big\{ \frac{\partial}{\partial x_i} \{ \Gamma_{jk}(x - y; \lambda)\} n_i(x) - \Phi_k(x - y) n_j(x) \Big\} 1_{\EE(y,\varepsilon)}(x) g_j(x) \d\sigma(x), 
  \end{align*}
  for $y \in \partial\Omega$ since $\overline{\Gamma_{jk}(x - y; \lambda)} = \Gamma_{jk}(x - y; \bar\lambda)$.

  In the next step we will go from truncated operators to principal value operators. 
  For this to work we will look for suitable majorants.
  If $x \in \partial\Omega$ we estimate
  \begin{align*}
    |(\K_{\bar\lambda}^{(\varepsilon)}f)_j (x)|
    &= \Big| \int_{|x - y| > \varepsilon} \frac{\partial}{\partial x_i} \{ \Gamma_{jk}(x - y; \lambda) \} f_k(y) \d\sigma(y) n_i(x)  \\
    &\qquad - \int_{|x - y| > \varepsilon} \Phi_k(x - y) f_k(y) n_j(x) \d\sigma(y) \Big| \\
    &\leq T_\lambda^*(f)(x) + A^*(f)(x).
  \end{align*}
  Now dominated convergence gives
  \begin{align*}
    \lim_{\varepsilon \to 0} \langle K_{\bar\lambda}^{(\varepsilon)} f, g \rangle = \langle K_{\bar\lambda} f, g\rangle.
  \end{align*}
  A similar argument gives
  \begin{align*}
    \lim_{\varepsilon \to 0} \langle f, {K_{\bar\lambda}^{(\varepsilon)}}^{(*)} g \rangle = \langle f,K_{\bar\lambda}^{(*)}  g\rangle,
  \end{align*}
  where
  \begin{align*}
    \big( ( K_{\bar\lambda}^* g\big)_k(y)
    = \pv\int_{\partial\Omega} \Big\{ \frac{\partial}{\partial x_i} \{ \Gamma_{kj}(x - y; \lambda)\} n_i(x) - \Phi_k(x - y) n_j(x) \Big\} g_j(x) \d\sigma(x). 
  \end{align*}
  Note that we have used the symmetry of $(\Gamma_{\alpha\beta})$.

  The last part now consists of proving that the equality \eqref{eq:nontangentialLimitDoubleLayer} holds.
  To simplify the calculations and make Lemma \ref{lem:traceFormulas} more accessible note that on the one hand
  \begin{align*}
    \int_{\partial\Omega} \frac{\partial}{\partial y_i} \{ \Gamma_{jk}(y - x; \lambda) \} n_i(y) f_k(y) \d\sigma(y)
    &= -\int_{\partial\Omega} \frac{\partial}{\partial x_i} \{ \Gamma_{jk}(x - y; \lambda) \} n_i(y) f_k(y) \d\sigma(y)\\
    &= -\frac{\partial}{\partial x_i} \slp(n_i f)_j(x)
  \end{align*}
  and on the other hand
  \begin{align*}
    -\int_{\partial\Omega} \Phi_j(y - x) n_k(y) f_k(y) \d\sigma(y)
    = \int_{\partial\Omega} \Phi_l(x - y) \delta_{lj} n_k(y) f_k(y) \d\sigma(y)
    = \slp_\Phi(\tilde f^j)(x), 
  \end{align*}
  where $\tilde f^j_l = \delta_{lj} n_k f_k$.
  For $x \in \partial\Omega$ we can now calculate
  \begin{align*}
    &\Big( \int_{\partial\Omega} \frac{\partial}{\partial y_i} \{ \Gamma_{jk}(y - \cdot\, ; \lambda) \} n_i(y) f_k(y) \d\sigma(y) \Big)_\pm(x)\\
    &\qquad= - \big( \frac{\partial}{\partial x_i} \slp_\lambda(n_i f)_j\big)_\pm(x) \\
    &\qquad= \mp \frac{1}{2} \{ n_i(x) n_i(x) f_j(x) - n_j(x) n_i(x) n_k(x) n_i(x) f_k(x)  \} \\
    &\qquad\quad - \pv\int_{\partial\Omega} \frac{\partial}{\partial x_i} \{\Gamma_{jk}(x - y; \lambda) \} n_i(y) f_k(y) \d\sigma(y) \\
    &\qquad= \mp \frac{1}{2} \{ f_j(x) - n_j(x) n_k(x) f_k(x) \} \\
    &\qquad\quad + \pv\int_{\partial\Omega} \frac{\partial}{\partial y_i} \{\Gamma_{jk}(x - y; \lambda) \} n_i(y) f_k(y) \d\sigma(y),
  \end{align*}
  where we used trace formula \eqref{eq:traceFormula}.
  A similar procedure for the second integral part of the double layer potential gives
  \begin{align*}
    &-\Big(\int_{\partial\Omega} \Phi_j(y - \cdot\,) n_k(y) f_k(y) \d\sigma(y)\Big)_\pm(x) \\
    &\qquad= \big(\slp_\Phi(\tilde f^j)\big)_\pm(x)  \\
    &\qquad= \mp \frac{1}{2} n_k(x) \tilde f^j_k(x) - \pv \int_{\partial\Omega} \Phi_k(x - y) \tilde f^j_k(x) \d\sigma(y) \\
    &\qquad= \mp \frac{1}{2} n_j(x) n_k(x) f_k(x) - \pv \int_{\partial\Omega} \Phi_j(x - y) n_k(x) f_k(x) \d\sigma(y)
  \end{align*}
  Putting everything together we get
  \begin{align*}
    (u_j)_\pm(x) = \mp\frac{1}{2} f_j(x) + (K_{\bar\lambda}^* f)_j(x)
  \end{align*}
  which proves the claim.
\end{proof}

\chapter{Rellich Estimates}
\label{chap:4}

In this section, we will establish Rellich-type estimates for the Stokes resolvent problem which will be used in the following chapter to prove the invertibility of the operators $\pm(1/2)I + \K_{\lambda}$ and their adjoints from Theorems  \ref{thm:jumpConditions} and \ref{thm:nontangentialLimitDoubleLayer}.

We will for this entire section always assume that $\Omega$ is a bounded Lipschitz domain in $\R^d$, $d \geq 2$, with connected boundary.
Furthermore, we will use the shorthand notation
\begin{align*}
    \| \cdot \|^{}_{\partial} \coloneqq \| \cdot \|_{\Ell^2(\partial\Omega; \C^k)}, \quad k \in \N,
\end{align*}
and we will tacitly use the summation convention whenever it is applicable.

The following Theorem states formulates the aforementioned Rellich estimates and is the central result of this chapter:

\begin{thm}
  \label{thm:rellich}
  Let $\lambda \in \Sigma_\theta$ and $|\lambda| \geq \tau$, where 
  $\tau \in (0,1)$.
  Let $(u,\phi)$ be a smooth solution to the Stokes resolvent problem in $\Omega$ and suppose that $(\nabla u)^* \in \Ell^2(\partial\Omega)$ and $(\phi)^* \in \Ell^2(\partial\Omega)$.
  Furthermore, assume that $\nabla u$, $\phi$ have nontangential limits almost everywhere on $\partial\Omega$.
  Then
  \begin{align}
    \label{eq:rellich1}
    \begin{alignedat}{1}
        &\|\nabla u\|^{}_\partial + \bigg\|\phi - \bigg\{ \frac{1}{{r_0}^{d - 1}}\int_{\partial\Omega} \phi \d\sigma \bigg\} \bigg\|^{}_\partial \\[1em]
      &\hspace{2cm} \leq C\, \bigg\{ \big\|\nabla_{\mathrm{tan}} u\big\|^{}_\partial + |\lambda|^{1/2} \|u\|^{}_\partial + |\lambda| \, \|u \cdot n\|_{\HH^{-1}(\partial\Omega)} \bigg\}
    \end{alignedat}
  \end{align}
  and
  \begin{align}
    \label{eq:rellich2}
    \|\nabla u\|^{}_\partial + |\lambda|^{1/2} \|u\|^{}_\partial + |\lambda| \, \|u \cdot n\|_{\HH^{-1}(\partial\Omega)} + \|\phi\|^{}_\partial
    \leq C \, \Big\|\frac{\partial u}{\partial \nu} \Big\|^{}_\partial,
  \end{align}
  where $\frac{\partial}{\partial \nu}$ denotes the conormal derivative, and $C$ depends only on $d$, $\tau$, $\theta$ and the Lipschitz character of $\Omega$.
\end{thm}

\begin{rem}
  \label{rem:shenNontangential}
  The assumptions on $u$ in Theorem \ref{thm:rellich} are sufficient for $u$ to have a nontangential limit and a square integrable maximal function $(u)^*$. 
  Indeed for $d = 2$ we have $(u)^* \in \Ell^\infty(\partial\Omega)$, for $d = 3$ we have $(u)^* \in \Ell^p(\partial\Omega), p \in (1,\infty)$, and for $d \geq 3$ we have $(u)^* \in \Ell^p(\partial\Omega)$, $p \in \big(1, 2 (d - 1) / (d - 3) \big)$.
  A proof of these facts can be found in Shen's notes \cite[Prop. 7.1.3]{Shen2017}.
\end{rem}

We will now prepare the proof of Theorem \ref{thm:rellich} by proving several helpful lemmata.
The first lemma deals with so called \emph{Rellich identities} for solutions to the Stokes resolvent system.

\begin{lem}
  \label{lem:rellichIdentity}
  Under the same conditions on $(u,\phi)$ as in Theorem \ref{thm:rellich}, we have
  \begin{align}
    \int_{\partial\Omega} h_k n_k |\nabla u|^2 \d\sigma 
    &= 2 \Ret \int_{\partial\Omega} h_k \frac{\partial \overline u_i}{\partial x_k} \bigg( \frac{\partial u}{\partial \nu} \bigg)_i \d\sigma + \int_\Omega \div(h)\, |\nabla u|^2 \d x \nonumber\\[0.5em]
    &\quad - 2 \Ret \int_\Omega \frac{\partial h_k}{\partial x_j} \cdot \frac{\partial u_i}{\partial x_k} \cdot \frac{\partial \overline u_i}{\partial x_j} \d x + 2 \Ret \int_\Omega \frac{\partial h_k}{\partial x_i} \cdot \frac{\partial u_i}{\partial x_k} \, \overline \phi \d x \nonumber\\[0.5em]
    &\quad - 2 \Ret \int_\Omega h_k \frac{\partial u_i}{\partial x_k} \cdot \overline {\lambda u_i} \d x \label{eq:rellichIdentity}
  \end{align}
  and
  \begin{align}
    \int_{\partial\Omega} h_k n_k |\nabla u|^2 \d\sigma
    &= 2\Ret \int_{\partial\Omega} h_k \frac{\partial \overline u_i}{\partial x_j} \bigg\{ n_k \frac{\partial u_i}{\partial x_j} - n_j \frac{\partial u_i}{\partial x_k} \bigg\} \d\sigma \nonumber\\[0.5em]
    &\quad + 2 \Ret \int_{\partial\Omega} h_k \, \overline \phi \, \bigg\{ n_i \frac{\partial u_i}{\partial x_k} - n_k \frac{\partial u_i}{\partial x_i} \bigg\} \d\sigma - \int_\Omega \div(h)\, |\nabla u|^2 \d x \nonumber\\[0.5em]
    &\quad + 2 \Ret \int_\Omega \frac{\partial h_k}{\partial x_j} \cdot \frac{\partial u_i}{\partial x_k} \cdot \frac{\partial \overline u_i}{\partial x_j} \d x - 2 \Ret \int_\Omega \frac{\partial h_k}{\partial x_i} \cdot \frac{\partial u_i}{\partial x_k} \, \overline \phi \d x \nonumber\\[0.5em]
    &\quad + 2 \Ret \int_\Omega h_k \frac{\partial u_i}{\partial x_k} \cdot \overline{ \lambda u_i} \d x\,, \label{eq:rellichIdentity2}
  \end{align}
  where $h = (h_1, \dots,h_d) \in \CC_0^1(\R^d, \R^d)$.
\end{lem}

\begin{proof}
  The proof of the stated identities reduces to several applications of the divergence theorem once we establish its applicability.
  To this end, we want to make Proposition \ref{prop:approximationArgument} available. We note that the assumptions given in Theorem \ref{thm:rellich} are sufficient for this purpose and we will verify them, once they are used.

  Let's expand the first summand in \eqref{eq:rellichIdentity} using the definition of conormal derivatives
  \begin{align*}
    2\Ret \int_{\partial\Omega} h_k \frac{\partial \overline u_i}{\partial x_k} \bigg( \frac{\partial u}{\partial \nu} \bigg)_i \d\sigma
    &= 2 \Ret \int_{\partial\Omega} h_k \frac{\partial \overline u_i}{\partial x_k} \cdot \frac{\partial u_i}{\partial x_j}  n_j \d\sigma  - 2\Ret \int_{\partial\Omega} h_k \frac{\partial \overline u_i}{\partial x_k} \, \phi\, n_i \d x \\
    &\eqqcolon I_1 - I_2.
  \end{align*}
  The divergence theorem is applicable for $I_1$ as $h$ is bounded and defined everywhere and the integrand has nontangential limits that can be dominated by $|(\nabla u)^*|^2 \in \Ell^2(\partial\Omega)$.
  Therefore, we find using the divergence theorem and the product rule:
  \begin{align*}
    I_1
    &= 2\Ret \int_\Omega \frac{\partial}{\partial x_j} \bigg\{ h_k \frac{\partial \overline u_i}{\partial x_k} \cdot \frac{\partial u_i}{\partial x_j} \bigg\} \d x \\[0.5em]
    &= 2\Ret \int_\Omega \frac{\partial h_k}{\partial x_j} \cdot \frac{\partial \overline u_i}{\partial x_k} \cdot \frac{\partial u_i}{\partial x_j} \d x 
    + 2\Ret \int_{\Omega} h_k \frac{\partial^2 \overline u_i}{\partial x_j \partial x_k} \cdot \frac{\partial u_i}{\partial x_j} \d x  \\[0.5em]
    &\hphantom{= 2\Ret \int_\Omega \frac{\partial h_k}{\partial x_j} \cdot \frac{\partial \overline u_i}{\partial x_k} \cdot \frac{\partial u_i}{\partial x_j} \d x } \;
    + 2 \Ret \int_\Omega h_k \frac{\partial \overline u_i}{\partial x_k} \cdot \frac{\partial^2 u_i}{\partial x_j^2} \d x \\
    &\eqqcolon I_3 + I_4 + I_5.
  \end{align*}
  For $I_5$, we use the fact that $u$ solves the Stokes resolvent problem which gives
  \begin{align*}
     I_5 
     &= 2 \Ret \int_\Omega h_k \frac{\partial \overline u_i}{\partial x_k}\cdot  \frac{\partial \phi}{\partial x_i} \d x + 2\Ret \int_\Omega h_k \frac{\partial\overline u_i}{\partial x_k} \lambda u_i \d x \\
     &\eqqcolon I_6 + I_7 .
  \end{align*}
  Now we want to apply the divergence theorem, i.e. Proposition \ref{prop:approximationArgument} to integral $I_2$.
  This is possible since $h$ is defined everywhere and bounded, $(\partial_k u_i) \cdot \phi$ has a nontangential limit and can be bounded by $\big( |(\nabla u)^*| |(\phi)^*| \big)$ which is integrable due to Hölder's inequality as $(\nabla u)^*$ and $(\phi)^*$ are square integrable by assumption.
  Thus the divergence theorem is applicable and yields together with the product rule:
  \begin{align*}
    I_2
    &= 2 \Ret \int_\Omega \frac{\partial}{\partial x_i} \Big\{ h_k \frac{\partial \overline u_i}{\partial x_k} \phi \Big\} \d x \\[0.5em]
    &= 2 \Ret \int_\Omega \frac{\partial h_k}{\partial x_i} \cdot \frac{\partial \overline u_i}{\partial x_k} \, \phi \d x + 2 \Ret \int_\Omega h_k \frac{\partial^2 \overline u_i}{\partial x_j \partial x_k} \, \phi \d x + 2 \Ret \int_\Omega h_k \frac{\partial \overline u_i}{\partial x_k} \cdot \frac{\partial \phi}{\partial x_i} \d x\\
    &\eqqcolon I_8 + I_9 + I_6.
  \end{align*}
  One term that hasn't come up so far, the second summand of the right hand side in \eqref{eq:rellichIdentity}, will now be expanded:
  \begin{align*}
    \int_\Omega \div(h)\,  |\nabla u|^2 \d x
    &= \int_\Omega \div( h\,|\nabla u|^2) \d x - \int_\Omega h_k \frac{\partial}{\partial x_i} \Big\{ |\nabla u|^2 \Big\} \d x \\
    &\eqqcolon I_{10} - I_{11}.
  \end{align*}
  Expanding the Integral $I_{11}$ gives us the identity
  \begin{align*}
    I_{11}
    &= \int_\Omega h_i \frac{\partial}{\partial x_i} \bigg\{ \frac{\partial u_k}{\partial x_j} \cdot \frac{\partial \overline u_k}{\partial x_j} \bigg\} \d x
    = \int_\Omega h_i  \bigg\{ \frac{\partial^2 u_k}{\partial x_i \partial x_j} \cdot \frac{\partial \overline u_k}{\partial x_j} + \frac{\partial u_k}{\partial x_j} \cdot \frac{\partial^2 \overline u_k}{\partial x_i \partial x_j} \bigg\} \d x= I_4.
  \end{align*}
  If we now put everything together, the right hand side of \eqref{eq:rellichIdentity} reads
  \begin{align*}
    &(I_1 - I_2) + (I_{10} - I_{11}) - I_3 + I_8 - I_7 \\
    &\quad= (I_3 + I_4 + I_6 + I_7) - (I_8 + I_9 + I_6) + I_{10} - I_{11} - I_3 + I_8 - I_7 = I_{10}.
  \end{align*}
  Noting that by the divergence theorem, which is applicable with the same justification as for the integral $I_1$, we have
  \begin{align*}
    I_{10} = \int_{\partial\Omega} h_k n_k |\nabla u|^2 \d \sigma.
  \end{align*}
  Thus, the first identity is proven.

  In order to prove identity \eqref{eq:rellichIdentity2}, we show that the expression we get from considering (\eqref{eq:rellichIdentity} + \eqref{eq:rellichIdentity2}) holds, i.e. we show the identity
  \begin{align*}
    2 \int_{\partial\Omega} h_k n_k |\nabla u|^2 \d\sigma
    &= 2 \Ret \int_{\partial\Omega} h_k \frac{\partial \overline u_i}{\partial x_k} \bigg( \frac{\partial u}{\partial \nu} \bigg)_i \d \sigma \\[0.5em]
    &\quad + 2 \Ret \int_{\partial\Omega} h_k \frac{\partial\overline u_i}{\partial x_j} \bigg\{ n_k \frac{\partial u_i}{\partial x_j} - n_j \frac{\partial u_i}{\partial x_k} \bigg\} \d \sigma \\[0.5em]
    &\quad+ 2 \Ret \int_{\partial\Omega} h_k \, \overline \phi\,  \bigg\{ n_i \frac{\partial u_i}{\partial x_i} - n_k \frac{\partial u_i}{\partial x_i} \bigg\} \d\sigma.
  \end{align*}
  To this end, note that the left side of the identity equals $2 I_{10}$, whereas the right hand side can be written as 
  \begin{align*}
    (I_1 - I_2) + 2 (I_{10} - I_1) + (I_2 - 0),
  \end{align*}
  where we also used the fact that $\div u = \partial_i u_i = 0$.
\end{proof}

Consider the operators $\partial_{\tau_{jk}}$ which act on compactly supported continuously differentiable functions $\psi$ in the neighborhood of $\partial\Omega$ by
\begin{align}
  \label{eq:defnTangDerivative}
  \partial_{\tau_{jk}} \psi \coloneqq n_j \frac{\partial \psi}{\partial x_k} \bigg|_{\partial\Omega} - n_k \frac{\partial \psi}{\partial x_j} \bigg|_{\partial\Omega}, \quad j,k = 1,\dots,d.
\end{align}
These operators show up in identity \eqref{eq:rellichIdentity2} as
\begin{align*}
  \partial_{\tau_{kj}} u_i = \Big\{ n_k \frac{\partial u_i}{\partial x_j} - n_j \frac{\partial u_i}{\partial x_k} \Big\}
  \quad\text{and}\quad
  \partial_{\tau_{ik}} u_i =  \Big\{ n_i \frac{\partial u_i}{\partial x_k} - n_k \frac{\partial u_i}{\partial x_i} \Big\}
\end{align*}
and have been introduced by Mitrea and Wrigth \cite[p. 16]{mitreaWright}. 
These operators are called \emph{first-order tangential derivative operators} and relate to the tangential gradient, which has been introduced in 
\eqref{eq:tangentialGradient}, in the following way:
\begin{align}
  \label{eq:relTanGrad}
  \big(\nabla_{\mathrm{tan}} \psi\big)_j = \frac{\partial \psi}{\partial x_j} - n_k n_j \frac{\partial \psi}{\partial x_k} = n_k \, \partial_{\tau_{kj}} \psi.
\end{align}
The tangential derivative operators come with a helpful  ``integration by parts'' rule and can be used to define Sobolev spaces on the boundary $\partial\Omega$:
For $f \in \Ell^1_\loc(\Omega)$, we start by defining antilinear functionals on $\CC_0^\infty(\R^d)$ by setting
\begin{align*}
    \partial_{\tau_{kj}} f \colon \CC_0^\infty(\R^d) \ni \psi \mapsto \int_{\partial\Omega} f \, \overline{ \partial_{\tau_{jk}} \psi} \d\sigma.
\end{align*}
Now the weak tangential derivatives of $f$ are given by those functionals which are regular, i.e. elements of $\Ell^1_\loc(\partial\Omega)$.
In this case, the following integration by parts formula holds:
\begin{align*}
    \int_{\partial\Omega} f\, \overline{( \partial_{\tau_{jk}} \psi)} \d\sigma = \int_{\partial\Omega} (\partial_{\tau_{kj}} f)\, \overline{\psi} \d\sigma.
\end{align*}
Note that the minus sign that usually comes with the integration by parts is hidden in the new order of indices.
For $p \in (1,\infty)$, we define the corresponding Sobolev space via
\begin{align*}
  \WW^{1,\,p}(\partial\Omega) = \Big\{ f \in \Ell^p(\partial\Omega) \colon \partial_{\tau_{jk}} f \in \Ell^p(\partial\Omega), \; j,k = 1,\dots,d \Big\}.
\end{align*}

We extend our detour by the following basic lemma on a reverse triangle inequality for elements of the sector $\Sigma_\theta$.
A powerful generalization of this Lemma can be found in Tolksdorf \cite[Lem. 5.2.4]{tolksdorf}.

\begin{lem}
  \label{lem:lambdaIneq}
  Let $\theta \in (0,\pi/2)$.
  Then there exists $\alpha$ depending only on $\theta$ such that for all $\lambda \in \Sigma_\theta$ the following inequality holds:
  \begin{align*}
    \Re(\lambda) + \alpha \, \big|\Im(\lambda)\,\big| \geq |\lambda|.
  \end{align*}
\end{lem}

\begin{proof}
  For the moment being, suppose $|\lambda| = 1$ Then, we have $\Re(\lambda) = \cos(\varphi)$ and $\Im(\lambda) = \sin(\varphi)$ with $|\varphi| \in (0, \pi - \theta)$. 
  Set
  \begin{align*}
    \alpha = \frac{1 - \cos(\pi - \theta)}{\sin(\pi - \theta)} \geq \frac{1 - \cos(|\varphi|)}{\sin(|\varphi|)}.
  \end{align*}
  If $\varphi = |\varphi|$, this gives the inequality
  \begin{align*}
    \Re(\lambda) + \alpha \, |\Im(\lambda)\,| = \cos(\varphi) + \alpha \sin(\varphi) \geq 1.
  \end{align*}
  Conversely, if $\varphi = -|\varphi|$, then we have by the symmetry properties of $\sin$ and $\cos$ that
  \begin{align*}
    \Re(\lambda) + \alpha\, |\Im(\lambda)\,| = \cos(-\varphi) + \alpha \sin(-\varphi) \geq 1.
  \end{align*}
  For arbitrary $\lambda$ the claim follows by considering the normalized value $(\lambda / |\lambda|)$.
\end{proof}

The next lemma enables us to handle the solid integrals in \eqref{eq:rellichIdentity} and \eqref{eq:rellichIdentity2}.

\begin{lem}
  \label{lem:laxMilgramIneq}
  Under the same assumptions on $(u,\phi)$ and $\lambda$ as in Theorem \ref{thm:rellich}, we have
  \begin{align}
    \label{eq:laxMilgramIneq}
    \int_\Omega |\nabla u|^2 \d x + |\lambda| \int_\Omega |u|^2 \leq C \,\Big\| \frac{\partial u}{\partial \nu} \Big\|_\partial  \|u\|_\partial,
  \end{align}
  where $C$ depends only on $\theta$.
\end{lem}

\begin{proof}
  Inserting the solution $u$ into the the variational problem of the Stokes resolvent problem gives us
  \begin{align}
    \label{eq:variationalStokes}
    \int_\Omega -\Delta u \cdot \overline u \d x + \lambda \int_\Omega u \cdot \overline u \d x= - \int_\Omega \nabla \phi \cdot \overline u \d x.
  \end{align}
  Rewriting the first term of equation \eqref{eq:variationalStokes} leads to 
  \begin{align*}
    -\int_{\Omega} \frac{\partial^2 u_j}{\partial x_i \partial x_i}  \; \overline u_j \d x
    = -\int_{\Omega} \frac{\partial }{\partial x_i} \bigg\{ \overline u_j \frac{\partial u_j}{\partial x_i} \bigg\} \d x + \int_\Omega \frac{\partial u_j}{\partial x_i} \cdot \frac{\partial \overline u_j}{\partial x_i} \d x.
    \end{align*}
  Note that since $u$ is solenoidal, we have for the third term of equation \eqref{eq:variationalStokes}
  \begin{align*}
    - \int_{\Omega} \frac{\partial \phi}{\partial x_i} \; \overline u_i \d x = - \int_{\Omega} \frac{\partial}{\partial x_i} \Big\{ \phi\,  \overline u_i \Big\} \d x.
  \end{align*}
  Now we want to transform the first and third of the above solid integrals into boundary integrals through Proposition \ref{prop:approximationArgument}.
  By the assumptions formulated in Theorem \ref{thm:rellich}, $\phi$ and $\nabla u$ have a nontangential limit and for both nontangential maximal functions the inclusion $(\phi)^*, (\nabla u)^* \in \Ell^2(\partial\Omega)$ holds. 
  Furthermore, according to Remark \ref{rem:shenNontangential}, also $u$ has a nontangential limit and the nontangential maximal function satisfies $(u)^* \in \Ell^2(\partial\Omega)$. 
  Therefore, the function $|\phi\,  \overline u_i |$ may be dominated by $|(\phi)^* (u)^*| \in \Ell^2(\partial\Omega)$ and the function $|(\partial_j u_i) u_i|$ may be dominated by $|(\nabla u)^* (u)^*|$, respectively. Thus, the door to Proposition \ref{prop:approximationArgument} has been opened which allows to transform equation \eqref{eq:variationalStokes} into
  \begin{align*}
    \int_\Omega |\nabla u|^2 \d x - \int_{\partial\Omega} \frac{\partial u}{\partial n} \cdot \overline u \d \sigma + \lambda \int_\Omega |u|^2 \d x = - \int_{\partial\Omega}  \phi n \cdot \overline u \, \d \sigma.
  \end{align*}
  We can rearrange the terms of this identity and use the definition of conormal derivatives, see equation \eqref{eq:conormalDerivative}, to derive
  \begin{align}
    \label{eq:testedStokes}
    \int_\Omega |\nabla u|^2 \d x + \lambda \int_\Omega |u|^2 \d x = \int_{\partial\Omega} \frac{\partial u}{\partial \nu} \cdot \overline u \, \d \sigma.
  \end{align}
  If we now take the real and imaginary part of \eqref{eq:testedStokes} and sum them up with the prefactor $\alpha(\theta) > 0$ from Lemma \ref{lem:lambdaIneq}, we get
  \begin{align*}
    \int_\Omega |\nabla u|^2 \d x + \Big\{ \Re(\lambda) + \alpha\, \big|\Im(\lambda)\,\big|\, \Big\} \int_\Omega |u|^2 \d x
    \leq (1 + \alpha) \, \bigg| \int_{\partial\Omega} \frac{\partial u}{\partial \nu} \cdot \bar u \d \sigma \, \bigg|\,.
  \end{align*}
  Lemma \ref{lem:lambdaIneq} now gives
  \begin{align*}
    \int_\Omega |\nabla u|^2 \d x + |\lambda| \int_\Omega |u|^2 \d x \leq C\,  \bigg| \int_{\partial\Omega} \frac{\partial u}{\partial \nu} \cdot \bar u \d \sigma\, \bigg|\,,
  \end{align*}
  with $C = (1 + \alpha)$ from which we readily derive estimate \eqref{eq:laxMilgramIneq} after applying the Cauchy-Schwartz inequality.
\end{proof}

The next lemma combines Rellich identities \eqref{eq:rellichIdentity} and \eqref{eq:rellichIdentity2} with estimate \eqref{eq:laxMilgramIneq}.

\begin{lem}
  Under the same assumptions on $(u,\phi)$ and $\lambda$ as in Theorem \ref{thm:rellich}, we have
  \begin{align}
    \label{eq:gradEstimateRellich}
    \| \nabla u\|_\partial \leq C_\varepsilon\, \,\Big\| \frac{\partial u}{\partial \nu} \Big\|_\partial + \varepsilon \,\Big\{ \|\nabla u\|_\partial + \|\phi\|_\partial +  |\lambda|^{1/2} \| u\|_\partial \Big\}
  \end{align}
  and
  \begin{align}
    \label{eq:gradEstimateRellich2}
    \|\nabla u\|_\partial \leq C_\varepsilon\, \Big\{ \|\nabla_{\mathrm{tan}} u \|_\partial + |\lambda|^{1/2} \| u\|_\partial \Big\} + \varepsilon\, \Big\{ \|\nabla u\|_\partial + \|\phi\|_\partial \Big\}
  \end{align}
  for all $\varepsilon \in (0,1)$, where $C_\varepsilon$ depends only on $d$, $\theta$, $\tau$, $\varepsilon$ and the Lipschitz character of $\Omega$.
\end{lem}

\begin{proof}
  Let $h = (h_1, \dots, h_d) \in \CC_0^1(\R^d, \R^d)$ with $h_k n_k \geq c > 0$ on $\partial \Omega$ as given by Therorem \ref{thm:smoothApproximation} v). 
  The idea of the proof of the desired estimates \eqref{eq:gradEstimateRellich} and \eqref{eq:gradEstimateRellich2} is to first use the Rellich identities from Lemma \ref{lem:rellichIdentity} with this particular $h$ to estimate $\|\nabla u\|_\partial$ and then to bound the resulting right hand side by providing individual estimates.

  Before we start, note that we have $\Delta \phi = 0$ on the one hand and for the nontangential maximal function $(\phi)^* \in \Ell^2(\partial\Omega)$ on the other hand. According to Shen \cite[p. 410]{Shen2012}, a result from Dahlberg \cite{dahlberg77} gives the estimation
  \begin{align}
    \label{eq:dahlbergEstimate}
    \int_\Omega |\phi|^2 \d x \leq C \, \|(\phi)^* \|_\partial^2 \leq C\, \|\phi\|_\partial^2.
  \end{align}

  We will now prove the first estimate \eqref{eq:gradEstimateRellich}.
  In view of identity \eqref{eq:rellichIdentity}, we have
  \begin{align}
    \label{eq:normRellich}
    \begin{alignedat}{1}
    \|\nabla u\|_\partial^2
    &\leq C \, \bigg\{ \|\nabla u\|_\partial \, \Big\| \frac{\partial u}{\partial \nu} \Big\|_\partial + \int_\Omega |\nabla u|^2 \d x \\
    &\hspace{1.5cm}+ \int_\Omega |\nabla u| \, |\phi| \d x + |\lambda| \int_\Omega |\nabla u|\, |u| \d x \bigg\},
    \end{alignedat}
  \end{align}
  where the first term follows from the Cauchy-Schwartz inequality and $C$ only depends on $d$ and the Lipschitz character of $\Omega$.

  For now, we keep the first term of \eqref{eq:normRellich} as it is, the second term can be handled via Lemma \ref{lem:laxMilgramIneq}. 
  The goal for the remaining two integrals will be to bound each of them by a product of norms $\|\cdot\|_\partial$.
To this end, for the third integral we calculate
  \begin{align}
    \int_\Omega |\nabla u| |\phi| \d x
    \leq \Big( \int_\Omega |\nabla u|^2 \d x \Big)^{1/2} \Big( \int_\Omega |\phi|^2 \d x \Big)^{1/2}
    \leq C \, \Big\| \, \frac{\partial u}{\partial \nu}\, \Big\|_\partial^{1/2} \|u\|_\partial^{1/2} \, \|\phi\|_\partial\,, \label{eq:nablaPhi}
  \end{align}
  where the first step is due to the Cauchy-Schwartz inequality and the second step combines estimate \eqref{eq:laxMilgramIneq} with estimate \eqref{eq:dahlbergEstimate}.

  The last integral of \eqref{eq:normRellich} can be estimated as follows:
  \begin{align}
\label{eq:lambdaNablaU}
    |\lambda| \int_\Omega |\nabla u| \, |u| \d x 
    &\leq  \frac{|\lambda|^{3/2}}{2} \int_\Omega |u|^2 \d x + \frac{|\lambda|^{1/2}}{2 } \int_\Omega |\nabla u|^2 \d x \leq C \, \Big\| \, \frac{\partial u}{\partial \nu} \, \Big\|_\partial \, |\lambda|^{1/2}  \big\| u\,\big\|_\partial, 
  \end{align}
  where in the first step we used the weighted Young inequality and in the second step we applied estimate \eqref{eq:laxMilgramIneq}.
  Putting everything together, we calculate
  \begin{align*}
    \| \nabla u\|_\partial^2 
    &\leq C\, \bigg\{ \| \nabla u\|_\partial \, \Big\|\,\frac{\partial u}{\partial \nu} \, \Big\|_\partial 
    + \Big\|\,\frac{\partial u}{\partial \nu}\,\Big\|_\partial \|u\|^{}_\partial 
    + \Big\|\,\frac{\partial u}{\partial \nu}\,\Big\|_\partial^{1/2} \|u\|_\partial^{1/2} \|\phi\|^{}_\partial 
+ \Big\|\,\frac{\partial u}{\partial \nu}\,\Big\|_\partial \, |\lambda|^{1/2} \big\|  u \big\|^{}_\partial\bigg\}
  \end{align*}
  If we now use the assumption $|\lambda| \geq \tau$ which allows us to bound $\|u\|_\partial$ via
  \begin{align*}
      \|u\|^{}_\partial \leq \frac{|\lambda|^{1/2}}{\tau^{1/2}} \|u\|^{}_\partial = C \, |\lambda|^{1/2} \|u\|^{}_{\partial},
  \end{align*}
  the desired estimate \eqref{eq:gradEstimateRellich} now follows applying Young's weighted inequality with an $\varepsilon$ and the norm equivalence on finite dimensional vector spaces. 
  Note that for the product of three norms from inequality \eqref{eq:nablaPhi} we need to apply the Young inequality twice:
  \begin{align*}
      \Big\| \, \frac{\partial u}{\partial \nu} \, \Big\|_\partial^{1/2} \|u\|_\partial^{1/2} \|\phi\|^{}_\partial 
      &\leq \Big\{\, \frac{1}{4 \varepsilon}\, \Big\|\, \frac{\partial u}{\partial \nu}\, \Big\|^{}_\partial 
  + \varepsilon \|u\|^{}_\partial \Big\} \, \|\phi\|^{}_\partial  \\[0.5em]
  &\leq \frac{1}{32\, \varepsilon^3} \, \Big\|\, \frac{\partial u}{\partial \nu} \, \Big\|_\partial^2 + \frac{\varepsilon}{2} \|\phi\|_\partial^2 + \frac{1}{2} \|u\|_\partial^2 + \frac{\varepsilon^2}{2} \|\phi\|_\partial^2 \\[0.5em]
  &\leq C_\varepsilon\, \Big\{ \; \Big\|\, \frac{\partial u}{\partial\nu} \, \Big\|_\partial^2 + \|u\|_\partial^2 \,\Big\} + \varepsilon \|\phi\|_\partial^2,
  \end{align*}
  where for the last inequality we used the fact that $\varepsilon < 1$.

  For inequality \eqref{eq:gradEstimateRellich2}, we use the Rellich identity \eqref{eq:rellichIdentity2} and the relation \eqref{eq:relTanGrad} to obtain the estimate
  \begin{align}
    \label{eq:onTheWay}
    \begin{alignedat}{1}
    \| \nabla u\|_\partial^2 
    &\leq C \bigg\{ \|\nabla_{\mathrm{tan}} u\|^{}_\partial 
    \Big\{ \|\nabla u\|^{}_\partial + \|\phi\|^{}_\partial \Big\}   \\
    &\hspace{2cm} + \int_\Omega |\nabla u|^2 \d x + \int_\Omega |\nabla u| |\phi| \d x + |\lambda| \int_\Omega |\nabla u| |u|\d x \bigg\}.
  \end{alignedat}
  \end{align}
  As before we estimate the three terms on the right side of \eqref{eq:onTheWay} using \eqref{eq:laxMilgramIneq}, \eqref{eq:nablaPhi} and \eqref{eq:lambdaNablaU}, respectively, and obtain the estimate
  \begin{align*}
    \| \nabla u\|_\partial^2 
    &\leq C\, \bigg\{ \|\nabla_{\mathrm{tan}} u\|^{}_\partial 
    \Big\{ \|\nabla u\|^{}_\partial + \|\phi\|^{}_\partial \Big\}   \\
    &\hspace{2cm} + \Big\| \frac{\partial u}{\partial \nu} \Big\|^{}_\partial \| u\|^{}_\partial
+ \Big\| \frac{\partial u}{\partial \nu} \Big\|_\partial^{1/2} \|u\|_\partial^{1/2} \|\phi\|^{}_\partial + \Big\| \frac{\partial u}{\partial \nu} \Big\|^{}_\partial \, |\lambda|^{1/2} \big\| u  \big\|^{}_\partial \bigg\}.
  \end{align*}
  If we now use the Young inequality with an $\varepsilon$, we get
  \begin{align*}
    \|\nabla u\|_\partial^2 
    &\leq C_\varepsilon\, \Big\{ \, \|\nabla_{\mathrm{tan}} u\|_\partial^2 + |\lambda|^{1/2} \| u\|_\partial^2 \, \Big\}  + \varepsilon\, \bigg\{ \, \|\nabla u\|_\partial^2 + \|\phi\|_\partial^2 + \frac{1}{4}\,\Big\|\frac{\partial u}{\partial \nu} \Big\|_\partial^2 \, \bigg\}.
  \end{align*}
  The claim now follows if we use the definition of the conormal derivative and the norm equivalence on finite dimensional vector spaces.
\end{proof}

We prove one last lemma before we tackle the central theorem of this chapter.
The following lemma will not depend on the lemmata which were proven in the preceding part of this chapter, as the approach to derive the desired boundary estimates will be different:
We will not rely upon the rellich identities from Lemma \ref{lem:rellichIdentity} but directly part from a variational formulation of the Stokes resolvent problem on the boundary. 
Furthermore, we will work with the following Theorem about the regularity problem and the Neumann problem for the Laplacian on bounded Lipschitz domains
\begin{thm}
  \label{thm:jerisonKenig}
  Let $\Omega \subset \R^d$, $d \geq 2$, be a bounded Lipschitz domain. Then the following statements hold:
  \begin{enumerate}[a)]
    \item Given $f \in \HH^1(\partial\Omega)$, there exists a unique $\psi \in$ with $(\nabla \psi)^* \in \Ell^2(\partial\Omega)$ such that $\Delta u = 0$ in $\Omega$ and $\psi$ converges nontangentially to $f$ a.e.. Furthermore, the estimate 
      \begin{align*}
          \|\nabla\psi\|^{}_\partial \leq C\, \|f\|_{\HH^1(\partial\Omega)}
      \end{align*}
      holds with a constant only depending on the Lipschitz character of $\Omega$.
    \item Given $f \in \Ell^2(\partial\Omega)$ with $\int_{\partial\Omega} f \d\sigma = 0$, there exists a harmonic function $\psi$ on $\Omega$ with $\frac{\partial\psi}{\partial n} = f$ a.e.. Furthermore, the estimate
      \begin{align*}
          \| \psi\|_{\HH^1(\partial\Omega)} \leq C\, \|f\|^{}_\partial
      \end{align*}
      holds with C only depending on the Lipschitz character of $\Omega$.
  \end{enumerate}
\end{thm}

\begin{proof}
  According to Kenig, for the differential operator $\Delta$ the \emph{regularity problem}, $(R)_2$, and the \emph{Neumann problem}, $(N)_2$,  are solvable for data $f \in \Ell^2(\partial\Omega)$, see \cite[Thm. 2.1.10]{kenigBook}.
  Checking the definitions of $(R)_2$, see \cite[Defn. 1.7.10]{kenigBook}, and $(N)_2$, see \cite[Defn. 1.7.9]{kenigBook}, the claimed properties of follow from Theorems 1.8.2 and 1.8.3 in \cite[Chap. 1]{kenigBook}.
  We also refer to the works of Jerison and Kenig on the $\Ell^2$-regularity problem \cite{jerisonKenig2} and the Neumann problem \cite{jerisonKenig}.
\end{proof}

\begin{lem}
  Assume that $(u,\phi)$ satisfies the same conditions as in Theorem \ref{thm:rellich}.
  Then,
  \begin{align}
    \label{eq:phiDashintPhi}
    \bigg\|\phi - \bigg\{ \frac{1}{{r_0}^{d - 1}}\int_{\partial\Omega} \phi \d\sigma \bigg\} \bigg\|^{}_\partial 
    \leq C \Big\{ \|\nabla u\|^{}_\partial + |\lambda| \, \|u\cdot n\|_{\HH^{-1}(\partial\Omega)} \Big\}
  \end{align}
  and
  \begin{align}
    \label{eq:lambdaun}
    |\lambda|\, \| u\cdot n \|_{\HH^{-1}(\partial\Omega)} \leq C \Big\{ \|\phi\|^{}_\partial + \|\nabla u\|^{}_\partial \Big\},
  \end{align}
  where $C$ depends only on $d$ and the Lipschitz character of $\Omega$.
\end{lem}

\begin{proof}
  Our first goal will be to show that without loss of generality we may assume that $\Delta u = \nabla \phi + \lambda u$ on $\partial\Omega$.
  The central player will once again be Theorem \ref{thm:smoothApproximation}.
  To this end let $(\Omega_j)$ be a sequence of approximating $\CC^\infty$ domains as in Theorem \ref{thm:smoothApproximation} and suppose that
  \ref{eq:phiDashintPhi} holds for all $\Omega_j$, i.e. we have proven the inequality
  \begin{align}
    \label{eq:phiDashintPhij}
    \bigg\|\, \phi - \bigg\{ \frac{1}{{r_0}^{d - 1}}\int_{\partial\Omega} \phi \d\sigma \bigg\} \, \bigg\|_{\Ell^2(\partial\Omega_j; \C^d)}
     \leq C \Big\{ \|\nabla u\|_{\Ell^2(\partial\Omega_j; \C^{d\times d})} + |\lambda| \, \|u\cdot n\|_{\HH^{-1}(\partial\Omega_j)} \Big\}
  \end{align}
  and an inequality analogous to \eqref{eq:lambdaun}.
  Note that the mean value integral of $\phi$ should remain unchanged.
  It is now crucial that the constants in these inequalities only depend on the Lipschitz character of $\Omega$ and not on other geometric properties of the domain.
  Thus $C$ does not depend on $j$ and we may take the limit $j \to \infty$ using dominated convergence.

  For the rest of the proof we will thus assume that $(u,\phi)$ satisfies the Stokes resolvent problem in a domain $\Omega'$ for some $\overline \Omega \subseteq \Omega'$.
  In particular we have $\Delta u = \nabla \phi + \lambda u$  on $\partial\Omega$.
  Multiplying this identity  on $\partial\Omega$ with the outer normal vector $n$ and using the triangle inequality gives the following set of estimates:
  \begin{align}
    \label{eq:stokesEquationH1}
    \begin{alignedat}{1}
    \|\nabla \phi \cdot n\|_{\HH^{-1}(\partial\Omega)} 
    &\leq \|\Delta u \cdot n \|_{\HH^{-1}(\partial\Omega)} + |\lambda|\, \| u\cdot n\|_{\HH^{-1}(\partial\Omega)}, \\
    |\lambda|\, \| u\cdot n \|_{\HH^{-1}(\partial\Omega)} 
    &\leq \|\Delta u \cdot n \|_{\HH^{-1}(\partial\Omega)} + \|\nabla \phi \cdot n \|_{\HH^{-1}(\partial\Omega)}.
    \end{alignedat}
  \end{align}
  This looks almost like the desired pair of inequalities.
  We will now show that on the one hand the estimate
  \begin{align}
    \label{eq:deltaun}
    \|\Delta u \cdot n\|_{\HH^{-1}(\partial\Omega)}
    \leq C\, \|\nabla u\|^{}_\partial
  \end{align}
  and on the other hand the estimate
  \begin{align}
    \label{eq:nablaPhin}
    c \, \bigg\|\,\phi - \bigg\{ \frac{1}{{r_0}^{d - 1}}\int_{\partial\Omega} \phi \d\sigma \bigg\} \, \bigg\|^{}_\partial 
    %c \| \phi - \dashint_{\partial\Omega} \phi \d\sigma \|_\partial
    \leq \|\nabla \phi \cdot n\|_{\HH^{-1}(\partial\Omega)}
    \leq C\,  \|\phi\|^{}_\partial
  \end{align}
  holds.
  Using these two estimates applied to the respecive terms of \eqref{eq:stokesEquationH1}, we can directly verify \eqref{eq:phiDashintPhi} and \eqref{eq:lambdaun}.

  In order to prove \eqref{eq:deltaun}, we note the following identity of differential operators
  \begin{align*}
    \Delta u \cdot n = n_i \frac{\partial^2 u_i}{\partial x_j^2} = \Big\{ n_i \frac{\partial}{\partial x_j} - n_j \frac{\partial}{\partial x_i} \Big\} \frac{\partial u_i}{\partial x_j},
  \end{align*}
  where we used the fact that $\div u = 0$ in $\overline \Omega$.
  The expression between the brackets is the first-order tangential derivative $\partial_{\tau_{ij}}$, see \eqref{eq:defnTangDerivative}. 
  We can thus see that
  \begin{align*}
    \big| \langle \Delta u \cdot n, u \rangle \big| = \big| \langle \nabla u, \nabla_{\mathrm{tan}} u \rangle\big| \leq 2\, \|\nabla u\|_\partial^2\,,
  \end{align*}
  where we compared the tangential gradient with the usual gradient via the triangle inequality and applied the Cauchy-Schwartz estimate.  
  Identifying $\Delta u \cdot n$ with an element of $\HH^{-1}(\partial\Omega)$, the estimate \eqref{eq:deltaun} follows. 

  For the proof of estimate \eqref{eq:nablaPhin}, we will use $\Ell^2$-estimates for the Neumann and regularity problems for the Laplace equation in Lipschitz domains.
  Jerison and Kenig showed that for $g \in \Ell^2(\partial\Omega)$ with mean value zero the \emph{Neumann problem} for Laplace's equation on the Lipschitz domain $\Omega$ has a unique solution $\psi$ with  $(\nabla \psi)^* \in \Ell^2(\partial\Omega)$, $\frac{\partial \psi}{\partial n} = g$ on $\partial \Omega$ nontangentially and the solution fulfills the estimate $\|\psi\|_{\HH^1(\partial\Omega)} \leq C \|g\|^{}_\partial$, see Theorem \ref{thm:jerisonKenig}.
  By Proposition \ref{prop:approximationArgument}, we see that
  \begin{align*}
    \int_{\partial\Omega} \phi \, \frac{\partial\psi}{\partial x_i} n_i \d\sigma  
    = \int_\Omega \frac{\partial}{\partial x_i} \bigg\{ \phi \frac{\partial\psi}{\partial x_i} \bigg\} d x 
    = \int_\Omega \frac{\partial \phi}{\partial x_i} \cdot \frac{\partial \psi}{\partial x_i} \d x 
    = \int_{\partial\Omega} \psi \, \frac{\partial \phi}{\partial x_i} n_i \d \sigma.
  \end{align*}
  We can then use this identity and the estimate of $\psi$ against the data $g$ to derive
  \begin{align}
    \label{eq:dualityPhi}
    \bigg| \int_{\partial\Omega} \phi \, g \d \sigma\, \bigg|
    &\leq \Big\|\, \frac{\partial \phi}{\partial n}\, \Big\|_{\HH^{-1}(\partial\Omega)} \| \psi \|_{\HH^1(\partial\Omega)} 
    \leq C\, \Big\|\, \frac{\partial \phi}{\partial n} \,\Big\|_{\HH^{-1}(\partial\Omega)} \| g\|^{}_\partial.
  \end{align}
  
  Now, if we set $\overline g = \phi - \tilde \phi$, with $\tilde \phi \coloneqq |\partial \Omega|^{-1} \int_{\partial\Omega} \phi \d \sigma$, we arrive at the following estimate:
  \begin{align*}
      \| \phi - \tilde\phi \|_\partial^2
    = \int_{\partial\Omega} (\phi - \tilde \phi) \, \overline{(\phi - \tilde\phi)} \d\sigma 
    = \int_{\partial\Omega} \phi\, \overline{(\phi - \tilde\phi)} \d\sigma 
    \leq C\, \Big\|\, \frac{\partial\phi}{\partial n}\, \Big\|_{\HH^{-1}(\partial\Omega)} \|\phi - \tilde \phi\|^{}_\partial,
  \end{align*}
  where in the last step we used \eqref{eq:dualityPhi}.
  This together with Lemma \ref{lem:compareBoundaryWithBall} proves the left side of inequality \eqref{eq:nablaPhin}.

  For the right side of inequality \eqref{eq:nablaPhin}, we work in a similar way.
  We will use results for the \emph{regularity problem} of Laplace's equation by Jerison and Kenig, see Theorem \ref{thm:jerisonKenig}:
  Given $f \in \HH^1(\partial\Omega)$, there exists a harmonic function $\psi$ in $\Omega$ such that $(\nabla\psi)^* \in \Ell^2(\partial\Omega)$ and $\psi = f$ on $\partial\Omega$ nontangentially. Furthermore, the estimate $\|\nabla \psi\|^{}_\partial \leq C\, \|f\|_{\HH^1(\partial\Omega)}$ holds.
  As for \eqref{eq:dualityPhi}, we calculate
  \begin{align*}
    \bigg| \int_{\partial\Omega} \frac{\partial\phi}{\partial n} f \d\sigma\, \bigg|
    = \bigg| \int_{\partial\Omega} \phi\, \frac{\partial \psi}{\partial n} \d \sigma\, \bigg| 
    \leq \|\phi\|^{}_{\partial} \|\nabla \psi \|^{}_\partial 
    \leq C\, \|\phi\|^{}_\partial \|f\|_{\HH^1(\partial\Omega)},
  \end{align*}
  Interpreting elements from $\Ell^2(\partial\Omega)$ as functionals on $\HH^1(\partial\Omega)$, we obtain
  \begin{align*}
      \Big\|\,\frac{\partial\phi}{\partial n}\, \Big\|_{\HH^{-1}(\partial\Omega)} \leq C\, \|\phi\|^{}_\partial. 
  \end{align*}
  This finishes the proof of inequality \eqref{eq:nablaPhin}.
\end{proof}

\begin{rem}
  \label{rem:harmonicEstimate}
  A careful look at the proof of inequality \eqref{eq:nablaPhin} reveals that the estimate
  \begin{align*}
      c\, \| \phi \|^{}_\partial \leq \| \nabla \phi \cdot n \|_{\HH^{-1}(\partial\Omega)},
  \end{align*}
  holds for \emph{all} harmonic functions $\phi$ with vanishing mean on $\partial\Omega$ and not only those that correspond to the pressure term of the Stokes resolvent problem.
\end{rem}

After all this preparation we have accquainted enough tools and are now able to prove Theorem \ref{thm:rellich}.

\begin{proof}[Proof of Theorem \ref{thm:rellich}]
  For the proof of estimate \eqref{eq:rellich1}, we can assume without loss of generality that $\int_{\partial\Omega} \phi \d \sigma = 0$.

  We start by proving estimate \eqref{eq:rellich1}. 
  Using \eqref{eq:phiDashintPhi} to bound the second summand in \eqref{eq:rellich1} and then \eqref{eq:gradEstimateRellich2} for $\nabla u$, we get
  \begin{align*}
      \|\nabla u\|^{}_\partial + \|\phi\|^{}_\partial
      &\leq C\, \Big\{ \|\nabla u\|^{}_\partial + |\lambda| \|u \cdot n \|_{\HH^1(\partial\Omega)} \Big\} \\[0.5em]
      &\leq C_\varepsilon\, \Big\{ \big\|\nabla_{\mathrm{tan}} u \big\|^{}_\partial + |\lambda|^{1/2} \|u\|^{}_\partial + |\lambda| \| u\cdot n\|_{\HH^{-1}(\partial\Omega)} \Big\} 
      + C \,\varepsilon \,\Big\{ \|\nabla u\|^{}_\partial + \|\phi\|^{}_\partial \Big\}
  \end{align*}
  for all $\varepsilon \in (0,1)$.
  Chosing $\varepsilon$ such that $C\, \varepsilon < (1/2)$, we can rearrange the above inequality and obtain estimate \eqref{eq:rellich1}.

  Estimate \eqref{eq:rellich2} will need more effort to be proven.
  In order to obtain the desired estimate, we will divide the left hand side of inequality \eqref{eq:rellich2} into two groups as the following line suggests and then bound both groups separately.
  \begin{align*}
      \Big\{ \|\nabla u\|^{}_\partial + |\lambda|\, \|u \cdot n\|_{\HH^{-1}(\partial\Omega)} + \|\phi\|^{}_\partial \Big\} + |\lambda|^{1/2} \|u\|^{}_\partial \eqqcolon G_1 + G_2.
  \end{align*}
  We start with inequality \eqref{eq:lambdaun} and derive
  \begin{align*}
    G_1
    \leq C\, \Big\{ \|\nabla u\|^{}_\partial + \| \phi\|^{}_\partial \Big\}
    \leq C \, \Big\{\, \Big\|\,\frac{\partial u}{\partial \nu}\, \Big\|^{}_\partial + \|\nabla u\|^{}_\partial \Big\},
  \end{align*}
  where in the last step we used the definition of conormal derivatives.
  If we now apply \eqref{eq:gradEstimateRellich}, we get
  \begin{align*}
    G_1
    \leq C_\varepsilon\, \Big\| \, \frac{\partial u}{\partial \nu} \, \Big\|^{}_\partial + \varepsilon \, \Big\{ \|\nabla u\|^{}_\partial + \|\phi\|^{}_\partial + \| |\lambda|^{1/2} u \|^{}_\partial \Big\}
  \end{align*}
  for all $\varepsilon \in (0,1)$.
  Choosing $\varepsilon$ appropriately yields
  \begin{align}
    \label{eq:partOfRellich2}
    \|\nabla u\|^{}_\partial + \|\phi\|^{}_\partial + |\lambda| \, \|u \cdot n\|_{\HH^{-1}(\partial\Omega)}
    \leq C\, \Big\{ \, \Big\|\,\frac{\partial u}{\partial \nu} \, \Big\|^{}_\partial + |\lambda|^{1/2} \| u\|^{}_\partial \Big\},
  \end{align}
  and the first group has been successfully bounded.

  Now we need to estimate $G_2$.
  For this, we will work with the next identity which is a consequence of Green's theorem:
  \begin{align}
    \int_{\partial\Omega} h_k n_k |u|^2 \d \sigma
    %= \int_{\Omega} \frac{\partial}{\partial x_k} \big( h_k |u|^2 ) \d x
    &= \int_{\Omega} \frac{\partial h_k}{\partial x_k}\, |u|^2 \d x + \int_{\Omega} h_k \frac{\partial |u|^2}{\partial x_k}  \d x \nonumber\\[0.5em]
    \label{eq:hknkgreen}
    &= \int_\Omega \div(h)\, |u|^2 \d x + 2 \Re \int_\Omega h_k \frac{\partial \bar u_i}{\partial x_k} u_i \d x,
  \end{align}
  where $h \in \CC_0^1(\R^d, \R^d)$.
  This calculation is valid since Remark \ref{rem:shenNontangential} assures the existence of nontangential limits of $u$ and furthermore gives $(u)^* \in \Ell^2(\partial\Omega)$. Thus, $( h_k n_k |u|^2 )$ can be dominated by the integrable function $(\|h\|_\infty |(u)^*|^2)$ and the claim follows from Proposition \ref{prop:approximationArgument}.

  Next, in estimate \eqref{eq:hknkgreen}, we choose $h \in \CC_0^1(\R^d, \R^d)$ with $h_k n_k \geq c > 0$ on $\partial\Omega$, which is possible due to Theorem \ref{thm:smoothApproximation} and leads us to the estimate 
  \begin{align}
    \label{eq:estupartial}
    \|u\|_\partial^2 \leq C\, \Big\{ \int_\Omega |u|^2 \d x +  \int_\Omega |u| |\nabla u| \d x\, \Big\}.
  \end{align}
  Multiplying  \eqref{eq:estupartial} with $|\lambda|$ and using \eqref{eq:laxMilgramIneq}, we obtain 
  \begin{align*}
    |\lambda|\, \|u\|_\partial^2 
    &\leq C\, \Big\{\, |\lambda| \int_\Omega |u|^2 \d x + |\lambda|^{1/2}  \int_\Omega (|\lambda|^{1/2} |u|) |\nabla u| \d x\, \Big\} \nonumber\\[0.5em]
    %&\leq C \| \frac{\partial u}{\partial \nu} \|_\partial \|u\|_\partial + |\lambda|^{1/2} C \int_\Omega (|\lambda|^{1/2} |u|) |\nabla u| \d x \\
    &\leq C\, \Big\{ \, \Big\|\, \frac{\partial u}{\partial \nu} \, \Big\|^{}_\partial \|u\|^{}_\partial + |\lambda|^{1/2} \bigg( \int_\Omega |\lambda| |u|^2 \bigg)^{1/2} \bigg( \int_\Omega (|\nabla u|^2 \d x \bigg)^{1/2} \Big\} \\[0.5em]
    &\leq C\, \Big\|\,\frac{\partial u}{\partial \nu} \, \Big\|^{}_\partial \, |\lambda|^{1/2}  \|u\|^{}_\partial.
  \end{align*}
  Note that for the last estimate we also used the fact that $|\lambda| \geq \tau$ helps us to bound $\|u\|^{}_\partial$ by $C |\lambda|^{1/2} \|u\|^{}_\partial$.
  Rearranging terms in the last estimate, we now derive
  \begin{align}
    \label{eq:lambda12u}
    |\lambda|^{1/2} \|  u\|^{}_\partial \leq C\, \Big\|\, \frac{\partial u}{\partial \nu} \,\Big\|^{}_\partial.
  \end{align}
  Estimate \eqref{eq:rellich2} follows directly from \eqref{eq:partOfRellich2} in combination with \eqref{eq:lambda12u} and this concludes our proof.
\end{proof}

Shen proved that under reasonable assumptions a theorem similar to Theorem~\ref{thm:rellich} also holds for exterior domains, see \cite[Thm. 4.6]{Shen2012}.
%It is important to note that in the case $d = 2$ solutions $u$ that are given as a single layer potential do not fulfill the stated requirements on the decay.

\begin{thm}
  \label{thm:rellichExterior}
  Let $\lambda \in \Sigma_\theta$ and $|\lambda| \geq \tau$, where $\tau \in (0,1)$.
  Let $(u,\phi)$ be a solution of the Stokes resolvent problem in $\Omega_- = \R^d \setminus \overline\Omega$.
  Suppose additionally that $(\nabla u)^*$, $(\phi)^* \in \Ell^2(\partial\Omega)$ and that $\nabla u$, $\phi$ have nontangential limits almost everywhere on $\partial\Omega$.
  Furthermore, let for $|x| \to \infty$
  \begin{align*}
    |\phi(x)| + |\nabla u(x)| = O(|x|^{1 - d}) \quad\text{and}\quad 
    u(x) = \begin{cases} O(|x|^{2 - d}) &\quad\text{if } d \geq 3, \\ o(1) &\quad\text{if } d = 2. \end{cases}
  \end{align*}
  Then the estimates
  \begin{align}
    \label{eq:rellich1ext}
    \|\nabla u\|^{}_\partial + \|\phi\|^{}_\partial
    \leq C\, \Big\{ \|\nabla_{\mathrm{tan}} u\|^{}_\partial + |\lambda|^{1/2} \|u\|^{}_\partial + |\lambda|\, \|u \cdot n\|_{\HH^{-1}(\partial\Omega)} \Big\}
  \end{align}
  and
  \begin{align}
    \label{eq:rellich2ext}
    \|\nabla u\|^{}_\partial + |\lambda|^{1/2} \|u\|^{}_\partial + |\lambda| \, \|u \cdot n\|_{\HH^{-1}(\partial\Omega)} + \|\phi\|^{}_\partial
    \leq C \, \Big\|\,\frac{\partial u}{\partial \nu} \,\Big\|^{}_\partial
  \end{align}
  hold, where $C$ depends only on $d$, $\tau$, $\theta$ and the Lipschitz character of $\Omega$.
\end{thm}

\chapter{Solving the $\Ell^2$ Dirichlet Problem}
\label{chap:5}

This section is all about the application of the method of layer potentials to solve the $\Ell^2$ Dirichlet problem \hyperref[eq:dirProblem]{$(\mathrm{Dir}_\lambda)$} for the Stokes resolvent system.
The results from Chapter~\ref{chap:3} on the jump relations and the corresponding operators, see Theorem~\ref{thm:jumpConditions} and Theorem~\ref{thm:nontangentialLimitDoubleLayer}, together with the Rellich estimates from Chapter~\ref{chap:4} will be the essential ingredients.

In the first part of this chapter, we will show that the operators $(1/2)I + \K_\lambda$ and their adjoints are isomorphisms on suitable Hilbert spaces.
Then, in the second part of this chapter, we will build on this information to construct solutions to $(\mathrm{Dir}_\lambda)$ via double layer potentials.
%Furthermore we will establish a uniform $\Ell^p$ estimate for the nontangential maximal function which will be important for the proof of our central theorem.

For the remainder of this chapter, let $\Omega$ always denote a bounded Lipschitz domain in $\R^d$, $d \geq 2$, with connected boundary.
We will use $\Ell^2_n(\partial\Omega)$ to denote the function space
\begin{align*}
  \Ell^2_n(\partial\Omega) \coloneqq \Big\{ f \in \Ell^2(\partial\Omega; \C^d) \colon \int_{\partial\Omega} f \cdot n \d \sigma = 0 \Big\},
\end{align*}
and $\Ell_0^2(\partial\Omega; \C^d)$  to denote the function space of $\Ell^2$ functions with mean value zero.
As before, $\| \cdot \|^{}_\partial$ stands for the norm of $\Ell^2(\partial\Omega; \C^k), k \in \N$.

The following results will build on the application of Theorem \ref{thm:rellichExterior}.
Therefore, the next lemma shows that solutions to the Stokes resolvent problem that are given by single layer potentials fulfill the requirements of the theorem.

\begin{lem}
  \label{lem:requirements}
  Let $\lambda \in \Sigma_\theta$ and $(u,\phi)$ be given by \eqref{eq:defSingleLayer} and \eqref{eq:defSingleLayerPressure}, respectively.
  Then the following holds for $|x| \to \infty\,$:
  \begin{align*}
    |\phi(x)| + |\nabla u(x)| = O(|x|^{1 - d}) \quad\text{and}\quad 
    u(x) = \begin{cases} O(|x|^{2 - d}) &\quad\text{if } d \geq 3, \\ o(1) &\quad\text{if } d = 2. \end{cases}
  \end{align*}
\end{lem}

\begin{proof}
  If $d \geq 2$, then an application of the dominated convergence theorem gives that $|\phi(x)| + |\nabla u(x)| = O(|x|^{1 - d})$ as $|\phi(x)| = O(|\Phi_k(x)|) = O(|x|^{1 - d})$ by \eqref{eq:fundamentalVectorPressure}.
 Furthermore, we have $|\nabla u(x)| = O(|x|^{1 - d})$ by estimate \eqref{eq:fundamentalMatrixEstimate}.

  If $d \geq 3$, then the first part of Lemma \ref{lem:estimateHelmholtzDerivatives} and an application of the dominated convergence theorem gives $u(x) = O(|x|^{2 - d})$.
  If $d = 2$, consider the definition of the fundamental matrix \eqref{eq:fundamentalMatrixStokes}. According to the first part of Lemma \ref{lem:estimateHelmholtzDerivatives} and the asymptotic behavior of the second derivative of $\log(|x|)$, we only have to worry about the first summand in \eqref{eq:fundamentalMatrixStokes}. 
  But the asymptotic behavior of the fundamental solution to the scalar Helmholtz equation is already available thanks to the second part of Lemma \ref{lem:estimateHelmholtzDerivatives}. 
  Through dominated convergence, the same asymptotic behavior holds for $u(x)$.
\end{proof}

We will now derive bounds on the inverse of the operator $(1/2) I + \K_\lambda$ from Chapter~\ref{chap:3}. 

\begin{lem}
  Let $\lambda \in \Sigma_\theta$ and $|\lambda| \geq \tau$, where $\tau \in (0,1)$.
  Suppose that $|\partial\Omega| = 1$.
  Then $(1/2)I + \K_\lambda$ is an isomorphism on $\Ell^2(\partial\Omega; \C^d)$ and
  \begin{align}
    \label{eq:inverseEstimate}
    \| f\|^{}_\partial \leq C \, \big\|\, \big( (1/2) I + \K_\lambda \big) f\, \, \big\|^{}_\partial \quad\text{for any } f \in \Ell^2(\partial\Omega; \C^d),
  \end{align}
  where $C$ depends only on $d$, $\theta$, $\tau$ and the Lipschitz character of $\Omega$.
\end{lem}

\begin{proof}
  We start with $f \in \Ell^2(\partial\Omega; \C^d)$ and the corresponding single layer potentials $u = \slp_\lambda(f)$ and $\phi = \slp_\Phi(f)$ given by \eqref{eq:defSingleLayer} and \eqref{eq:defSingleLayerPressure}.
  We saw in Chapter 3 that $(u,\phi)$ solves the Stokes resolvent problem in $\R^d \setminus \partial\Omega$ and got from Lemma \ref{lem:nontangentialMaximalFunctions} with $p = 2$ for the nontangential maximal functions that $(\nabla u)^*$, $(\phi)^* \in \Ell^2(\partial\Omega)$.
  We furthermore saw in Lemma \ref{lem:traceFormulas} that $\nabla u$ and $\phi$ have nontangential limits almost everywhere on $\partial\Omega$. 
  Finally, in Theorem \ref{thm:jumpConditions} we saw that $\nabla_{\mathrm{tan}} u_+ = \nabla_{\mathrm{tan}} u_-$ and derived the jump relation $\big(\frac{\partial u}{\partial \nu}\big)_\pm = (\pm (1/2) I + \K_\lambda) f$.

  Let us assume for a moment that the estimate
  \begin{align}
    \label{eq:negNablaPhi}
    \| \nabla u_- \|^{}_\partial + \| \phi_- \|^{}_\partial \leq C \, \Big\| \bigg( \frac{\partial u}{\partial \nu} \bigg)_+ \Big\|^{}_\partial .
  \end{align}
  holds for a constant $C$ depending only on $d$, $\theta$, $\tau$ and the Lipschitz character of $\Omega$. 
  Using \eqref{eq:negNablaPhi}, we can prove \eqref{eq:inverseEstimate}:
  Note that the jump relation gives us $f = \big(\frac{\partial u}{\partial \nu} \big)_+ -  \big( \frac{\partial u}{\partial \nu} \big)_-$.
  With the definition of the conormal derivative and estimate \eqref{eq:negNablaPhi}, we calculate that
  \begin{align*}
    \| f\|^{}_\partial 
    &\leq  \Big\|    \bigg( \frac{\partial u}{\partial \nu} \bigg)_+ \Big\|^{}_\partial + \Big\|  \bigg( \frac{\partial u}{\partial \nu} \bigg)_- \Big\|^{}_\partial  \\[0.5em]
    &\leq \Big\|     \bigg( \frac{\partial u}{\partial \nu} \bigg)_+  \Big\|^{}_\partial + \Big\| \bigg( \frac{\partial u}{\partial n}   \bigg)_- \Big\|^{}_\partial + \| \phi_- \|^{}_\partial  \\[0.5em]
    &\leq C\, \Big\| \bigg( \frac{\partial u}{\partial \nu} \bigg)_+ \Big\|^{}_\partial
    = C\, \big\| \big((1/2)I + K_\lambda\big) f\, \big\|^{}_\partial.
  \end{align*}

  In order to prove \eqref{eq:negNablaPhi}, note that due to Lemma \ref{lem:requirements} we can use Theorem \ref{thm:rellichExterior} to derive
  \begin{align}
    \|\nabla u_-\|^{}_\partial + \|\phi_-\|^{}_\partial
    &\leq C\, \Big\{ \big\|\nabla_{\mathrm{tan}} u_- \big\|^{}_\partial + |\lambda|^{1/2} \| u_-\|^{}_\partial + |\lambda|\, \| n \cdot u_-\|_{\HH^{-1}(\partial\Omega)} \Big\} \nonumber\\[0.5em]
    \label{eq:nablauminus}
    &= C \, \Big\{ \big\|\nabla_{\mathrm{tan}} u_+ \big\|^{}_\partial + |\lambda|^{1/2} \| u_+ \|^{}_\partial + |\lambda|\, \|n \cdot u_+ \|_{\HH^{-1}(\partial\Omega)} \Big\},
  \end{align}
  where we used the fact that $u_+ = u_-$ and $\nabla_{\mathrm{tan}} u_+ = \nabla_{\mathrm{tan}} u_-$ on $\partial\Omega$. The former fact is a consequence of the continuity of the single layer potential across $\partial\Omega$, see Mitrea and Wright \cite[Prop. 4.7]{mitreaWright}.
  Inequality \eqref{eq:rellich2} of Theorem \ref{thm:rellich} now allows us to estimate the right hand side of \eqref{eq:nablauminus} by $C \,\|\big(\frac{\partial u}{\partial \nu} \big)_+ \|^{}_\partial$ and thus the desired estimate \eqref{eq:negNablaPhi} follows.

  Let's now work on the invertibility of $(1/2) I + \K_\lambda$. In the case $\lambda = 0$, Mitrea and Wright showed in \cite[Eq. (5.166)]{mitreaWright} that $(1/2) I + \K_0$ as an operator on $\Ell^2(\partial\Omega; \C^d)$ has a one dimensional null space and as range the space $\Ell_0^2(\partial\Omega; \C^d)$.
  Thus, $(1/2) I + \K_0$ has Fredholm index $0$.
  %This remains true if we replace $\Ell^2(\partial\Omega; \R^d)$ by $\Ell^2(\partial\Omega; \C^d)$ as this just corresponds to a complexification of the vector space and the operator. 
  Since the operator $\K_\lambda - \K_0$ is compact on $\Ell^2(\partial\Omega; \C^d)$ by Lemma \ref{lem:compactness}, we deduce that the operator
  \begin{align*}
    (1/2)I + \K_\lambda = (1/2)I + \K_0 + (\K_\lambda - \K_0)
  \end{align*}
  has the Fredholm index zero as well for all $\lambda \in \Sigma_\theta$.
  Now inequality \eqref{eq:inverseEstimate} gives that $(1/2)I + \K_\lambda$ is injective and thus the Fredholm index of zero implies that it is also surjective and hence an isomorphism.
\end{proof}

The next lemma works with the counterpart of $(1/2)I + \K_\lambda$.

\begin{lem}
  \label{lem:inverseEstimate}
  Let $\lambda \in \Sigma_\theta$.
  Then $-(1/2)I + \K_\lambda$ is a Fredholm operator on $\Ell^2(\partial\Omega; \C^d)$ with index zero and
  \begin{align}
    \label{eq:inverseEstimate2}
    \| f\|^{}_\partial \leq C \, \big\| \, \big( - (1/2) I + \K_\lambda \big) f\, \big\|^{}_\partial \quad\text{for all } f \in \Ell^2_n(\partial\Omega),
  \end{align}
  where $C$ depends only on $d$, $\theta$ and the Lipschitz character of $\Omega$.
\end{lem}

\begin{proof}
  Let us assume without loss of generality that $|\partial\Omega| = 1$.
  In the case $\lambda = 0$, Mitrea and Wright showed in \cite[Eq. (5.166)]{mitreaWright} that the Fredholm index of the operator $-(1/2)I + \K_0$ on $\Ell^2(\partial\Omega; \C^d)$ is zero and estimate \eqref{eq:inverseEstimate2} holds.
  %As in the previous proof, this still remains true if we complexify the operator making it a Fredholm operator with index zero on  $\Ell^2(\partial\Omega; \C^d)$.
  Since $\K_\lambda - \K_0$ is compact on $\Ell^2(\partial\Omega; \C^d)$ and the Fredholm index remains unchanged under compact perturbations, we know that the Fredholm index of $-(1/2)I + \K_\lambda$ on $\Ell^2(\partial\Omega; \C^d)$ is zero for all $\lambda \in \Sigma_\theta$.
  This proves the first claim of the lemma.

  Now let $\tau < (2 \diam(\Omega)^2 + 1)^{-1}$ and $|\lambda| < \tau$.
  We claim that
  \begin{align*}
  \big\|\, (\K_\lambda - \K_0) f\, \big\|^{}_\partial \leq C \, |\lambda|^{1/2} \|f\|^{}_\partial.
  \end{align*}
  In order to prove this inequality, we want to apply Young's inequality from Lemma \ref{lem:young}. To this end, we start by estimating
  \begin{align*}
      \big\|\,  (\K_\lambda - \K_0) f\, \big\|^{}_\partial 
    \leq \sup_{p \in \partial\Omega} \Big\| \nabla_x \Big\{ \Gamma(p - \cdot\,; \lambda) - \Gamma(p - \cdot\,; 0) \Big\} \Big\|_{\Ell^1(\partial\Omega)} \| f\|_{\Ell^2(\partial\Omega)}.
  \end{align*}
  In the next step, we prove that for $p \in \partial\Omega$ the integral over the gradients of $\Gamma$ can be estimated independent of $p$.
  This is straightforward using Lemma \ref{lem:central} as Corollary \ref{cor:differenceFundamentalSolutionStokes} gives us
  \begin{align*}
      &\int_{\partial\Omega} | \nabla_x \{ \Gamma(p - y; \lambda) - \Gamma(p - y; 0) \} | \d\sigma(y)\\[0.5em]
      &\quad\leq C \, |\lambda|^{1/2} \int_{\partial\Omega} \frac{1}{|p - y|^{d - 2}} \d \sigma(y) \\[0.5em]
      &\quad= C \, |\lambda|^{1/2} \, \int_{\partial\Omega \cap \BB(p,\, r_0/4)} \frac{1}{|p - y|^{d - 2}} \d\sigma(y) + C |\lambda|^{1/2}\, \int_{\partial\Omega \setminus \BB(p,\, r_0/4)} \frac{1}{|p - y|^{d - 2}} \d\sigma(y) \\[0.5em]
      &\quad\leq C \, |\lambda|^{1/2} \,\big( r_0/4 + 4^{2 - d} r_0^{d - 2} |\partial\Omega| \big),
  \end{align*}
  where $r_0$ is the radius from the definition of Lipschitz domains.
  But as $|\partial\Omega| = 1$ and $r_0$ can be related to $|\partial\Omega|$ by Lemma \ref{lem:compareBoundaryWithBall}, we get
  \begin{align*}
      \int_{\partial\Omega} |\nabla_x \{\Gamma(p - y; \lambda) - \Gamma(p - y; 0) \} | \d\sigma(y) \leq C \, |\lambda|^{1/2}
  \end{align*}
  with a constant $C$ that only depends on $d$, $\theta$ and the Lipschitz character of $\Omega$.
  Note that by the choice of $\tau$ the estimate from Corollary \ref{cor:differenceFundamentalSolutionStokes} applies on the whole domain of integration.

  For $f \in \Ell^2_n(\partial\Omega)$, we can now estimate
  \begin{align*}
    \| f\|^{}_\partial 
    &\leq C\,  \big\|\,(-(1/2)I + \K_0) f \, \big\|^{}_\partial  \\
    &\leq C\,  \big\|\,(-(1/2)I + \K_\lambda)f\, \big\|^{}_\partial + \big\|\, (\K_\lambda - \K_0)f\, \big\|^{}_\partial \\
    &\leq C\,  \big\|\,(-(1/2)I + \K_\lambda)f\, \big\|^{}_\partial + C \, |\lambda|^{1/2} \|f\|^{}_\partial,
  \end{align*}
  with a constant $C$ which depends only on $d$, $\theta$ and the Lipschitz character of $\partial\Omega$.
  Choosing $\tau$ even smaller allows us to rearrange the terms in the above estimate such that estimate \eqref{eq:inverseEstimate2} holds for $\lambda \in \Sigma_\theta$ and $|\lambda| < \tau$, with $\tau$ depending on $d$, $\theta$ and the Lipschitz character of $\Omega$.

  Now leave $\tau$ fixed and consider the case $|\lambda| \geq \tau$.
  This case will be handled using the Rellich estimates from Section 4.
  We use the facts that for $\nabla_{\mathrm{tan}} u$ and $u$ the inner and outer nontangential limits coincide and apply Theorems \ref{thm:rellich} and \ref{thm:rellichExterior} to conclude that
  \begin{align*}
      &\|\nabla u_+ \|_\partial + \Big\| \, \phi_+ - \frac{1}{r_0^{d - 1}}\int_{\partial\Omega} \phi_+ \d\sigma \, \Big\|^{}_\partial \\[0.5em]
      &\qquad\leq C\, \Big\{ \big\|( \nabla_{\mathrm{tan}} u)_+ \big\|_\partial^{} + |\lambda|^{1/2} \|u_+\|^{}_\partial + |\lambda|\,  \| u_+\cdot n\|_{\HH^{-1}(\partial\Omega)} \Big\} \\[0.5em]
      &\qquad = C\, \Big\{ \big\|( \nabla_{\mathrm{tan}} u)_- \big\|_\partial^{} + |\lambda|^{1/2} \|u_-\|^{}_\partial + |\lambda| \, \| u_-\cdot n\|_{\HH^{-1}(\partial\Omega)} \Big\}
      \leq C\, \Big\| \bigg(\frac{\partial u}{\partial \nu} \bigg)_- \Big\|^{}_\partial.
  \end{align*}
  We can now use this inequality to estimate $\| \big( \frac{\partial u}{\partial \nu} \big)_+ \|_\partial$ via
  \begin{align*}
      \Big\| \bigg( \frac{\partial u}{\partial \nu} \bigg)_+ \Big\|^{}_\partial
      &\leq \Big\|\bigg(\frac{\partial u}{\partial n} \bigg)_+ \Big\|^{}_\partial + C\, \| \phi_+ \|^{}_\partial \\[0.5em]
      &\leq C\,  \Big\{  \| (\nabla u)_+ \|^{}_\partial + \Big\|\, \phi_+ - \frac{1}{r_0^{d - 1}} \int_{\partial\Omega} \phi_+ \d\sigma \, \Big\|^{}_\partial + \Big| \, \frac{1}{r_0^{d - 1}} \int_{\partial\Omega} \phi_+ \d \sigma \, \Big| \; \Big\}\\[0.5em]
      &\leq C\,  \Big\{ \, \Big\| \bigg(\frac{\partial u}{\partial \nu} \bigg)_- \Big\|^{}_\partial + \Big| \int_{\partial\Omega} \phi_+ \d\sigma \, \Big|\,\Big\}.
  \end{align*}
  Considering the jump relation \eqref{eq:nontangentialConormalDerivative} and the previous estimate, we get that
  \begin{align}
      \|f \|^{}_\partial 
      &\leq \Big\|\Big( \frac{\partial u}{\partial \nu} \Big)_+ \Big\|^{}_\partial + \Big\| \Big( \frac{\partial u}{\partial \nu} \Big)_- \Big\|^{}_\partial  \nonumber\\[0.5em]
      &\leq C \,\Big\{ \,  \Big\| \Big( \frac{\partial u}{\partial \nu} \Big)_- \Big\|^{}_\partial + \Big| \int_{\partial\Omega} \phi_+ \d\sigma \, \Big|\, \Big\} \nonumber\\[0.5em]
    \label{eq:estimatef}
    &\leq C \, \Big\{ \big\| (-(1/2)I + \K_\lambda) f \big\|^{}_\partial +  \Big| \int_{\partial\Omega} \phi_+ \d\sigma \, \Big|\,\Big\}.
  \end{align}
  We now are left with the term $\int_{\partial\Omega} \phi_+ \d\sigma$ that needs to be estimated.
  To this end, note that multiplying the conormal derivatives of $u$ by $n$ gives
  \begin{align*}
    \bigg( \frac{\partial u}{\partial \nu} \bigg)_+ \cdot n
    = \bigg( \frac{\partial u_i}{\partial x_j}\bigg)_+ n_i n_j - \phi_+
    = n_j \Big( n_i \frac{\partial}{\partial x_j} - n_j \Big( \frac{\partial}{\partial x_i} \Big) u_i \Big)_+ - \phi_+,
  \end{align*}
  where for the second equality we used that $\div(u) = 0$ in $\Omega$ and thus this also holds for the nontangential limit.
  Using the above identity, it follows that
  \begin{align}
    \Big| \int_{\partial\Omega} \phi_+ \d\sigma \, \Big|
    &\leq \Big| \int_{\partial\Omega} \Big( \frac{\partial u}{\partial \nu} \Big)_+ \cdot n \d\sigma \, \Big|  + C\, \big\|\nabla_{\mathrm{tan}} u \big\|^{}_\partial \nonumber\\[0.5em]
    &\leq \Big| \int_{\partial\Omega} \Big( \frac{\partial u}{\partial \nu} \Big)_- \cdot n \d\sigma \, \Big| +  C\, \big\| \nabla_{\mathrm{tan}} u\big\|^{}_\partial \nonumber\\[0.5em]
    \label{eq:estimatephiplus}
    &\leq C\, \Big\| \Big( \frac{\partial u}{\partial \nu} \Big)_- \Big\|^{}_\partial,
  \end{align}
  where in the second step, we used the jump relation to exchange $\big(\frac{\partial u}{\partial \nu}\big)_+ \cdot n$ by $\big(\frac{\partial u}{\partial \nu} \big)_- + f \cdot n$ and then used the fact $f \in \Ell^2_n(\partial\Omega)$.
  The third step follows from Theorem \ref{thm:rellichExterior} considering that $\big\|\nabla_{\mathrm{tan}} u \big\|^{}_\partial \leq C\, \|\nabla u \|_{\partial}$.
  Now, extending estimate \eqref{eq:estimatef} by \eqref{eq:estimatephiplus} gives
  \begin{align*}
      \| f\|^{}_\partial 
      \leq C \, \big\| ( - (1/2) I + \K_\lambda) f \big\|^{}_\partial + C\, \Big\| \Big(\frac{\partial u}{\partial \nu} \Big)_- \Big\|^{}_\partial 
    \leq C \, \big\| \big( - (1/2) I + \K_\lambda) f\big\|_{\partial},
  \end{align*}
  where we used the jump relation \eqref{eq:nontangentialConormalDerivative} again.
  This proves estimate \eqref{eq:inverseEstimate2} in the case $|\lambda| \geq \tau$ and thus concludes the proof.
\end{proof}

With the following lemma that looks like a reverse trace theorem, we will later show the uniqueness of solutions to the $\Ell^2$ Dirichlet problem to the Stokes resolvent system.

\begin{lem}
  \label{lem:l2unique}
  Let $\lambda \in \Sigma_\theta$ and $(u,\phi)$ be a solution to the Stokes resolvent problem in $\Omega$.
  Furthermore, suppose that the nontangential limit of $u$ exists almost everywhere on $\partial\Omega$ and that $(u)^* \in \Ell^2(\partial\Omega)$.
  Then 
  \begin{align}
    \label{eq:OmegaBoundaryEstimate}
    \int_\Omega |u|^2 \d x \leq C \int_{\partial\Omega} |u|^2 \d\sigma,
  \end{align}
  where $C$ depends only on $d$, $\theta$ and the Lipschitz character of $\Omega$.
\end{lem}

\begin{proof}
  We use the approximation theorem \ref{thm:smoothApproximation} and approximate $\Omega$ by a sequence of smooth domains with uniform Lipschitz characters from inside.
  It suffices to prove \eqref{eq:OmegaBoundaryEstimate} for elements of this sequence of domains as $(u)^* \in \Ell^2(\partial\Omega)$.
  As a consequence, we will assume for the rest of the proof that $\Omega$ is smooth and that $u$, $\phi$ are smooth in $\overline\Omega$.
  %The fact that $(u)^* \in \Ell^2(\partial\Omega)$ is necessary if we want to take the limit of the approximating sequence of domains, as it enables us to apply the dominated convergence theorem.
  Let $(w,\psi) \in \HH_0^1(\Omega; \C^d) \times \HH^1(\Omega)$ be a solution to the inhomogenous system
  \begin{align}
    \label{eq:inhomogenousStokes}
    \begin{alignedat}{1}
      - \Delta w + \lambda w + \nabla \psi = \overline u &\text{ in } \Omega, \\
      \div(w) = 0 &\text{ in } \Omega.
    \end{alignedat}
  \end{align}
  In fact, the regularity theory for the Stokes equation gives us that $w$ and $\psi$ are even smooth in $\Omega$ as $u$ is smooth.
  It follows from testing \eqref{eq:inhomogenousStokes} against $u$ that
  \begin{align}
    \label{eq:testingInhomogenousStokes}
    \int_\Omega |u|^2 \d x
    &= \int_\Omega u \cdot \big\{ - \Delta w + \lambda w + \nabla \psi \big\} \d x.
    %&= - \int_{\partial\Omega} u \cdot \big\{ \frac{\partial w}{\partial n} - \psi n \big\} \d \sigma
  \end{align}
  The left hand side of \eqref{eq:testingInhomogenousStokes} gives the starting point for the proof of inequality \eqref{eq:OmegaBoundaryEstimate}.
  
  Using one of Green's identities, see \cite[Thm. 3, App. C.2]{evans}, on the first summand and the fact that $u$ is the solution to the Stokes resolvent problem gives that
  \begin{align*}
    \int_\Omega -u \cdot \Delta w \d x
    &= \int_\Omega -w \cdot \Delta u \d x - \int_{\partial\Omega} u \cdot \frac{\partial w}{\partial n} \d \sigma, \\[0.5em]
    &= \int_\Omega w \cdot ( -\lambda u - \nabla \phi) \d x - \int_{\partial\Omega} u \cdot \frac{\partial w}{\partial n} \d \sigma \\[0.5em]
    &= \int_\Omega -\lambda w \cdot u \d x - \int_{\partial\Omega} u \cdot \frac{\partial w}{\partial n} \d \sigma, 
  \end{align*}
  where in the last step we used integration by parts and the fact that $w$ vanishes on $\partial\Omega$ and is divergence free:
  \begin{align*}
    \int_\Omega w \cdot \nabla\phi \d x = -\int_\Omega \div(w) \, \phi \d x + \int_{\partial\Omega} \phi \, w \cdot n \d\sigma = 0.
  \end{align*}
  For the third summand in \eqref{eq:testingInhomogenousStokes}, we do the same with the only difference that the second integral does not vanish:
  \begin{align*}
    \int_\Omega u \cdot \nabla\psi \d x 
    = -\int_\Omega \div(u)\, \psi \d x + \int_{\partial\Omega} \psi \, u \cdot n \d\sigma 
    = \int_{\partial\Omega} \psi \, u \cdot n \d\sigma.
  \end{align*}
  Putting everything together gives
  \begin{align}
    \int_\Omega |u|^2 \d x 
    &= \Big| \int_{\partial\Omega} u \cdot \Big\{ \frac{\partial w}{\partial n} - \psi n \Big\} \d\sigma\, \Big|
    \leq \| u\|^{}_\partial \, \Big\{ \|\nabla w\|^{}_\partial + \|\psi\|^{}_\partial \Big\} \label{eq:u2estimate}
  \end{align}
  by the Cauchy-Schwartz inequality.
  As the pressure $\psi$ is only specified modulo additive constants, we may as well assume that $\int_{\partial\Omega} \psi \d\sigma = 0$.
  Furthermore, by the Schwartz theorem we see from \eqref{eq:inhomogenousStokes} that $\Delta \psi = \div(\overline u) = 0$ in $\Omega$. 
  As stated in Remark \ref{rem:harmonicEstimate}, this allows us to use the results from the proof of \eqref{eq:nablaPhin} by setting $\phi = \psi$ to conclude that
  \begin{align}
      \|\psi\|^{}_\partial 
    &\leq C\, \| \nabla \psi \cdot n\|_{\HH^{-1}(\partial\Omega)} \nonumber
    \intertext{and since $w$ has vanishing trace on $\partial\Omega$ we can use that property together with the fact that  $(w,\psi)$ solves \eqref{eq:inhomogenousStokes} to further estimate}
    &\leq C \, \Big\{ \| \Delta w \cdot n \|_{\HH^{-1}(\partial\Omega)} + \| u \cdot n \|_{\HH^{-1}(\partial\Omega)} \Big\} \nonumber\\[0.5em]
    &\leq C \, \Big\{ \|  \nabla w\|^{}_\partial + \|u\|^{}_\partial  \Big\}, \label{eq:psiEstimate}
  \end{align}
  where for the last estimate we used \eqref{eq:deltaun} which is applicable since $\div w = 0$ on $\Omega$.
  If we combine inequalities \eqref{eq:u2estimate} and \eqref{eq:psiEstimate}, we get
  \begin{align}
    \int_\Omega |u|^2 \d x \leq C \, \Big\{ \|u\|^{}_\partial \|\nabla w \|^{}_\partial + \|u\|_\partial^2\Big\}. \label{eq:u2estimate2}
  \end{align}
  We are left with estimating the first term in \eqref{eq:u2estimate2}.
  In fact, it will suffice to show the following inequality
  \begin{align}
    \label{eq:w2estimate}
    \int_{\partial\Omega} |\nabla w|^2 \d\sigma \leq C \int_\Omega |u|^2 \d x + C \int_{\partial\Omega} |u|^2 \d \sigma
  \end{align}
  since by the weighted Young inequality for real numbers this would make the estimate
  \begin{align*}
    C \| u\|^{}_\partial \|\nabla w\|^{}_\partial \leq \frac{1}{2} \int_\Omega |u|^2 \d x + C \int_{\partial\Omega} |u|^2 \d \sigma
  \end{align*}
  available which after rearranging terms yields \eqref{eq:OmegaBoundaryEstimate}.

  In order to derive estimate \eqref{eq:w2estimate}, we need the Rellich-type identity
%  \begin{align}
%    \label{eq:rellichIdentity3}
%    \begin{alignedat}{1}
%    \int_{\partial\Omega} h_k n_k |\nabla w|^2 \d \sigma
%    &= \int_\Omega \div(h)\, |\nabla w|^2 \d x + 2 \Re \int_{\Omega} h_j \frac{\partial \psi}{\partial x_j} \frac{\overline w_j}{\partial x_k}\d x   \\[0.5em]
%    &\quad + 2 \Re \int_\Omega h_k \lambda w_j \frac{\partial\overline w_j}{\partial x_k} + 2 \Re \int_\Omega h_k \overline u_j \frac{\partial \overline w_j}{\partial x_k} \d x. 
%    \end{alignedat}
%  \end{align} 
  \begin{samepage}
  \begin{align}
    \int_{\partial\Omega} h_k n_k |\nabla w|^2 \d\sigma
    &= - \int_\Omega \div(h)\, |\nabla w|^2 \d x 
    + 2 \Ret \int_\Omega \frac{\partial h_k}{\partial x_j} \cdot \frac{\partial w_i}{\partial x_k} \cdot \frac{\partial \overline w_i}{\partial x_j} \d x  \nonumber\\[0.5em]
    &\quad - 2 \Ret \int_\Omega \frac{\partial h_k}{\partial x_i} \cdot \frac{\partial w_i}{\partial x_k} \, \overline \psi \d x 
    + 2 \Ret \int_\Omega h_k \frac{\partial w_i}{\partial x_k} \cdot \overline{ \lambda w_i} \d x  \nonumber\\[0.5em]
    &\quad + 2 \Ret \int_{\Omega} h_k \frac{\partial \overline w_i}{\partial x_k} \, \overline u_i \d x\,, \label{eq:rellichIdentity3}
  \end{align}
  \end{samepage}
  where $h = (h_1, \dots,h_d) \in \CC_0^1(\R^d, \R^d)$.
  Note that since all involved quantities are smooth up to the boundary, integration by parts is allowed and the proof of the stated  Rellich identity boils down to a formal calculation.
  The proof is analogous to the proof of Rellich identity \eqref{eq:rellichIdentity2} and one can prove that for $(w,\psi)$ an equality like \eqref{eq:rellichIdentity2} holds with one extra term, namely
  \begin{align*}
      - 2\Re \int_\Omega h_k \frac{\partial \overline w_i}{\partial x_k} \, \overline u_i \d x.
  \end{align*}
  This term will show up in the proof of Lemma~\ref{lem:rellichIdentity} in the calculation of $I_5$ when one uses the fact that $(w,\psi)$ solves problem \eqref{eq:inhomogenousStokes}.
  Finally, the first two terms on the right hand side of \eqref{eq:rellichIdentity2} vanish due to the integration by parts rule for tangential derivatives, see \eqref{eq:intByPartsTang}, and the fact that $w$ is equal to $0$ on  $\partial\Omega$.

  Now, let $h \in \CC_0^1(\R^d, \R^d)$ with $h_k n_k \geq c > 0$ on $\partial \Omega$ be the function from Theorem \ref{thm:smoothApproximation}.
%  Then the divergence theorem gives
%  \begin{align*}
%    \int_{\partial\Omega} h_k n_k |\nabla w|^2 \d\sigma
%    = \int_{\Omega} \div(h\, |\nabla w|^2) \d x 
%    = \int_\Omega \div(h)\, |\nabla w|^2 \d x + \int_\Omega h_k \frac{\partial}{\partial x_k} \Big\{ |\nabla w|^2 \Big\} \d x
%  \end{align*}
%  and we can rewrite the second summand as
%  \begin{align*}
%    \int_{\Omega} h_k \frac{\partial}{\partial x_k} \Big\{ |\nabla w|^2 \Big\} \d x
%    &= \int_{\Omega} h_k \frac{\partial}{\partial x_k} \Big\{ \frac{\partial w_j}{\partial x_i} \cdot \frac{\partial \overline w_j}{\partial x_i}\Big\} \d x \\[0.5em]
%    &= 2 \Re \int_\Omega h_k \frac{\partial^2 w_j}{\partial x_k \partial x_i} \cdot  \frac{\partial \overline w_j}{\partial x_i} \d x \\[0.5em]
%    + 2 \Re \int_
%    &= - 2 \Re \int_\Omega h_k \frac{\partial}{\partial x_k} \Big\{ \frac{\partial^2}{\partial x_i^2} w_j \Big\}\, \overline w_j \d x 
%    + \int_{\partial\Omega} \frac{\partial h_k}{\partial x_i} \cdot \frac{\partial^2 w_j}{\partial x_k \partial x_i} \, \overline w_j \d x\\[0.5em]
%    &= 2 \Re \int_\Omega h_k ( \Delta w_j ) \frac{\partial \overline w_j}{\partial x_k} \d x + 0 \\[0.5em]
%    &= 2 \Re \int_\Omega h_k \big( \frac{\partial \psi}{\partial x_j} + \lambda w_j - \overline u_j \big) \frac{\partial \overline w_j}{\partial x_k} \d x,
%  \end{align*}
%  where in addition to integration by parts we used \eqref{eq:inhomogenousStokes} and the fact that $w = 0$ on $\partial\Omega$.
  We apply the triangle inequality to the Rellich-type identity \eqref{eq:rellichIdentity3} to obtain
  \begin{align}
\label{eq:rellichInequality3}
      \begin{alignedat}{2}
    \int_{\partial\Omega} |\nabla w|^2 \d \sigma 
    &\leq C \, \Big\{ \int_\Omega |\nabla w|^2 \d x + \int_\Omega |\nabla w | |\psi| \d x \\[0.5em]
    &\qquad + |\lambda| \int_\Omega |\nabla w| |w| \d x + \int_\Omega |\nabla w| |u| \d x \Big\}. 
\end{alignedat}
  \end{align}

  The next step consists in deriving estimates which are compatible with the right hand side of \eqref{eq:rellichInequality3}.
  Testing the first equation of \eqref{eq:inhomogenousStokes} with $\overline w$, integration by parts and Lemma~\ref{lem:lambdaIneq} give us as in the proof of Lemma~\ref{lem:laxMilgramIneq} the estimation
  \begin{align*}
    \int_\Omega |\nabla w|^2 \d x + |\lambda| \int_\Omega |w|^2 \d x \leq C \int_\Omega |w| |u| \d x.
  \end{align*}
  The next step consists in using the previous inequality and the Poincar\'{e} inequality to estimate
  \begin{align*}
    \int_\Omega |\nabla w|^2 \d x + (1 + |\lambda| ) \int_\Omega |w|^2 \d x 
    &\leq  (1 + C) \int_\Omega |\nabla w|^2 \d x + |\lambda| \int_\Omega |w|^2 \d x \\[0.5em]
    &\leq C \int_\Omega |w| |u| \d x \\[0.5em]
    &\leq C \, \bigg( \int_\Omega |w|^2 \d x \bigg)^{1/2} \bigg( \int_\Omega |u|^2 \d x \bigg)^{1/2}\\[0.5em]
    \intertext{where for the last step we used the Cauchy-Schwartz inequality. The weighted Young inequality for real numbers allows us to further estimate}
    &\leq \frac{C}{4\varepsilon} \int_\Omega |u|^2 \d x + C \, \varepsilon \int_\Omega |w|^2 \d x \\[0.5em]
    &= \frac{\tilde C}{1 + |\lambda|} \int_\Omega |u|^2 \d x + \frac{1}{2} ( 1 + |\lambda| ) \int_\Omega |w|^2 \d x
  \end{align*}
  if we set $\varepsilon = \frac{(1 + |\lambda|)}{2 C}$.
  Rearranging terms, we can produce our next estimate
  \begin{align}
    \int_\Omega |\nabla w|^2 \d x + ( 1 + |\lambda| ) \int_\Omega |w|^2 \d x
    &\leq \frac{C}{1 + |\lambda|} \int_\Omega |u|^2 \d x.\label{eq:w2lambdaEstimate}
  \end{align}
  Now it's time to harvest: 
  %Using estimate \eqref{eq:rellichInequality3} together with \eqref{eq:w2lambdaEstimate} gives
  Using the weighted Young inequality, we see that we can simplify the right hand side of \eqref{eq:rellichInequality3} to
  \begin{align*}
    \int_{\partial\Omega} |\nabla w|^2 \d\sigma 
    &\leq C_\varepsilon\,  (1 + |\lambda|) \int_{\Omega} |\nabla w|^2 \d x + C\,  |\lambda| \int_\Omega |w|^2 \d x \\[0.5em]
    &\qquad\qquad+ C \int_\Omega |u|^2 \d x + \varepsilon \int_\Omega |\psi|^2 \d x  \\
    &\leq C_\varepsilon \int_\Omega |u|^2 \d x + \varepsilon \int_\Omega |\psi|^2 \d x,
  \end{align*}
  where the second inequality is thanks to estimate \eqref{eq:w2lambdaEstimate}.
  The first term on the right hand side is already fine for \eqref{eq:w2estimate}.
  For the second one we use the estimate $\|\psi\|^{}_{\Ell^2(\Omega)} \leq C\, \|\psi\|^{}_\partial$ and inequality \eqref{eq:psiEstimate} and arrive at
  \begin{align*}
    \varepsilon \int_\Omega |\psi|^2 \d x \leq \varepsilon \, C \int_{\partial\Omega} |\nabla w|^2 \d \sigma + C_\varepsilon \int_{\partial\Omega} |u|^2 \d\sigma.
  \end{align*}
  Choosing $\varepsilon = \frac{1}{2 C}$ and rearranging finally gives the desired estimate \eqref{eq:w2estimate}.
  This concludes our proof.
\end{proof}

The next Theorem states the important fact that in $\Ell^2$ the Dirichlet Stokes resolvent problem has a unique solution.

\begin{thm}
  \label{thm:exAndUniqueSolution}
  Let $\Omega$ be a bounded Lipschitz domain in $\R^d$, $d \geq 2$, with connected boundary and let $\lambda \in \Sigma_\theta$.
  For all $g \in \Ell^2_n(\partial\Omega)$ there exists a unique $u$ and harmonic function $\phi$ which is unique up to constants such that $(u,\phi)$ satisfies \eqref{eq:stokesResolventProblem}, $(u)^* \in \Ell^2(\partial\Omega)$ and $u = g$ on $\partial\Omega$ in the sense of nontangential convergence.
  Moreover, the estimate $ \| (u)^* \|^{}_\partial \leq C \, \| g\|^{}_\partial$ holds and $u$ may be represented by the double layer potential $\dlp_\lambda(f)$ with $\|f\|^{}_\partial \leq C \, \|g\|^{}_\partial$, where in both cases $C$ depends only on $d$, $\theta$ and the Lipschitz character of $\Omega$.
\end{thm}

\begin{proof}
  By Lemma \ref{lem:l2unique}, we already know that the problem under consideration admits at most one solution.
  Therefore, we only have to worry about the existence of a solution.
  In Chapter \ref{chap:3}, it was already established that a solution to the Stokes resolvent problem is given by the double layer potentials $u \coloneqq \dlp_\lambda(f)$ and $\phi \coloneqq \dlp_\Phi(f)$.
  It was also shown that in this way one also solves the $\Ell^2$ Dirichlet problem with boundary data $\dlp_\lambda(f)_+$.
  From Theorem~\ref{thm:nontangentialLimitDoubleLayer} we know that we find this nontangential limit as $((-1/2)I + \K_\lambda^*)(f)$.
  Thus, the central idea of this proof will be to invert the operator $(-1/2)I + \K_\lambda^*$ in order to find the \emph{right} $f$ to plug into the double layer potentials in order to attain the given boundary data $g \in \Ell^2_n(\partial\Omega)$ as a nontangential limit.
  
  We first note that due to Lemma \ref{lem:inverseEstimate} the operator
  \begin{align*}
    T \colon \Ell^2(\partial\Omega; \C^d) \to \Ell^2(\partial\Omega; \C^d), \quad\text x \mapsto -(1/2)x  + \K_{\bar \lambda}x,
  \end{align*}
  is a Fredholm operator on $\Ell^2(\partial\Omega; \C^d)$ with index $0$ and thus the same is true for its adjoint 
  \begin{align*}
    T^* \colon \Ell^2(\partial\Omega; \C^d) \to \Ell^2(\partial\Omega; \C^d),\quad x \mapsto  -(1/2)x + \K_{\bar\lambda}^*x.
  \end{align*}
  In the following paragraphs, we will show that $T^*$ has a bounded inverse.

  We know that for all $f \in \Ell^2(\partial\Omega; \C^d)$ we have $\div(\dlp_\lambda f) = 0$ and therefore 
  \begin{align*}
    \int_{\partial\Omega} T^* f \cdot n \d\sigma =\int_{\partial\Omega} u_+ \cdot n \d\sigma =  0 
  \end{align*}
  holds, where for the first equality we applied Theorem \ref{thm:nontangentialLimitDoubleLayer}. The second equality uses the fact that since $(u)^*$ is integrable, Proposition \ref{prop:approximationArgument} and hence the divergence theorem are available.
  This gives $\Im(T^*) \subseteq \Ell_n^2(\partial\Omega)$.
  Now on the one hand we have
  \begin{align*}
    \operatorname{span}(n)
    = \Ell_n^2(\partial\Omega)^\perp
    \subseteq \Im(T^*)^\perp
    = \ker(T)
  \end{align*}
  and on the other hand, as $T$ is injective on $\Ell^2_n(\partial\Omega)$ by \eqref{eq:inverseEstimate2}, we have that $ \operatorname{span}(n) \supseteq \ker(T)$.
  This yields $\operatorname{span}(n) = \ker(T)$.
  We can use this equality and show that
  \begin{align*}
    \Ell^2_n(\partial\Omega) = \ker(T)^\perp = \overline{\Im(T^*)} = \Im(T^*),
  \end{align*}
  where for the last equality we used the fact that the range of $T^*$ is closed, as usual for Fredholm operators.
  With the same argument we can show for $T$ that
  \begin{align*}
    \ker(T^*)^\perp = \overline{\Im(T)} = \Im(T).
  \end{align*}
  Now let 
  \begin{align*}
    X = \Ell^2_n(\partial\Omega)\quad\text{ and }\quad Y = \Im(T).
  \end{align*}
  Both spaces are closed subspaces of the Hilbert space $\Ell^2(\partial\Omega; \C^d)$ and therefore again Hilbert spaces.
  Consequently, the operator
  \begin{align*}
    K_{Y,X}' \colon Y \to X, \quad x \mapsto T^* x
  \end{align*}
  is invertible by the continuous inverse theorem.
  At this point we could already establish the solvability of the $\Ell^2$ Dirichlet problem.
  But before we do this, in order to derive the additional estimates which were stated in the theorem, we want to bound the operator norm of $\big(K_{Y,X}'\big)^{-1}$ by a constant that does not depend on $\lambda$ but on the sectoriality parameter $\theta$.
  To this end, let us introduce the operator
  \begin{align*}
    K_{X,Y} \colon X \to Y, \quad x \mapsto T x \,.
  \end{align*}
  Now, for $x \in X$ and $y \in Y$ we have that
  \begin{align*}
    \Big\langle x, K_{X,Y}^* y \Big\rangle_X
    & =\Big\langle K_{X,Y}^{} x, y  \Big\rangle_Y
      =\Big\langle T x, y        \Big\rangle_Y
      =\Big\langle x, T^*y       \Big\rangle_X
      =\Big\langle x, K_{Y,X}' y \Big\rangle_X
  \end{align*}
  which shows that $K_{Y,X}' = K_{X,Y}^*$ on $Y$.
  With the above definitions at hand, Lemma \ref{lem:inverseEstimate} states that $K_{X,Y}^{}$ is an invertible operator on $Y$ with operator norm of the inverse bounded from above by some $C > 0$ and $C$ depends at most on $d$, $\theta$ and the Lipschitz character of $\Omega$.
  Therefore, for all $f \in Y = \Im(T) = \Im((1/2)I + \K_{\bar \lambda})$ we have
  \begin{align}
    \label{eq:dualityArgument}
    \| f\|_\partial^{}
    &= \Big\| (K_{Y,X}')^{-1} K_{Y,X}' f \Big\|_\partial^{}
    \leq  \Big\| (K_{X,Y}^*)^{-1} \Big\|_{X,Y}^{}\, \Big\| K_{X,Y}^* f\Big\|_\partial^{}
    \leq C \, \big\|(-(1/2)I + \K_{\bar \lambda}^*)f \big\|_\partial^{}.
  \end{align}
%  Considering once again estimate \eqref{eq:inverseEstimate2} and a duality argument we get that
%  \begin{align}
%    \|f\|_\partial \leq C \| T^* f\|_\partial
%  \end{align}
%  for all $f \in \Im(-(1/2)I + \K_\lambda)$, as
%  \begin{align*}
%    \| (T^*)^{-1} \|_{ \Ell^2_n(\partial\Omega), R(-(1/2)I + \K_{\bar\lambda}),} = \| T^{-1} \|_{R(-(1/2) I + \K_{\bar\lambda}), \Ell^2_n(\partial\Omega)}.
%  \end{align*}

  We are now in position to derive the missing estimates which were stated in the theorem.
  For $g \in \Ell^2_n(\partial\Omega)$, let $f \in \Im(T)$ with $T^* f = g$.
  Fix this $f$ and let $(u,\phi)$ be the respective double layer potentials which were defined in equations \eqref{eq:defDoubleLayer} and \eqref{eq:defDoubleLayerPressure}.
  Then $u_+ = T^* f =  g$ on $\partial\Omega$ by Theorem \ref{thm:nontangentialLimitDoubleLayer}.
  Additionally, we have that
  \begin{align*}
    \|(u)^* \|_\partial^{} \leq C \, \|f\|_\partial^{} \leq C\, \| g\|_\partial^{}
  \end{align*}
  where we used inequality \eqref{eq:lpBoundednessUNontangentialMax} and \eqref{eq:dualityArgument}. In particular, this gives $(u)^* \in \Ell^2(\partial\Omega)$.
  Consequently, all claims of the Theorem have been proven.
\end{proof}

The next theorem can in some sense be regarded as a reverse trace theorem and will play an important role for the proof of the needed reverse H\"{o}lder inequality in the forthcoming chapter.

\begin{thm}
  \label{thm:reverseTrace}
  Let $\Omega$ be a bounded Lipschitz domain in $\R^d$, $d \geq 2$, with connected boundary.
  Let $u \in \HH^1(\Omega; \C^d)$ and $\pi \in \Ell^2(\Omega)$ satisfy the Stokes resolvent problem in $\Omega$ for some $\lambda \in \Sigma_\theta$.
  Then
  \begin{align}
    \label{eq:reverseTrace}
    \Big( \int_\Omega |u|^{p_d} \d x \Big)^{1/{p_d}} \leq C\, \Big( \int_{\partial\Omega} |u|^2 \d\sigma \Big)^{1/2},
  \end{align}
  where $p_d = \frac{2d}{d - 1}$ and $C$ depends only on $d$, $\theta$ and the Lipschitz character of $\Omega$.
\end{thm}

\begin{proof}
  We start by citing the following central inequality
  \begin{align}
    \label{eq:weizhangestimate}
    \Big( \int_\Omega |u|^{p_d} \d x \Big)^{1/{p_d}} \leq C\, \Big( \int_{\partial\Omega} |(u)^*|^2 \d \sigma \Big)^{1/2},
  \end{align}
  where $C$ only depends on $d$ and the Lipschitz constant of $\Omega$.
  The proof of \eqref{eq:weizhangestimate} was carried out by Wei and Zhang in \cite[Lem. 3.3]{weiZhang} and can also be found in Shen's paper \cite[p. 418f.]{Shen2012}.

  Up to now, we do not know whether the right hand side of this inequality equals infinity.
  We only know that $(u,\phi)$ solves the Stokes resolvent problem and that it implicitly solves a Dirichlet problem with boundary data given as $f = \operatorname{Tr}_{\partial\Omega} (u) \in \Ell^2(\partial\Omega; \C^d)$.
  The following part of this proof will show that $u$ coincides with the solution $w \coloneqq \dlp_\lambda(g)$, $g \in \Ell^2(\partial\Omega; \C^d),$  of the $\Ell^2$ Dirichlet problem with boundary data $f$ as given by Theorem \ref{thm:exAndUniqueSolution}. 
  If this is the case, then the knowledge about $(w)^*$ that is contained in Theorem \ref{thm:exAndUniqueSolution} will help us to complete estimate \eqref{eq:weizhangestimate}.

  In order to show that $u = w$ on $\Omega$, consider the sequence $(\Omega_j)_{j \in \N}$ of smooth domains that approximates $\Omega$ from inside as described by Theorem \ref{thm:smoothApproximation}.
  Then, an application of Lemma \ref{lem:l2unique} shows
  \begin{align}
    \label{eq:reverseTraceUnique}
    \int_{\Omega_j} |u - w|^2 \d x \leq C \int_{\partial\Omega_j} |u - w|^2 \d \sigma_j,
  \end{align}
  where $C$ does not depend on $j$ but on the Lipschitz character of $\Omega$.
  Furthermore, the trace theorem on bounded Lipschitz domains gives us for all $h \in \HH^1(\Omega; \C^d)$ that
  \begin{align}
    \label{eq:traceLipschitzDomain}
     \|h\|_\partial^2 \leq C\, \|h\|_{\HH^1(\Omega; \C^d)}^2,
  \end{align}
  where $C$ only depends on $d$ and the Lipschitz character of $\Omega$, see Wei and Zhang \cite[Lem. 2.2]{weiZhang}.
  Now let $\varepsilon > 0$ be given. 
  From the theory of Sobolev spaces it is known that there exists $\varphi_\varepsilon \in \CC^\infty(\overline\Omega; \C^d)$ such that $\big\|\varphi_\varepsilon - u\big\|_{\HH^1(\Omega; \C^d)}^2 \leq {\varepsilon}/({6 \tilde C})$, see Adams and Fournier \cite[Thm. 3.18]{adams}.
  Here, $\tilde C$ denotes the fixed constant from \eqref{eq:traceLipschitzDomain} which is uniform for all $\Omega_j$ and $\Omega$.
  Thanks to Theorem \ref{thm:smoothApproximation}, we know that 
  \begin{align*}
    \int_{\partial\Omega_j} |\varphi_\varepsilon - w|^2 \d\sigma_j \to \int_{\partial\Omega} |\varphi_\varepsilon - u|^2 \d\sigma, \text{ as } j \to \infty
  \end{align*}
  since $w = f$ on $\partial\Omega$ in the sense of nontangential convergence.
  Therefore, we choose $J$ large enough such that for all $j \geq J$ we have
  \begin{align*}
    \int_{\partial\Omega_j} |\varphi_\varepsilon - w|^2 \d\sigma_j \leq \int_{\partial\Omega} |\varphi_\varepsilon - u|^2 \d\sigma + \frac{\varepsilon}{6}\,.
  \end{align*}
  Plugging everything together, this gives us the estimation
  \begin{align*}
    &\int_{\partial\Omega_j} |u - w|^2 \d\sigma_j  \\
    &\qquad\leq 2 \, \bigg\{ \int_{\partial\Omega_j} |u - \varphi_\varepsilon |^2 \d\sigma_j + \int_{\partial\Omega_j} |\varphi_\varepsilon - w|^2 \d\sigma_j \bigg\}
    \leq 4 \tilde C \, \big\| u - \varphi_\varepsilon \big\|_{\HH^1(\Omega)}^2 + \frac{\varepsilon}{3} 
    \leq \varepsilon.
  \end{align*}
  Gluing together this inequality and \eqref{eq:reverseTraceUnique}, we see that $w = u$ holds a.e. in $\Omega$ once we take the limit $j \to \infty$ and remember that $\varepsilon$ was chosen arbitrarily.
  In particular we have $(u)^* = (w)^*$ a.e. on $\partial\Omega$.
  Now, Theorem~\ref{thm:exAndUniqueSolution} gives
  \begin{align*}
%    \| (u)^* \|^{}_\partial 
%    = \| (w)^* \|^{}_\partial \
%    \leq C\, \| f\|^{}_\partial 
%    = C\, \|u\|_\partial,
    \int_{\partial\Omega} |(u)^*|^2 \d\sigma 
    = \int_{\partial\Omega} |(w)^*|^2 \d \sigma
    \leq C \int_{\partial\Omega} |f|^2 \d\sigma 
    = C \int_{\partial\Omega} |u|^2 \d\sigma
  \end{align*}
  where $C$ depends on $d$, $\theta$, and the Lipschitz character of $\Omega$ and we have completed the proof.
\end{proof}

\begin{rem}
  At this point, the choice $p = {2d}({d - 1})^{-1}$ in Theorem \ref{thm:reverseTrace} may seem arbitrary.
  Taking a closer look at the proof by Wei and Zhang in \cite[Lem. 3.3]{weiZhang}, one sees that the choice of $p$ is due to 2 facts: (1) For the dual exponent we have $p' = \frac{2d}{d + 1}$. (2) It holds $\frac{1}{p} - \frac{1}{p'} = \frac{1}{d}$ and thus the \emph{Hardy-Littlewood-Sobolev theorem on fractional integration} may be applied to estimate the $p'$-norm of the \emph{Riesz potential} $I_1(f)$ of a function $f \in \Ell^p(\R^d)$ by the $p$-norm of $f$, see Grafakos \cite[Thm 6.1.3]{grafakos2009modern}.
\end{rem}

In the following remark and the forthcoming chapter, we will make use of an integration argument which can be considered as an application of the following Theorem on \emph{integration along slices}. A proof of this result can be found in Federer \cite[Thm. 3.2.12]{federer}.

\begin{thm}[Co-area formula]
  \label{thm:coarea}
  If $f \colon \R^d \to \R$, $d \geq 2$, is Lipschitz continuous and $\mathfrak{g}$ the representative of a function $g \in \Ell^2(\R^d)$, then
  \begin{align}
    \label{eq:coarea}
    \int_{\R^d} \mathfrak{g}(x) \Big[ \sum_{i = 1}^d |\partial_i f(x)|^2 \Big]^{1/2} \d x = \int_\R \int_{f^{-1}(y)} \mathfrak{g}(x) \d m_{d - 1}(x) \d y,
  \end{align}
  where $m_{d - 1}$ denotes the $(d - 1)$-dimensional Hausdorff measure on $\R^d$.
\end{thm}

Note that the $(d-1)$-dimensional Hausdorff measure on $\R^d$ is comparable to the surface measure $\sigma$ of a bounded Lipschitz domain $\Omega \subseteq \R^d$.
When using Theorem~\ref{thm:coarea} in estimations, we will therefore always be working with the surface measure only.

\begin{rem}
  Let $(u,\phi)$ be a solution of the Stokes resolvent system in the domain $\BB(x_0,\, r) \subseteq \R^d$.
  Then the estimate
  \begin{align}
    \label{eq:interiorEstimateDoubleLayer}
    |\nabla^l u(x_0) | \leq \frac{C_l}{r^{d - 1 + l}} \bigg( \int_{\BB(x_0,\, r)} |u(x)|^2 \d x \bigg)^{1/2}
  \end{align}
  holds for all $l \in \N_0$, where $C_l$ only depends on $d$, $l$ and $\theta$:
  Without loss of generality, we may rescale and translate and thus assume that $x_0 = 0$ and $r = 2$.
  Let $t \in (1,2)$. 
  By Theorem \ref{thm:exAndUniqueSolution}, we know that a solution to the Stokes resolvent system on $\BB(0,\, t) \subsetneq \BB(0,\,2)$ with boundary values $g_t \coloneqq \operatorname{Tr}_{\partial\BB(0,\,t)} (u) \in \Ell^2(\partial\BB(0,\,t); \C^d)$ is given by a boundary layer potential $\dlp_\lambda(f_t)$.
  We can use this fact to derive the desired estimate via
  \begin{align*}
    |\nabla^l u(0)|^2 
    &\leq C\, \bigg( \int_{\partial \BB(0,\, t)} \Big\{ \, \big|\nabla_x^{l+1} \Gamma(y; \lambda)\big|  + \big|\nabla_x^l \Phi(y) \big| \, \Big\} \, |f_t(y)| \d\sigma(y)\bigg)^2 \\[0.5em]
    &\leq C\, \bigg( \int_{\partial \BB(0,\,t)} \frac{|f_t(y)|}{t^{d - 1 + l}}  \d\sigma(y)\bigg)^2 \\[0.5em]
    &\leq C\, \int_{\partial \BB(0,\,t)} |f_t(y)|^2  \d\sigma(y) \\[0.5em]
    &\leq C\, \int_{\partial \BB(0,\,t)} |u(y)|^2 \d\sigma(y), 
  \end{align*}
  where in the last step we used the estimate of $f_t$ against the ``data'' from Theorem \ref{thm:exAndUniqueSolution}.
  Integrating this inequality in $t$ over the interval $(1,2)$ and using the co-area formula \eqref{eq:coarea} with Lipschitz function $|\cdot|$ gives 
  \begin{align*}
    |\nabla^l u(0)|^2 \leq C \int_{\BB(0,\,2)} |u(x)|^2 \d x.
  \end{align*}
  Note that for this argument to work it is crucial that $C$ does only depend on $d$ and the Lipschitz character of $\Omega$.
  Now the claim follows readily.
\end{rem}



% Literaturverzeichnis (beginnt auf einer ungeraden Seite)
\cleardoublepage
\nocite{*}
%\bibliographystyle{geralpha}
\bibliographystyle{acm}
\phantomsection
\addcontentsline{toc}{chapter}{Bibliography}
\bibliography{thesis}

      
% ggf. hier Tabelle mit Symbolen 
% (kann auch auf das Inhaltsverzeichnis folgen)
  
% Stichwortverzeichnis (beginnt auf einer ungeraden Seite)
\cleardoublepage
\printindex
\addcontentsline{toc}{chapter}{Stichwortverzeichnis}
  
% Eidesstattliche Erklaerung.
\cleardoublepage 
\thispagestyle{empty}
\section*{Erklärung zur Abschlussarbeit gemäß § 22 Abs. 7 \\und § 23 Abs. 7 APB TU Darmstadt}
\addcontentsline{toc}{chapter}{Erk\"arung zur Abschlussarbeit}

Hiermit versichere ich, Fabian Gabel, die vorliegende Master-Thesis gemäß § 22 Abs. 7 APB der TU Darmstadt ohne Hilfe Dritter und nur mit den angegebenen Quellen und Hilfsmitteln angefertigt zu haben. Alle Stellen, die Quellen entnommen wurden, sind als solche kenntlich gemacht worden. Diese Arbeit hat in gleicher oder ähnlicher Form noch keiner Prüfungsbehörde vorgelegen. 

Mir ist bekannt, dass im Falle eines Plagiats (§ 38 Abs.2 APB) ein Täuschungsversuch vorliegt, der dazu führt, dass die Arbeit mit 5,0 bewertet und damit ein Prüfungsversuch verbraucht wird. Abschlussarbeiten dürfen nur einmal wiederholt werden.

Bei der abgegebenen Thesis stimmen die schriftliche und die zur Archivierung eingereichte elektronische Fassung gemäß § 23 Abs. 7 APB überein.

%\section*{Thesis Statement pursuant to \S23 paragraph 7 of APB TU Darmstadt}
%\section*{English translation for information purposes only:\\
%Thesis Statement pursuant to § 22 paragraph 7 and § 23 paragraph 7 of APB TU Darmstadt}
%
%I herewith formally declare that I, Fabian Gabel, have written the submitted thesis independently pursuant to § 22 paragraph 7 of APB TU Darmstadt. I did not use any outside support except for the quoted literature and other sources mentioned in the paper. I clearly marked and separately listed all of the literature and all of the other sources which I employed when producing this academic work, either literally or in content. This thesis has not been handed in or published before in the same or similar form.
%
%I am aware, that in case of an attempt at deception based on plagiarism (§38 Abs. 2 APB), the thesis would be graded with 5,0 and counted as one failed examination attempt. The thesis may only be repeated once.
%
%In the submitted thesis the written copies and the electronic version for archiving are pursuant to § 23 paragraph 7 of APB identical in content.
\vspace{20pt}


\noindent
Darmstadt, Datum/Date\vspace{30pt}

% Unterschrift (handgeschrieben)

\noindent
Unterschrift des Autors/Signature of author

\end{document}
